% Utilisation des concepts de base du développement logiciel — Informatique 3ème — Cours signé OnBuch+
% Programme officiel MINESEC. Compiler avec : tectonic N15-utilisation-concepts-base-developpement-logiciel.tex
\newcommand{\DOCMATIERE}{Informatique}
\newcommand{\DOCNIVEAU}{3ème}
\newcommand{\DOCMODULE}{Informatique – classe de 3ème}
\newcommand{\DOCLECON}{N15}
\newcommand{\DOCTITRE}{Utilisation des concepts de base du développement logiciel}
% OnBuch+ — préambule « cours signés » ENRICHI (compilé avec tectonic / XeLaTeX)
% Système visuel « Pop » d'OnBuch+ : fond crème (jamais blanc), bordures encre
% épaisses, ombres « dures » décalées, titres Archivo Black en capitales.
% Version MATHÉMATIQUES : graphes (pgfplots/tikz), boîtes pédagogiques
% (définition, propriété, méthode, attention, le savais-tu, exercice résolu,
% exercice, TP, bilan), page de garde et signature OnBuch+ ancrée en bas.
%
% Macros attendues dans main.tex AVANT \input{preamble} :
%   \DOCMATIERE  \DOCNIVEAU  \DOCMODULE  \DOCLECON (numéro)  \DOCTITRE
\documentclass[11pt]{article}

\usepackage{fontspec}
\usepackage[french]{babel}
\usepackage[a4paper,margin=2cm,top=3.4cm,bottom=2.8cm,headheight=1.4cm,footskip=1.3cm]{geometry}
\usepackage{amsmath,amssymb,amsfonts}
\usepackage{mathtools} % \xmapsto, \coloneqq... souvent utilisés par les modèles
\usepackage{cancel} % simplifications barrées (bilans de forces)
% \abs{z} (module, valeur absolue) et \norm{u} (norme), utilisés par les modèles
\providecommand{\abs}{}\let\abs\relax\DeclarePairedDelimiter{\abs}{\lvert}{\rvert}
\providecommand{\norm}{}\let\norm\relax\DeclarePairedDelimiter{\norm}{\lVert}{\rVert}
\usepackage[table]{xcolor} % [table] : fournit aussi \rowcolors (alternance de lignes)
\usepackage{pagecolor}
\usepackage{tikz}
\usetikzlibrary{arrows.meta,positioning,calc,shapes.geometric,shapes.misc,shapes.arrows,shapes.symbols,decorations.pathmorphing,decorations.pathreplacing,decorations.markings,patterns,shadows,mindmap,trees,fit,backgrounds,matrix,intersections,angles,quotes}
\usepackage{pgfplots}
\usepackage[version=4]{mhchem} % équations nucléaires \ce{^{235}_{92}U}
\usepackage[european,siunitx]{circuitikz} % schémas électriques
\pgfplotsset{compat=1.18}
\usepgfplotslibrary{fillbetween,groupplots} % aires entre courbes (calcul intégral)
\usepackage{enumitem}
\usepackage{fancyhdr}
% [most] charge toutes les bibliothèques tcolorbox (listings, minted, raster,
% documentation...) et gonfle inutilement la mémoire TeX sur un document long
% (~300 boîtes) : on ne charge que ce qui est réellement utilisé.
\usepackage[skins,breakable]{tcolorbox}
\usepackage{graphicx}
\usepackage{colortbl}
\usepackage{booktabs}
\usepackage{tabularx}
\usepackage{diagbox} % case d'en-tête diagonale des tableaux à double entrée
% Colonnes extensibles centrée / gauche / droite (C, L, R), souvent utilisées par les modèles
\newcolumntype{C}{>{\centering\arraybackslash}X}
\newcolumntype{L}{>{\raggedright\arraybackslash}X}
\newcolumntype{R}{>{\raggedleft\arraybackslash}X}
\usepackage{array}
\usepackage{multirow}
\usepackage{multicol}
\usepackage{siunitx}
% parse-numbers=false : accepte aussi \SI{2,5 \times 10^{-3}}{...}
\sisetup{locale=FR,output-decimal-marker={,},group-minimum-digits=4,parse-numbers=false}
\DeclareSIUnit\ohm{\ensuremath{\symup{\Omega}}} % U+2126 absent de la police
\DeclareSIUnit\division{div} % réglages d'oscilloscope (ms/div, V/div)
\DeclareSIUnit\tr{tr} % tours (vitesse de rotation, ex. \SI{1500}{\tr\per\minute})
\DeclareSIUnit\bit{bit}
\DeclareSIUnit\byte{o} % octet
\DeclareSIUnit\inch{po} % pouce
\DeclareSIUnit\ppp{ppp} % points par pouce (résolution d'image)
\usepackage{qrcode}
% Blocs de code source (PHP, C, HTML, JavaScript, SQL) — spécifique à la
% matière Informatique : listings + habillage aux couleurs OnBuch+ (voir le
% style \lstset ci-dessous, à côté des autres boîtes pédagogiques).
\usepackage{listings}
% \degree : commande fréquemment utilisée par les modèles pour un angle ou une
% température en dehors de \SI{}{} (non fournie par les paquets chargés).
\providecommand{\degree}{^{\circ}}
% \qed (fin de démonstration), utilisé par les modèles sans amsthm
\providecommand{\qed}{\hfill$\square$}
% \barre : double barre des valeurs interdites dans un tableau de variations (tkz-tab)
\providecommand{\barre}{\ensuremath{\|}}
% environnement proof (amsthm non chargé) utilisé par les modèles
\ifdefined\proof\else\newenvironment{proof}[1][Démonstration]{\par\noindent\textit{#1.}\ }{\qed\par}\fi
\usepackage{lastpage}
\usepackage{hyperref}
\hypersetup{hidelinks}

% ============================================================
% Palette « Pop » OnBuch+ — mêmes hex que la classe P (plus_ui.dart)
% ============================================================
\definecolor{poporange}{HTML}{F59321}
\definecolor{popdark}{HTML}{E07A0C}
\definecolor{popcream}{HTML}{FFF6E8}
\definecolor{popink}{HTML}{1C1714}
\definecolor{poppurple}{HTML}{7A5AE0}
\definecolor{popgreen}{HTML}{2E9E6A}
\definecolor{poppink}{HTML}{E0457B}
\definecolor{popblue}{HTML}{2F7DE1}
\definecolor{popgold}{HTML}{C98A12}
\definecolor{popmuted}{HTML}{8A7F76}
\definecolor{popline}{HTML}{EADFD2}
\definecolor{popgray}{HTML}{8A7F76}
\colorlet{popgrey}{popgray}
% Alias tolérants (le modèle nomme parfois les couleurs différemment)
\colorlet{popwhite}{white}
\colorlet{popblack}{popink}
\colorlet{popred}{poppink}
\colorlet{popyellow}{popgold}
% Teintes claires pour fonds de tableaux / zones
\colorlet{popblueL}{popblue!10!white}
\colorlet{poporangeL}{poporange!14!white}
\colorlet{popgreenL}{popgreen!12!white}
\colorlet{poppurpleL}{poppurple!12!white}
\colorlet{poppinkL}{poppink!12!white}
\colorlet{popgoldL}{popgold!14!white}

\pagecolor{popcream}

% ============================================================
% Typographie
% ============================================================
\defaultfontfeatures{Ligatures=TeX}
\setmainfont{PlusJakartaSans-Medium.ttf}[Path=../../../fonts/,BoldFont=PlusJakartaSans-Bold.ttf]
% Maths en sans-serif (Fira Math) avec chiffres et lettres droites de Plus
% Jakarta, pour que les formules se fondent dans le texte.
\usepackage[math-style=ISO,bold-style=ISO]{unicode-math}
\setmathfont{FiraMath-Regular.otf}[Path=../../../fonts/]
\setmathfont{PlusJakartaSans-Medium.ttf}[Path=../../../fonts/,range={up/num,up/{Latin,latin},\mathup}]
\usepackage{icomma} % virgule décimale sans espace en mode math
% Caractères mathématiques Unicode absents des polices de texte (Plus Jakarta,
% Archivo Black) : rendus par leur équivalent mathématique, en texte comme en
% formule — sinon ils s'affichent en carré vide (ex. « Arithmétique dans ℤ »).
\usepackage{newunicodechar}
\newunicodechar{ℝ}{\ensuremath{\mathbb{R}}}
\newunicodechar{ℤ}{\ensuremath{\mathbb{Z}}}
\newunicodechar{ℂ}{\ensuremath{\mathbb{C}}}
\newunicodechar{ℕ}{\ensuremath{\mathbb{N}}}
\newunicodechar{ℚ}{\ensuremath{\mathbb{Q}}}
\newunicodechar{𝔻}{\ensuremath{\mathbb{D}}}
\newunicodechar{α}{\ensuremath{\alpha}}
\newunicodechar{β}{\ensuremath{\beta}}
\newunicodechar{γ}{\ensuremath{\gamma}}
\newunicodechar{δ}{\ensuremath{\delta}}
\newunicodechar{ε}{\ensuremath{\varepsilon}}
\newunicodechar{θ}{\ensuremath{\theta}}
\newunicodechar{λ}{\ensuremath{\lambda}}
\newunicodechar{μ}{\ensuremath{\mu}}
\newunicodechar{σ}{\ensuremath{\sigma}}
\newunicodechar{τ}{\ensuremath{\tau}}
\newunicodechar{φ}{\ensuremath{\varphi}}
\newunicodechar{ψ}{\ensuremath{\psi}}
\newunicodechar{ω}{\ensuremath{\omega}}
\newunicodechar{Σ}{\ensuremath{\Sigma}}
% majuscules produites par \MakeUppercase dans les titres (α-aminés -> Α)
\newunicodechar{Α}{\ensuremath{\alpha}}
\newunicodechar{Β}{\ensuremath{\beta}}
\newunicodechar{Γ}{\ensuremath{\Gamma}}
\newunicodechar{Δ}{\ensuremath{\Delta}}
\newunicodechar{Ε}{\ensuremath{\varepsilon}}
\newunicodechar{Θ}{\ensuremath{\Theta}}
\newunicodechar{Λ}{\ensuremath{\Lambda}}
\newunicodechar{Μ}{\ensuremath{\mu}}
\newunicodechar{Τ}{\ensuremath{\tau}}
\newunicodechar{Φ}{\ensuremath{\Phi}}
\newunicodechar{Ψ}{\ensuremath{\Psi}}
\newunicodechar{Ω}{\ensuremath{\Omega}}
\newunicodechar{↦}{\ensuremath{\mapsto}}
\newunicodechar{∈}{\ensuremath{\in}}
\newunicodechar{∉}{\ensuremath{\notin}}
\newunicodechar{∧}{\ensuremath{\wedge}}
\newunicodechar{⇔}{\ensuremath{\Leftrightarrow}}
\newunicodechar{⟺}{\ensuremath{\Longleftrightarrow}}
\newunicodechar{⇒}{\ensuremath{\Rightarrow}}
\newunicodechar{⟹}{\ensuremath{\Longrightarrow}}
\newunicodechar{∘}{\ensuremath{\circ}}
\newunicodechar{∪}{\ensuremath{\cup}}
\newunicodechar{∩}{\ensuremath{\cap}}
\newunicodechar{≡}{\ensuremath{\equiv}}
\newunicodechar{⊂}{\ensuremath{\subset}}
\newunicodechar{⊥}{\ensuremath{\perp}}
\newunicodechar{∃}{\ensuremath{\exists}}
\newunicodechar{∀}{\ensuremath{\forall}}
\newunicodechar{∛}{\ensuremath{\sqrt[3]{\,}}}
\newunicodechar{∜}{\ensuremath{\sqrt[4]{\,}}}
\newunicodechar{⃗}{\ensuremath{{}^{\to}}}
\newunicodechar{ⁿ}{\ifmmode^{n}\else\textsuperscript{\ensuremath{n}}\fi}
\newunicodechar{ˣ}{\ifmmode^{x}\else\textsuperscript{\ensuremath{x}}\fi}
\newunicodechar{ᵇ}{\ifmmode^{b}\else\textsuperscript{\ensuremath{b}}\fi}
\newunicodechar{ʳ}{\ifmmode^{r}\else\textsuperscript{\ensuremath{r}}\fi}
\newunicodechar{ᵘ}{\ifmmode^{u}\else\textsuperscript{\ensuremath{u}}\fi}
\newunicodechar{ʸ}{\ifmmode^{y}\else\textsuperscript{\ensuremath{y}}\fi}
\newunicodechar{ᵃ}{\ifmmode^{a}\else\textsuperscript{\ensuremath{a}}\fi}
\newunicodechar{ᵉ}{\ifmmode^{e}\else\textsuperscript{\ensuremath{e}}\fi}
\newunicodechar{ᵏ}{\ifmmode^{k}\else\textsuperscript{\ensuremath{k}}\fi}
\newunicodechar{ᵖ}{\ifmmode^{p}\else\textsuperscript{\ensuremath{p}}\fi}
\newunicodechar{ᴺ}{\ifmmode^{N}\else\textsuperscript{\ensuremath{N}}\fi}
\newunicodechar{ᶜ}{\ifmmode^{c}\else\textsuperscript{\ensuremath{c}}\fi}
\newunicodechar{ʰ}{\ifmmode^{h}\else\textsuperscript{\ensuremath{h}}\fi}
\newunicodechar{ᵠ}{\ifmmode^{q}\else\textsuperscript{\ensuremath{q}}\fi}
\newunicodechar{ₙ}{\ifmmode_{n}\else\textsubscript{\ensuremath{n}}\fi}
\newunicodechar{ₐ}{\ifmmode_{a}\else\textsubscript{\ensuremath{a}}\fi}
\newunicodechar{ᵢ}{\ifmmode_{i}\else\textsubscript{\ensuremath{i}}\fi}
\newunicodechar{ᵣ}{\ifmmode_{r}\else\textsubscript{\ensuremath{r}}\fi}
\newunicodechar{ₖ}{\ifmmode_{k}\else\textsubscript{\ensuremath{k}}\fi}
\newunicodechar{ᵦ}{\ifmmode_{\beta}\else\textsubscript{\ensuremath{\beta}}\fi}
\newunicodechar{ₓ}{\ifmmode_{x}\else\textsubscript{\ensuremath{x}}\fi}
\newfontfamily\popdisplay{ArchivoBlack-Regular.ttf}[Path=../../../fonts/]
\newfontfamily\poplogo{SpaceGrotesk-Bold.ttf}[Path=../../../fonts/]
\newfontfamily\popsemi{PlusJakartaSans-SemiBold.ttf}[Path=../../../fonts/]
\newfontfamily\popxbold{PlusJakartaSans-ExtraBold.ttf}[Path=../../../fonts/]
\newfontfamily\popxboldls{PlusJakartaSans-ExtraBold.ttf}[Path=../../../fonts/,LetterSpace=3.0]

\setlist[enumerate]{leftmargin=1.7em,itemsep=5pt,topsep=4pt}
\setlist[itemize]{leftmargin=1.7em,itemsep=4pt,topsep=4pt,label={\color{poporange}\textbullet}}
\setlength{\parindent}{0pt}
\setlength{\parskip}{5pt}
\renewcommand{\arraystretch}{1.25}

% pgfplots : style Pop par défaut
\pgfplotsset{
  every axis/.append style={axis line style={popink,line width=1pt},tick style={popink},
    grid style={popline,line width=0.5pt},label style={font=\small},tick label style={font=\footnotesize}},
  popcurve/.style={poporange,line width=1.8pt,no markers,smooth},
}
% Même style disponible dans un \draw TikZ simple (hors pgfplots)
\tikzset{popcurve/.style={poporange,line width=1.8pt}}
\tikzset{popgrid/.style={popline,very thin}}
\colorlet{popcurve}{poporange} % le nom du style est parfois utilisé comme couleur

% ============================================================
% Pill / squiggle
% ============================================================
\newcommand{\poppill}[3]{%
  \tikz[baseline=-0.55ex]\node[draw=popink,line width=1.2pt,rounded corners=2.6pt,%
    fill=#2,inner xsep=6.5pt,inner ysep=3pt]{%
    \popxboldls\fontsize{7}{7}\selectfont\color{#3}\MakeUppercase{#1}};%
}
\newcommand{\popsquiggle}[1]{%
  \tikz[baseline=-0.4ex]{
    \draw[#1,line width=2.6pt,line cap=round]
      plot [smooth] coordinates {(0,0.09) (0.34,0.01) (0.68,0.21) (1.02,0.01) (1.36,0.10)};
  }%
}
% Mot-clé surligné (marqueur orange sous le texte)
\newcommand{\cle}[1]{{\popxbold\color{popdark}#1}}

% ============================================================
% En-tête / pied
% ============================================================
\pagestyle{fancy}
\fancyhf{}
\renewcommand{\headrulewidth}{0pt}
\renewcommand{\footrulewidth}{0pt}
\lhead{\raisebox{-0.15ex}{\poplogo\fontsize{13}{13}\selectfont\color{popink}OnBuch\color{poporange}+}}
\rhead{\poppill{\DOCMATIERE\ \textbullet\ \DOCNIVEAU\ \textbullet\ Leçon \DOCLECON}{white}{popink}}
\lfoot{\popxbold\fontsize{7.6}{7.6}\selectfont\color{popmuted}\MakeUppercase{Cours signé OnBuch+}\ \textbullet\ {\popsemi\DOCTITRE}}
\rfoot{\tikz[baseline=-0.6ex]\node[rounded corners=8.5pt,fill=popink,minimum height=17pt,minimum width=17pt,inner xsep=5pt,inner sep=0]{\popxbold\fontsize{8}{8}\selectfont\color{popcream}\thepage\,/\,\pageref*{LastPage}};}

% ============================================================
% Titres
% ============================================================
\newcommand{\chaptitre}[2]{%
  \begin{tcolorbox}[enhanced,colback=poporange,colframe=popink,boxrule=2.6pt,arc=11pt,
      drop shadow=popink,left=15pt,right=15pt,top=12pt,bottom=11pt,before skip=3pt,after skip=18pt]
    \poppill{#2}{popink}{popcream}\par\vspace{9pt}
    {\raggedright\hyphenpenalty=10000\exhyphenpenalty=10000\popdisplay\fontsize{22}{24}\selectfont\color{popcream}\MakeUppercase{#1}\par}
  \end{tcolorbox}
}
% Section numérotée : pastille orange avec le numéro + titre Archivo
\newcounter{popsec}
\newcommand{\coursec}[1]{%
  \stepcounter{popsec}\setcounter{popsub}{0}%
  \par\vspace{16pt}\noindent
  \begin{minipage}{\linewidth}
  \tikz[baseline=-0.7ex]\node[draw=popink,line width=1.6pt,fill=poporange,rounded corners=4pt,minimum size=22pt,inner sep=0]{\popdisplay\fontsize{12}{12}\selectfont\color{popcream}\thepopsec};%
  \hspace{8pt}{\popdisplay\fontsize{15}{17}\selectfont\color{popink}\MakeUppercase{#1}}\par
  \vspace{4pt}\noindent{\color{poporange}\rule{36pt}{3.6pt}}
  \end{minipage}\par\nopagebreak\vspace{8pt}
}
\newcounter{popsub}
\newcommand{\courssub}[1]{%
  \stepcounter{popsub}%
  \par\vspace{9pt}\noindent{\popxbold\fontsize{11.5}{13}\selectfont\color{popink}{\color{poporange}\thepopsec.\thepopsub}\hspace{6pt}#1}\par\nopagebreak\vspace{4pt}
}

% ============================================================
% Boîtes pédagogiques — même recette : fond blanc, bordure épaisse colorée,
% ombre dure encre, badge plein. Titre optionnel [#1].
% ============================================================
\tcbset{popbase/.style={breakable,enhanced,colback=white,boxrule=2.3pt,arc=9pt,
  shadow={3pt}{-3pt}{0pt}{popink},left=14pt,right=14pt,top=11pt,bottom=11pt,before skip=11pt,after skip=15pt,
  fontupper=\popsemi\fontsize{10.6}{14.5}\selectfont\color{popink},
  fontlower=\popsemi\fontsize{10.6}{14.5}\selectfont\color{popink}}}
% Coche dessinée (le glyphe ✓ n'existe pas dans les polices du document)
\newcommand{\popcheck}[1]{\tikz[baseline=-0.5ex]\draw[#1,line width=1.6pt,line cap=round,line join=round](-0.13,0.0)--(-0.04,-0.09)--(0.13,0.1);}
\providecommand{\coche}{\popcheck{popgreen}}
\providecommand{\resultat}[1]{\par\smallskip\noindent\fcolorbox{popgreen}{popgreenL}{\parbox{\dimexpr\linewidth-2\fboxsep-2\fboxrule\relax}{\textbf{Résultat.} #1}}\par\smallskip}
\newcommand{\popboxhead}[3]{\poppill{#1}{#2}{white}\ifx\relax#3\relax\else\hspace{6pt}{\popxbold\fontsize{10.6}{12}\selectfont\color{popink}#3}\fi\par\vspace{7pt}}

\newtcolorbox{aretenir}[1][]{popbase,colframe=popgreen,before upper={\popboxhead{À retenir}{popgreen}{#1}}}
\newtcolorbox{exemplebox}[1][]{popbase,colframe=poporange,before upper={\popboxhead{Exemple}{poporange}{#1}}}
\newtcolorbox{definition}[1][]{popbase,colframe=popblue,before upper={\popboxhead{Définition}{popblue}{#1}}}
\newtcolorbox{propriete}[1][]{popbase,colframe=poppurple,before upper={\popboxhead{Propriété}{poppurple}{#1}}}
\newtcolorbox{methode}[1][]{popbase,colframe=popgold,before upper={\popboxhead{Méthode}{popgold}{#1}}}
\newtcolorbox{attention}[1][]{popbase,colframe=poppink,before upper={\popboxhead{Attention}{poppink}{#1}}}
\newtcolorbox{experience}[1][]{popbase,colframe=popblue,colback=popblueL,before upper={\popboxhead{Expérience}{popblue}{#1}}}
% « Le savais-tu ? » : carte violette pleine (contexte, culture, Cameroun)
\newtcolorbox{savaistu}[1][]{popbase,colback=poppurple,colframe=popink,
  fontupper=\popsemi\fontsize{10.4}{14.2}\selectfont\color{white},
  before upper={\poppill{Le savais-tu\,?}{white}{poppurple}\ifx\relax#1\relax\else\hspace{6pt}{\popxbold\color{white}#1}\fi\par\vspace{7pt}}}
% Exercice résolu : énoncé en haut, \tcblower puis la solution sur fond crème
\newcounter{popexo}\newcounter{popexr}
\newtcolorbox{exoresolu}[1][]{popbase,colframe=poporange,colbacklower=popcream,
  segmentation style={popink,line width=1.2pt,dashed},
  before upper={\stepcounter{popexr}\popboxhead{Exercice résolu \thepopexr}{poporange}{#1}},
  before lower={\poppill{Solution}{popink}{popcream}\par\vspace{6pt}}}
% Exercice (non résolu en place) — la correction est dans la partie Corrigés
\newtcolorbox{exercice}[1][]{popbase,colframe=popink,
  before upper={\stepcounter{popexo}\popboxhead{Exercice \thepopexo}{popink}{#1}}}
% Corrigé
\newtcolorbox{corrige}[1][]{popbase,colframe=popgreen,colback=popgreenL,
  before upper={\popboxhead{Corrigé}{popgreen}{#1}}}
% Objectifs / prérequis / bilan
\newtcolorbox{objectifs}{popbase,colframe=popink,colback=poporangeL,
  before upper={\popboxhead{Objectifs de la leçon}{popink}{}}}
\newtcolorbox{prerequis}{popbase,colframe=popink,colback=popblueL,
  before upper={\popboxhead{Prérequis}{popblue}{}}}
\newtcolorbox{bilan}{popbase,colframe=popink,colback=white,boxrule=2.8pt,
  before upper={\popboxhead{Fiche bilan}{poporange}{L'essentiel à connaître pour l'examen}}}
% Figure encadrée avec légende
\newtcolorbox{popfigure}[1][]{enhanced,breakable=false,colback=white,colframe=popline,boxrule=1.4pt,arc=8pt,
  left=10pt,right=10pt,top=10pt,bottom=8pt,before skip=10pt,after skip=12pt,halign=center,
  lower separated=false,halign lower=center,
  fontlower=\popsemi\fontsize{9}{11.5}\selectfont\color{popmuted},
  #1}
\newcommand{\legende}[1]{\par\vspace{6pt}{\popsemi\fontsize{9}{11.5}\selectfont\color{popmuted}#1}}

% ============================================================
% Bloc de code source — \begin{lstlisting}[language=Php]...\end{lstlisting}
% (ou language=C, HTML, JavaScript, SQL ; option omise = pseudo-code brut).
% Même recette visuelle que les figures (\popfigure) : cadre encre épais,
% fond clair (popcream), légende possible juste après via \legende{...}
% comme pour une figure (listings ne sait pas arrondir ses coins : on garde
% un cadre rectangulaire, cohérent avec les autres tableaux du document).
% CONTRAT : les modèles ne doivent utiliser QUE \begin{lstlisting}[...] tel
% quel (pas de \begin{tcblisting}, pas de \verb pour un bloc multi-lignes).
% ============================================================
\lstdefinestyle{poplst}{
  backgroundcolor=\color{popcream},
  basicstyle=\popsemi\fontsize{8.6}{11}\selectfont\color{popink},
  keywordstyle=\bfseries\color{popblue},
  commentstyle=\itshape\color{popmuted},
  stringstyle=\color{popgreen},
  numberstyle=\tiny\color{popmuted},
  identifierstyle=\color{popink},
  frame=l,
  rulecolor=\color{popink},
  framerule=1.6pt,
  framesep=7pt,
  framexleftmargin=7pt,
  framexrightmargin=7pt,
  xleftmargin=2pt,
  xrightmargin=2pt,
  breaklines=true,
  breakatwhitespace=false,
  showstringspaces=false,
  showspaces=false,
  showtabs=false,
  tabsize=2,
  columns=flexible,
  keepspaces=true,
  numbers=none,
  aboveskip=10pt,
  belowskip=10pt,
  captionpos=b,
}
\lstset{style=poplst, language=PHP}
% La distribution TeX embarquée ne fournit pas le fichier de langue
% JavaScript de listings (lstlang*.dat) : on la redéfinit nous-mêmes
% pour que \begin{lstlisting}[language=JavaScript] fonctionne.
\lstdefinelanguage{JavaScript}{
  keywords={break,case,catch,class,const,continue,debugger,default,delete,do,
    else,export,extends,finally,for,function,if,import,in,instanceof,let,new,
    return,static,super,switch,this,throw,try,typeof,var,void,while,with,yield,
    async,await,of,get,set},
  keywordstyle=\bfseries\color{popblue},
  ndkeywords={true,false,null,undefined,NaN,Infinity},
  ndkeywordstyle=\bfseries\color{poporange},
  identifierstyle=\color{popink},
  sensitive=true,
  comment=[l]{//},
  morecomment=[s]{/*}{*/},
  string=[b]",
  morestring=[b]',
  morestring=[b]`,
}
% Même limitation pour CSS et les fichiers de configuration (apacheconf,
% nginx, ini...) : ils ne sont pas fournis. CSS et un style générique de
% type "config" (commentaires # ou ;, pas de mots-clés) couvrent les besoins.
\lstdefinelanguage{CSS}{
  keywords={color,background,background-color,margin,padding,border,
    border-radius,font,font-size,font-weight,font-family,width,height,
    display,position,top,left,right,bottom,float,clear,text-align,
    line-height,list-style,cursor,overflow,box-shadow,transition,
    flex,grid,align-items,justify-content,z-index},
  keywordstyle=\bfseries\color{popblue},
  identifierstyle=\color{popink},
  sensitive=false,
  comment=[s]{/*}{*/},
  string=[b]",
  morestring=[b]',
}
\lstdefinelanguage{apacheconf}{
  sensitive=false,
  morecomment=[l]{\#},
}
\lstdefinelanguage{Apache}{
  sensitive=false,
  morecomment=[l]{\#},
}
\lstdefinelanguage{text}{sensitive=false}
\lstdefinelanguage{powershell}{
  keywords={New-Object,New-SmbShare,Get-Acl,Set-Acl,Get-ChildItem,Set-Content,
    Get-Content,Write-Host,ForEach-Object,Where-Object,Select-Object,
    Test-Path,Remove-Item,Copy-Item,Move-Item,Start-Process,Stop-Process,
    If,Else,ElseIf,ForEach,While,Function,Return,Try,Catch},
  keywordstyle=\bfseries\color{popblue},
  identifierstyle=\color{popink},
  sensitive=false,
  comment=[l]{\#},
  string=[b]",
  morestring=[b]',
}
\lstdefinelanguage{ini}{
  sensitive=false,
  morecomment=[l]{;},
  morecomment=[l]{\#},
}

% ============================================================
% Signature OnBuch+
% ============================================================
\newcommand{\poplogoinline}[1]{{\poplogo\fontsize{#1}{#1}\selectfont\color{popink}OnBuch\color{poporange}+}}
\newcommand{\signatureonbuch}{%
  \begin{tcolorbox}[enhanced,colback=popink,colframe=popink,boxrule=2.2pt,arc=10pt,
      drop shadow=poporange,left=14pt,right=14pt,top=10pt,bottom=10pt,before skip=0pt,after skip=0pt]
    \tikz[baseline=-0.7ex]\node[circle,fill=popgreen,draw=popcream,line width=1.3pt,minimum size=22pt,inner sep=0]{\popcheck{white}};%
    \hspace{9pt}\begin{minipage}[c]{0.62\linewidth}
      {\popxbold\fontsize{12}{13}\selectfont\color{popcream}Signé\ \textbullet\ {\poplogo OnBuch\color{poporange}+}}\par\vspace{2pt}
      {\raggedright\popsemi\fontsize{8}{10.5}\selectfont\color{popline}Ludovic A.\\\MakeUppercase{Conforme au programme officiel MINESEC}\par}
    \end{minipage}\hfill
    {\poplogo\fontsize{22}{22}\selectfont\color{popcream}OnBuch\color{poporange}+}
  \end{tcolorbox}%
}

% ============================================================
% Page de garde
% #1 titre · #2 matière · #3 niveau · #4 module · #5 n° leçon · #6 durée ·
% #7 description / objectifs (texte ou itemize)
% ============================================================
\newcommand{\pagegarde}[7]{%
  \thispagestyle{empty}
  \newgeometry{a4paper,margin=1.7cm,top=1.6cm,bottom=1.4cm}%
  \begin{tikzpicture}[remember picture,overlay]
    \fill[poporange!18] ([xshift=-3.2cm,yshift=-3.6cm]current page.north east) circle (4.6cm);
    \fill[poppurple!14] ([xshift=2.4cm,yshift=7.6cm]current page.south west) circle (3.8cm);
  \end{tikzpicture}%
  {\poplogo\fontsize{30}{30}\selectfont\color{popink}OnBuch\color{poporange}+}\hfill
  \raisebox{0.6ex}{\poppill{Cours signé \textbullet\ Premium}{popink}{popcream}}\par
  \vspace{1pt}\popsquiggle{poppurple}\par
  \vspace{8pt}
  \begin{tcolorbox}[enhanced,colback=poporange,colframe=popink,boxrule=2.8pt,arc=13pt,
      drop shadow=popink,left=18pt,right=18pt,top=12pt,bottom=13pt,after skip=10pt]
    \poppill{#2 \textbullet\ #3}{popink}{popcream}\hspace{5pt}\poppill{#4}{white}{popink}\par\vspace{8pt}
    {\popdisplay\fontsize{14}{14}\selectfont\color{popink}LEÇON #5}\par\vspace{4pt}
    {\raggedright\hyphenpenalty=10000\exhyphenpenalty=10000\popdisplay\fontsize{25}{27}\selectfont\color{popcream}\MakeUppercase{#1}\par}
    \vspace{8pt}
    {\popxbold\fontsize{9.5}{9.5}\selectfont\color{popink}\MakeUppercase{Durée conseillée : #6}}
  \end{tcolorbox}
  \begin{tcolorbox}[enhanced,colback=white,colframe=popink,boxrule=2.2pt,arc=10pt,
      drop shadow=popink,left=16pt,right=16pt,top=10pt,bottom=10pt,after skip=8pt]
    \poppill{Dans cette leçon}{poppurple}{white}\par\vspace{5pt}
    \popsemi\fontsize{9.3}{12.6}\selectfont\color{popink}#7
  \end{tcolorbox}
  \noindent\begin{minipage}[t]{0.487\linewidth}
    \begin{tcolorbox}[enhanced,colback=popgreen,colframe=popink,boxrule=2.2pt,arc=10pt,
        drop shadow=popink,left=11pt,right=11pt,top=10pt,bottom=10pt,height=3.1cm,valign=center]
      \tcbox[colback=white,colframe=popink,boxrule=1.3pt,arc=5pt,boxsep=0pt,nobeforeafter,
        left=3pt,right=3pt,top=3pt,bottom=3pt,valign=center]{\qrcode[height=1.9cm]{https://play.google.com/store/apps/details?id=com.luvvix.onbuch&hl=fr}}%
      \hspace{8pt}\begin{minipage}[c]{0.5\linewidth}
        {\popdisplay\fontsize{11}{12.5}\selectfont\color{white}\MakeUppercase{Télécharge OnBuch}}\par\vspace{4pt}
        {\raggedright\popsemi\fontsize{8.4}{11}\selectfont\color{white}Gratuit \textbullet\ Google Play\\Tuteur IA, annales, concours, résultats\par}
      \end{minipage}
    \end{tcolorbox}
  \end{minipage}\hfill
  \begin{minipage}[t]{0.487\linewidth}
    \begin{tcolorbox}[enhanced,colback=poppurple,colframe=popink,boxrule=2.2pt,arc=10pt,
        drop shadow=popink,left=11pt,right=11pt,top=10pt,bottom=10pt,height=3.1cm,valign=center]
      \tikz[baseline=-0.7ex]\node[circle,fill=white,draw=popink,line width=1.3pt,minimum size=1.6cm,inner sep=0]{\popdisplay\fontsize{24}{24}\selectfont\color{poppurple}+};%
      \hspace{8pt}\begin{minipage}[c]{0.6\linewidth}
        {\popdisplay\fontsize{11}{12.5}\selectfont\color{white}\MakeUppercase{Passe à OnBuch+}}\par\vspace{4pt}
        {\raggedright\popsemi\fontsize{8.4}{11}\selectfont\color{white}Tous les cours signés, Tuteur IA illimité, fiches PDF — onglet OnBuch+ de l'app\par}
      \end{minipage}
    \end{tcolorbox}
  \end{minipage}
  \vspace{8pt}
  \popcontenu
  \vfill
  \signatureonbuch
  \restoregeometry
  \newpage
}

% Rangée « Ce cours contient » (page de garde)
\newcommand{\poptile}[3]{% #1 couleur #2 gros texte #3 légende
  \begin{tcolorbox}[enhanced,colback=#1,colframe=popink,boxrule=2pt,arc=9pt,shadow={2.5pt}{-2.5pt}{0pt}{popink},
      left=6pt,right=6pt,top=9pt,bottom=9pt,halign=center,valign=center,height=1.75cm,width=\linewidth,nobeforeafter]
    {\popdisplay\fontsize{12.5}{13}\selectfont\color{white}\MakeUppercase{#2}}\par\vspace{3pt}
    {\centering\popsemi\fontsize{7.8}{9.5}\selectfont\color{white}#3\par}
  \end{tcolorbox}}
\newcommand{\popcontenu}{%
  \par\noindent{\popxboldls\fontsize{7.5}{7.5}\selectfont\color{popmuted}CE COURS SIGNÉ CONTIENT}\par\vspace{7pt}
  \noindent\begin{minipage}[t]{0.235\linewidth}\poptile{popblue}{Cours}{complet, illustré, pas à pas}\end{minipage}\hfill
  \begin{minipage}[t]{0.235\linewidth}\poptile{poporange}{Méthodes}{et exercices résolus}\end{minipage}\hfill
  \begin{minipage}[t]{0.235\linewidth}\poptile{poppink}{Exercices}{type examen + corrigés}\end{minipage}\hfill
  \begin{minipage}[t]{0.235\linewidth}\poptile{popgreen}{Fiche bilan}{l'essentiel à retenir}\end{minipage}\par}

% Sommaire
\newtcolorbox{sommaire}{enhanced,colback=white,colframe=popink,boxrule=2.2pt,arc=10pt,
  drop shadow=popink,left=18pt,right=18pt,top=13pt,bottom=13pt,before skip=6pt,after skip=20pt,
  before upper={\poppill{Sommaire}{poppurple}{white}\par\vspace{9pt}\popsemi\fontsize{10.6}{14.5}\selectfont\color{popink}}}

% Signature de fin de cours — poussée tout en bas de la dernière page
\newcommand{\findecours}{%
  \par\vfill\nopagebreak
  \begin{center}\popsquiggle{poporange}\end{center}\vspace{4pt}
  \signatureonbuch
}
\AtBeginDocument{\def\llbracket{\mathopen{[\mkern-3.5mu[}}\def\rrbracket{\mathclose{]\mkern-3.5mu]}}} % intervalles d'entiers (glyphe ⟦ absent de la police)
% Glyphes absents de Fira Math : redéfinis à partir de glyphes disponibles
\AtBeginDocument{%
  \def\setminus{\mathbin{\backslash}}\let\smallsetminus\setminus
  \def\vdots{\mathord{\vbox{\baselineskip3.2pt\lineskiplimit0pt\kern4pt\hbox{.}\hbox{.}\hbox{.}}}}%
  \def\ddots{\mathinner{\mkern1mu\raise6.5pt\hbox{.}\mkern2mu\raise3.6pt\hbox{.}\mkern2mu\raise0.7pt\hbox{.}\mkern1mu}}%
  \def\triangle{\mathord{\scalebox{0.9}{$\symup{\Delta}$}}}%
  \def\nabla{\mathord{\raisebox{\depth}{\scalebox{1}[-1]{$\symup{\Delta}$}}}}%
  \def\checkmark{\popcheck{popdark}}%
  \def\star{\mathbin{\ast}}\def\bigstar{\mathord{\ast}}%
  \def\diamond{\mathbin{\scalebox{0.6}{$\Diamond$}}}%
  \def\bigcup{\mathop{\mathchoice{\raisebox{-0.3em}{\scalebox{1.7}{$\cup$}}}{\scalebox{1.25}{$\cup$}}{\cup}{\cup}}\displaylimits}%
  \def\bigcap{\mathop{\mathchoice{\raisebox{-0.3em}{\scalebox{1.7}{$\cap$}}}{\scalebox{1.25}{$\cap$}}{\cap}{\cap}}\displaylimits}%
  \def\lesseqqgtr{\mathrel{\lessgtr}}\def\bigoplus{\mathop{\oplus}\displaylimits}%
}
\providecommand{\ssi}{si et seulement si} % abréviation fréquente
\AtBeginDocument{\providecommand{\wideparen}{\overparen}} % arc de cercle

\begin{document}

\pagegarde{Utilisation des concepts de base du développement logiciel}{Informatique}{3ème}{Informatique – classe de 3ème}{N15}{1 séance (fiche de progression harmonisée)}{Cette leçon initie aux concepts fondamentaux du développement logiciel : cycle de vie, algorithmique, structures de contrôle et modularité. Elle permet de structurer une solution informatique pour résoudre un problème concret du quotidien.
\vspace{4pt}
\begin{multicols}{2}\small
\begin{itemize}
\item À la fin de cette leçon, je sais décrire les grandes étapes du cycle de vie d'un logiciel (analyse, conception, codage, test, maintenance).
\item À la fin de cette leçon, je sais distinguer un algorithme (logique) d'un programme (code source) et expliquer le rôle du compilateur ou de l'interpréteur.
\item À la fin de cette leçon, je sais identifier et utiliser les trois structures de contrôle fondamentales : séquence, sélection (SI) et itération (TANT QUE / POUR).
\item À la fin de cette leçon, je sais déclarer des variables adaptées (entier, réel, chaîne, booléen) et leur affecter des valeurs.
\item À la fin de cette leçon, je sais décomposer un problème en sous-problèmes en utilisant des fonctions ou procédures (modularité).
\item À la fin de cette leçon, je sais rédiger un algorithme en pseudo-code standard (style Bac/BEPC) pour résoudre un problème simple de la vie courante.
\end{itemize}\end{multicols}}

\begin{sommaire}
\begin{multicols}{2}
\begin{enumerate}
\item 1. Le cycle de vie d'un logiciel : de l'idée au produit
\item 2. De l'algorithme au programme : langages et traduction
\item 3. Les briques de base : Variables, Types et Structures de contrôle
\item 4. Modularité : Fonctions et Procédures (Diviser pour régner)
\item 5. Valider et documenter : Tests, Débogage et Qualité
\item Activité d'intégration
\item Méthodes et exercices résolus
\item Exercices
\item Corrigés des exercices
\item Fiche bilan
\end{enumerate}
\end{multicols}
\end{sommaire}

\begin{objectifs}
\begin{itemize}
\item À la fin de cette leçon, je sais décrire les grandes étapes du cycle de vie d'un logiciel (analyse, conception, codage, test, maintenance).
\item À la fin de cette leçon, je sais distinguer un algorithme (logique) d'un programme (code source) et expliquer le rôle du compilateur ou de l'interpréteur.
\item À la fin de cette leçon, je sais identifier et utiliser les trois structures de contrôle fondamentales : séquence, sélection (SI) et itération (TANT QUE / POUR).
\item À la fin de cette leçon, je sais déclarer des variables adaptées (entier, réel, chaîne, booléen) et leur affecter des valeurs.
\item À la fin de cette leçon, je sais décomposer un problème en sous-problèmes en utilisant des fonctions ou procédures (modularité).
\item À la fin de cette leçon, je sais rédiger un algorithme en pseudo-code standard (style Bac/BEPC) pour résoudre un problème simple de la vie courante.
\item À la fin de cette leçon, je sais construire un jeu de tests (valeurs normales, limites, invalides) pour valider un algorithme et justifier mes choix.
\item À la fin de cette leçon, je sais expliquer la différence entre erreur de syntaxe, erreur d'exécution et erreur logique (bug).
\end{itemize}
\end{objectifs}

\begin{prerequis}
\begin{itemize}
\item Notion d'algorithme et de pseudo-code (variables, affichage, lecture) vue en 4ème.
\item Utilisation basique d'un éditeur de texte et navigation dans les dossiers/fichiers.
\item Logique booléenne simple (VRAI/FAUX, ET, OU, NON) abordée en mathématiques.
\item Résolution de problèmes arithmétiques simples (opérations, pourcentages).
\end{itemize}
\end{prerequis}

\newpage

\coursec{Le cycle de vie d'un logiciel : de l'idée au produit}

À Yaoundé, Mme Ngo gère une tontine (\emph{Njangi}) de 20 membres. Chaque semaine, elle note manuellement les cotisations dans un cahier usé, calcule les intérêts à la main, détermine l'ordre des tours et prépare les reçus. Cette tâche lui prend des heures, la fatigue aidant, des erreurs s'y glissent : un montant mal reporté, un tour oublié pour Paul, un intérêt mal calculé pour Marie. Un jeune développeur, cousin d'un membre, lui propose de créer une application simple pour automatiser ce processus. Mme Ngo accepte, soulagée. Mais comment passer du cahier de comptes à un programme fiable ? Quelles sont les étapes indispensables pour garantir que l'application ne perdra pas l'argent de la tontine et respectera les règles tacites du groupe ? C'est tout l'objet du \textbf{cycle de vie d'un logiciel} : une démarche structurée, rigoureuse et itérative pour transformer un besoin réel en un outil numérique robuste.

\begin{definition}[Logiciel]
Un \cle{logiciel} est l'ensemble des programmes, procédures, règles et documentation relatifs au fonctionnement d'un système de traitement de l'information. Contrairement au matériel (hardware) que l'on peut toucher, le logiciel est immatériel : c'est la connaissance figée dans du code.
\end{definition}

Le \cle{développement logiciel} n'est pas une simple séance de frappe au clavier. C'est un processus d'ingénierie complet qui va de la naissance de l'idée jusqu'au retrait de l'application (maintenance, évolution, retrait). Pour ne rien oublier — ni le tour de Paul, ni l'intérêt de Marie —, les ingénieurs suivent des modèles de cycle de vie. Le plus classique, et le plus adapté à notre niveau d'initiation, est le \textbf{modèle en cascade (ou cycle en V simplifié)}.

\courssub{Le modèle en cascade : une vue d'ensemble structurée}

Imaginez une cascade : l'eau coule d'un bassin à l'autre, du haut vers le bas. Dans le modèle en cascade, chaque phase doit être terminée (et validée) avant de passer à la suivante. Si l'on découvre une erreur tardivement, il faut « remonter » la cascade, ce qui est coûteux. C'est pourquoi la rigueur en amont (analyse, conception) est vitale.

\begin{popfigure}
\begin{tikzpicture}[
    node distance=1.2cm and 1.5cm,
    phase/.style={rectangle, rounded corners, draw=popblue, thick, fill=popblueL, minimum width=2.8cm, minimum height=1.6cm, align=center, font=\small\bfseries},
    arrow/.style={-Stealth, thick, popdark},
    icon/.style={font=\Large, text=popblue}
]
% Nœuds principaux
\node[phase, align=center] (analyse){\\\textbf{1. Analyse}\\\textbf{des besoins}};
\node[phase, right=of analyse, align=center] (conception){\\\textbf{2. Conception}\\\textbf{(Algorithmes)}};
\node[phase, right=of conception, align=center] (codage){\\\textbf{3. Codage}\\\textbf{(Programmation)}};
\node[phase, right=of codage, align=center] (tests){\\\textbf{4. Tests}\\\textbf{Validation}};
\node[phase, right=of tests, align=center] (maintenance){\\\textbf{5. Maintenance}\\\textbf{Évolution}};
% Flèches principales (cascade)
\draw[arrow] (analyse.east) -- node[above, font=\tiny, sloped] {Cahier des charges} (conception.west);
\draw[arrow] (conception.east) -- node[above, font=\tiny, sloped] {Algorithmes / MCD} (codage.west);
\draw[arrow] (codage.east) -- node[above, font=\tiny, sloped] {Code source} (tests.west);
\draw[arrow] (tests.east) -- node[above, font=\tiny, sloped] {Logiciel livré} (maintenance.west);
% Boucles de retour (Maintenance -> phases précédentes)
\draw[arrow, poporange, dashed, bend left=40] (maintenance.north) to node[above, font=\tiny, text=poporange] {Correction / Nouvelle version} (analyse.north);
\draw[arrow, poporange, dashed, bend left=25] (tests.south) to[out=-90, in=-90] node[below, font=\tiny, text=poporange] {Bugs détectés} (codage.south);
\draw[arrow, poporange, dashed, bend left=15] (codage.south) to[out=-90, in=-90] node[below, font=\tiny, text=poporange] {Erreur de logique} (conception.south);
% Légendes des icônes (remplacement textuel si images absentes)
% Note: En compilation réelle, remplacer \includegraphics par du texte ou des nœuds TikZ purs.
% Ici on suppose des fichiers icons/search.png etc. Pour l'exemple, on met du texte.
\end{tikzpicture}
\legende{Le cycle de vie en cascade simplifié : flux descendant (flèches pleines) et boucles de rétroaction (flèches pointillées orange) pour correction et évolution.}
\end{popfigure}

\begin{aretenir}[Les 5 phases du cycle de vie (Modèle cascade simplifié)]
\begin{enumerate}
    \item \textbf{Analyse des besoins} : Comprendre \emph{quoi} faire et \emph{pourquoi}. Résultat : Cahier des charges.
    \item \textbf{Conception} : Déterminer \emph{comment} faire. Rédaction des algorithmes (pseudo-code), modélisation des données (MCD). Indépendant du langage.
    \item \textbf{Codage (Implémentation)} : Traduire l'algorithme dans un langage de programmation (C, Python, JavaScript...). Résultat : Fichiers sources.
    \item \textbf{Tests (Validation)} : Vérifier que le programme fait \emph{exactement} ce qui est demandé. Tests unitaires, tests d'intégration, tests utilisateur.
    \item \textbf{Maintenance} : Corriger les bugs découverts plus tard, adapter aux nouvelles règles (ex: changement de taux d'intérêt), ajouter des fonctionnalités.
\end{enumerate}
\end{aretenir}

\courssub{Phase 1 : Analyse des besoins — Comprendre le problème réel}

C'est l'étape la plus critique. « Un problème bien posé est à moitié résolu » (John Dewey). Pour la tontine de Mme Ngo, il ne suffit pas de dire « je veux une appli ». Il faut formaliser les règles métier, souvent implicites.

\begin{methode}[Analyser un besoin : la démarche en 3 étapes]
\begin{enumerate}
    \item \textbf{Identifier les données d'entrée (Entrées)} : Quelles informations le système reçoit-il ?
        \begin{itemize}
            \item Montant de la cotisation hebdomadaire (ex: 5\,000 FCFA).
            \item Nombre de membres (ex: 20).
            \item Liste des membres (Nom, Téléphone, Rang dans le tour).
            \item Taux d'intérêt appliqué aux retards (ex: 5\% par semaine de retard).
            \item Date de début de la tontine.
        \end{itemize}
    \item \textbf{Identifier les résultats attendus (Sorties)} : Que doit produire le système ?
        \begin{itemize}
            \item Planning des tours (Qui reçoit la caisse cette semaine ?).
            \item État de la caisse totale (Somme des cotisations + intérêts perçus).
            \item Suivi des retards par membre (Nombre de semaines, Montant dû).
            \item Reçus imprimables ou SMS de confirmation.
        \end{itemize}
    \item \textbf{Déterminer les traitements (Transformations)} : Quels calculs, tris, décisions transforment les entrées en sorties ?
        \begin{itemize}
            \item \textit{Calcul du tour} : Rotation circulaire sur la liste des membres.
            \item \textit{Calcul des intérêts} : Si cotisation payée après le vendredi $\rightarrow$ Pénalité = Cotisation $\times$ Taux $\times$ Semaines de retard.
            \item \textit{Mise à jour caisse} : Caisse = Caisse + Cotisations + Pénalités - Montant versé au bénéficiaire du tour.
            \item \textit{Gestion des exceptions} : Membre décédé, membre exclu, entrée d'un nouveau membre en cours de cycle.
        \end{itemize}
\end{enumerate}
\end{methode}

\begin{exemplebox}[Application à la tontine « Espoir » de Mme Ngo]
Le cahier des charges issu de l'analyse pourrait contenir cette règle métier :
\begin{quote}
\textbf{Règle R1 (Calcul pénalité) :} Si un membre verse sa cotisation avec $n$ semaines de retard ($n \ge 1$), il doit payer une pénalité $P = C \times t \times n$, où $C$ est la cotisation hebdomadaire (5\,000 FCFA) et $t$ le taux d'intérêt (0,05 soit 5\%).
\end{quote}
Cette règle est claire, non ambiguë, chiffrée. Elle servira de base aux tests plus tard.
\end{exemplebox}

\courssub{Phase 2 : Conception — L'algorithme, langage universel}

Une fois le \emph{quoi} connu, on définit le \emph{comment} logique, sans encore choisir C, Python ou JavaScript. On écrit un \cle{algorithme}.

\begin{definition}[Algorithme]
Un \cle{algorithme} est une suite \textbf{finie} et \textbf{non ambiguë} d'instructions élémentaires permettant de résoudre un problème ou d'obtenir un résultat à partir de données d'entrée. Il s'exprime en \textbf{pseudo-code} (français structuré, anglais courant) compréhensible par tout informaticien, indépendamment de la machine ou du langage.
\end{definition}

\begin{attention}[Algorithme $\neq$ Programme]
Ne confondez jamais les deux :
\begin{itemize}
    \item L'\textbf{algorithme} décrit la \textbf{logique} (la recette de cuisine). Il se lit par l'humain, il est universel, il n'a pas de syntaxe stricte imposée par un compilateur.
    \item Le \textbf{programme} est la \textbf{traduction} de l'algorithme dans un \textbf{langage de programmation précis} (C, Python, JavaScript...). Il obéit à une \textbf{syntaxe stricte} (points-virgules, accolades, majuscules/minuscules). Il s'exécute par la machine.
\end{itemize}
\textbf{Piège fréquent :} Écrire du pseudo-code avec la syntaxe du C (ex: \texttt{for(i=0;i<n;i++)}) en pensant faire de l'algorithmique. Non ! En algorithmique, on écrit : \texttt{Pour i de 0 à n-1 faire}.
\end{attention}

Voici la traduction de la Règle R1 en algorithmique standard (style français, proche de l'écriture au BEPC).
\begin{lstlisting}[language=, caption={Algorithme : CalculPenalite}]
Algorithme CalculPenalite
// But : Calculer la pénalité de retard pour un membre de la tontine
// Entrées : cotisation (Réel), tauxInteret (Réel), semainesRetard (Entier)
// Sortie : penalite (Réel)

Variables
    cotisation, tauxInteret, penalite : Réel
    semainesRetard : Entier
Début
    Afficher "Entrez la cotisation hebdomadaire (FCFA) : "
    Saisir cotisation
    Afficher "Entrez le taux d'intérêt (ex: 0.05 pour 5%) : "
    Saisir tauxInteret
    Afficher "Entrez le nombre de semaines de retard : "
    Saisir semainesRetard

    Si semainesRetard > 0 Alors
        penalite <- cotisation * tauxInteret * semainesRetard
        Afficher "Pénalité à payer : ", penalite, " FCFA"
    Sinon
        Afficher "Aucun retard, pas de pénalité."
        penalite <- 0
    FinSi
Fin
\end{lstlisting}


Remarquez la clarté : \texttt{Si ... Alors ... Sinon ... FinSi}, \texttt{<-} pour l'affectation, pas de \texttt{;} ni d'accolades. C'est lisible par un gestionnaire de tontine non informaticien.

\courssub{Phase 3 : Codage — Traduction dans un langage de programmation}

C'est le passage du « langage humain structuré » au « langage machine compréhensible par compilateur/interpréteur ». On choisit un langage adapté : \texttt{C} pour la performance/système, \texttt{Python} pour la rapidité de script, \texttt{JavaScript} pour le web.

Le même algorithme devient un programme C compilable. La rigueur syntaxique est obligatoire.
\begin{lstlisting}[language=C, caption={Programme C : calcul\_penalite.c}]
#include <stdio.h>   // Bibliothèque standard pour entrée/sortie
#include <stdlib.h>  // Pour EXIT_SUCCESS

int main() {
    // 1. Déclaration des variables (Typage fort en C)
    float cotisation, tauxInteret, penalite;
    int semainesRetard;

    // 2. Saisie des données (Entrées)
    printf("Entrez la cotisation hebdomadaire (FCFA) : ");
    scanf("%f", &cotisation);  // & obligatoire pour l'adresse mémoire

    printf("Entrez le taux d'interet (ex: 0.05 pour 5%%) : ");
    scanf("%f", &tauxInteret);

    printf("Entrez le nombre de semaines de retard : ");
    scanf("%d", &semainesRetard);

    // 3. Traitement (Structure de contrôle Si/Else)
    if (semainesRetard > 0) {
        penalite = cotisation * tauxInteret * semainesRetard;
        printf("Penalite a payer : %.2f FCFA\n", penalite);
    } else {
        printf("Aucun retard, pas de penalite.\n");
        penalite = 0.0f;
    }

    return EXIT_SUCCESS; // Fin normale du programme
}
\end{lstlisting}


\begin{propriete}[Compilation vs Interprétation]
Une fois le code source écrit (fichier \texttt{.c}, \texttt{.py}, \texttt{.js}), il doit devenir executable.
\begin{itemize}
    \item \textbf{Langage compilé (C, C++)} : Un \cle{compilateur} (ex: \texttt{gcc}) traduit \textbf{tout} le code source en \textbf{code machine} (binaire, fichier \texttt{.exe} ou \texttt{.out}) \textbf{avant} l'exécution. Erreurs = Erreurs de compilation (syntaxe). Exécution = Rapide.
    \item \textbf{Langage interprété (Python, JavaScript)} : Un \cle{interpréteur} lit, traduit et exécute \textbf{ligne par ligne} à la volée. Pas de fichier exécutable intermédiaire. Erreurs = Erreurs d'exécution (runtime). Développement = Plus interactif.
\end{itemize}
Pour notre tontine, un petit script Python ou une page web HTML/JS suffit amplement.
\end{propriete}

\courssub{Phases 4 et 5 : Tests et Maintenance — La qualité dans la durée}

Un programme qui compile (ou s'exécute sans planter) n'est pas un programme \textbf{correct}. Il doit faire \emph{exactement} ce que le cahier des charges demande.

\paragraph{Les Tests (Validation).}
\begin{itemize}
    \item \textbf{Tests unitaires} : On teste chaque fonction isolément. Ex: Appeler \texttt{CalculPenalite(5000, 0.05, 2)} $\rightarrow$ Attendu : 500 FCFA.
    \item \textbf{Tests d'intégration} : On teste l'enchaînement : Saisie $\rightarrow$ Calcul $\rightarrow$ Mise à jour caisse $\rightarrow$ Affichage reçu.
    \item \textbf{Tests utilisateur (Recette)} : Mme Ngo utilise l'appli sur une vraie semaine de tontine. Elle valide : « Oui, c'est bien le tour de Paul, l'intérêt de Marie est juste. »
\end{itemize}

\paragraph{La Maintenance.}
Le logiciel vit. Au bout de 6 mois, la tontine décide : « Le taux passe à 10\% » ou « On ajoute une cotisation exceptionnelle pour la fête de fin d'année ».
\begin{itemize}
    \item \textbf{Maintenance corrective} : Corriger un bug (ex: l'appli plantait si un membre payait 0 FCFA).
    \item \textbf{Maintenance adaptative} : Adapter à l'environnement (ex: changement de format de date, nouvelle version d'Android).
    \item \textbf{Maintenance évolutive} : Ajouter des fonctionnalités (ex: envoi SMS automatique via API Orange Money).
\end{itemize}
C'est pourquoi le code doit être \textbf{lisible}, \textbf{commenté} et \textbf{modulaire} (voir section 4 sur les fonctions).

\begin{exoresolu}[Test unitaire de la fonction de pénalité]
\textbf{Énoncé :} On dispose de la fonction C suivante :
\begin{lstlisting}[language=C]
float calculerPenalite(float cotisation, float taux, int semaines) {
    if (semaines > 0) return cotisation * taux * semaines;
    return 0.0f;
}
\end{lstlisting}
Proposer un jeu de tests unitaires complet (cas nominaux, cas limites, cas d'erreur) sous forme de tableau.

\textbf{Solution détaillée :}
Un bon jeu de tests couvre les branches du \texttt{if} et les valeurs extrêmes.

\begin{center}
\begin{tabularx}{\linewidth}{|l|X|X|X|X|}
\hline
\rowcolor{popblueL}
\textbf{Cas de test} & \textbf{Entrées (cotisation, taux, semaines)} & \textbf{Résultat attendu} & \textbf{Résultat obtenu} & \textbf{Statut} \\
\hline
Cas nominal 1 & (5000, 0.05, 2) & 500.0 & 500.0 & \textcolor{popgreen}{OK} \\
\hline
Cas nominal 2 & (10000, 0.10, 1) & 1000.0 & 1000.0 & \textcolor{popgreen}{OK} \\
\hline
Cas limite : Seuil & (5000, 0.05, 1) & 250.0 & 250.0 & \textcolor{popgreen}{OK} \\
\hline
Cas limite : Zéro semaine & (5000, 0.05, 0) & 0.0 & 0.0 & \textcolor{popgreen}{OK} \\
\hline
Cas limite : Semaines négatives (erreur métier) & (5000, 0.05, -2) & 0.0 (ou erreur) & 0.0 & \textcolor{poporange}{À revoir*} \\
\hline
Cas extrême : Grande valeur & (1000000, 0.5, 52) & 26\,000\,000.0 & 26\,000\,000.0 & \textcolor{popgreen}{OK} \\
\hline
\end{tabularx}
\end{center}
\textbf{Analyse :} Le cas « Semaines négatives » révèle une faille : l'algorithme ne protège pas contre une saisie aberrante (retour dans le temps). En maintenance, on ajoutera une vérification \texttt{if (semaines < 0) return -1; // Code erreur} ou une assertion. C'est ça, la valeur des tests : trouver les trous \emph{avant} le client.
\end{exoresolu}

\begin{savaistu}[Le Cameroun numérique suit le cycle de vie]
Au Cameroun, des startups à succès comme \textbf{GiftedMom} (suivi de grossesse par SMS en zone rurale) ou \textbf{Agro-Hub} (mise en relation agriculteurs/acheteurs) n'ont pas codé « à l'aveugle ». Elles ont passé des mois en \textbf{analyse de besoins} sur le terrain (parler aux sages-femmes, aux cultivateurs), ont conçu des \textbf{algorithmes robustes} (gestion hors-ligne, synchronisation différée), codé en \textbf{Java/Kotlin/Swift}, testé avec de vrais utilisateurs (tests d'acceptation) et assurent une \textbf{maintenance continue} (nouvelles versions Play Store régulières). Le cycle de vie n'est pas une théorie d'école : c'est la différence entre un prototype qui plante et un service qui sauve des vies ou nourrit des familles.
\end{savaistu}

\begin{experience}[TP Guidé : Simuler le cycle de vie sur papier (20 min)]
\textbf{Objectif :} S'approprier les 5 phases sans ordinateur.
\textbf{Matériel :} Feuille A4, stylo, cerveau.
\textbf{Consigne :} En binôme, concevez l'algorithme (pseudo-code) pour \textbf{« Déterminer le bénéficiaire du tour de la semaine »} dans la tontine de Mme Ngo.
\begin{enumerate}
    \item \textbf{Analyse (3 min) :} Listez les entrées (liste membres, index tour actuel), la sortie (nom du bénéficiaire), la règle (rotation circulaire, si membre sorti $\rightarrow$ passer au suivant).
    \item \textbf{Conception (7 min) :} Écrivez l'algorithme en pseudo-code clair (Variables, Début, Pour/Si, Fin). \emph{Ne codez pas en C !}
    \item \textbf{Codage mental (3 min) :} L'un dicte l'algo, l'autre « exécute » sur un exemple papier (5 membres : A, B, C, D, E ; tour actuel = 2 $\rightarrow$ C). Vérifiez.
    \item \textbf{Test (5 min) :} Testez les cas : Liste vide, 1 seul membre, membre bénéficiaire déjà sorti (ex: C a déjà eu son tour et est retiré de la liste active). Corrigez l'algo si bug.
    \item \textbf{Maintenance (2 min) :} Comment modifier l'algo si on veut gérer une « priorité urgence » (un membre peut demander à passer devant moyennant majoration) ?
\end{enumerate}
\textbf{Restitution :} Un binôme présente son algo au tableau, la classe joue le rôle de « testeur » en proposant des cas tordus.
\end{experience}

\coursec{De l'algorithme au programme : langages et traduction}

Dans la section précédente, nous avons vu qu'un logiciel naît d'un besoin, qu'on le conçoit (analyse, algorithmes) avant de le réaliser. Nous voici maintenant au cœur de cette réalisation : \textbf{comment transformer un algorithme écrit en français structuré (pseudo-code) en un programme que l'ordinateur comprend et exécute ?} C'est le passage de la logique humaine à l'exécution machine.

\courssub{Du pseudo-code au langage de programmation : pourquoi changer de notation ?}

Rappelons-nous : un \textbf{algorithme} est une suite finie et non ambiguë d'instructions pour résoudre un problème. En classe, nous l'écrivons en \textbf{pseudo-code} (ou \emph{Français Structuré}) : \texttt{SI ... ALORS ... SINON ... FIN SI}, \texttt{POUR i DE 1 À 10 FAIRE}, \texttt{AFFICHER}. Cette notation est \emph{indépendante de toute machine} : elle sert à \emph{réfléchir}, à \emph{communiquer} entre humains, à \emph{valider la logique} avant de coder.

\begin{definition}[Pseudo-code (Français Structuré)]
Notation algorithmique normalisée, proche du langage naturel, utilisant des mots-clés réservés (\texttt{DEBUT}, \texttt{FIN}, \texttt{SI}, \texttt{TANTQUE}, \texttt{POUR}, \texttt{AFFICHER}, \texttt{LIRE}, \texttt{<--}) pour décrire la logique d'un programme sans se soucier de la syntaxe stricte d'un langage machine. Il est \textbf{exécutable par l'esprit humain}, pas par le processeur.
\end{definition}

L'ordinateur, lui, ne comprend que le \textbf{binaire} (suites de 0 et 1). Écrire directement en binaire est impossible pour un humain (trop lent, trop d'erreurs, illisible). Nous avons donc inventé des \textbf{langages de programmation} : des notations formelles, strictes, textuelles, qui servent d'\emph{interface} entre la pensée algorithmique et le matériel.

\begin{propriete}[Niveaux de langage]
On classe les langages selon leur proximité avec le matériel (bas niveau) ou avec l'humain (haut niveau) :
\begin{itemize}
    \item \textbf{Langage machine (1\up{re} génération) :} Binaire pur (\texttt{10110000 01100001}). Exécuté directement par le CPU. Incompréhensible, non portable.
    \item \textbf{Assembleur (2\up{de} génération) :} Mnémoniques textuels (\texttt{MOV AL, 61h}). Correspondance 1-pour-1 avec le langage machine. Toujours dépendant du processeur (x86, ARM). Utilisé pour pilotes, systèmes embarqués, optimisation extrême.
    \item \textbf{Langages de haut niveau (3\up{me} génération et +) :} Proches de l'anglais et des maths (\texttt{if (x > 5) \{ ... \}}, \texttt{print("Grand")}). \textbf{Indépendants de la machine} (portables). Un même code source peut tourner sur Windows, Linux, Mac, téléphone, après traduction adaptée. Exemples : \cle{C}, \cle{C++}, \cle{Java}, \cle{Python}, \cle{JavaScript}, \cle{C\#}.
\end{itemize}
\end{propriete}

\begin{aretenir}[Choix pédagogique en 3\up{me}]
Au BEPC, nous manipulons principalement le \textbf{pseudo-code} (pour l'algorithmique pure) et nous \textbf{découvrons la syntaxe concrète en C et en JavaScript} (langages au programme officiel), avec Python souvent utilisé comme passerelle pédagogique pour l'interprétation. L'essentiel n'est pas de mémoriser la syntaxe par cœur, mais de comprendre \textbf{la correspondance structurelle} : un \texttt{SI} pseudo-code devient un \texttt{if} en C/JS/Python, une boucle \texttt{POUR} devient un \texttt{for} ou \texttt{while}, une variable s'écrit \texttt{x = 10} (Python/JS) ou \texttt{int x = 10;} (C).
\end{aretenir}

\courssub{La traduction : Compilation vs Interprétation}

Un \textbf{code source} (fichier texte \texttt{.c}, \texttt{.js}, \texttt{.py}, \texttt{.java}) ne s'exécute pas tout seul. Il doit être \emph{traduit} en instructions machine. Il existe deux grandes stratégies : la \textbf{compilation} et l'\textbf{interprétation}.

\begin{popfigure}
\begin{tikzpicture}[
    node distance=1.8cm and 1.2cm,
    box/.style={draw=popdark, thick, rounded corners, fill=popcream, minimum width=3.2cm, minimum height=1.2cm, align=center, font=\small},
    arrow/.style={-Stealth, thick, popblue},
    arrowlabel/.style={midway, fill=white, font=\scriptsize, text=popdark, inner sep=2pt}
]
% Compilation path
\node[box, align=center] (srcC){Code Source C\\\texttt{programme.c}};
\node[box, right=of srcC, align=center] (comp){Compilateur\\\texttt{gcc}};
\node[box, right=of comp, align=center] (exeC){Exécutable\\\texttt{a.out / programme.exe}};
\node[box, right=of exeC, align=center] (runC){Exécution CPU\\Résultats};
\draw[arrow] (srcC) -- node[arrowlabel] {Traduit \textbf{tout} le code} (comp);
\draw[arrow] (comp) -- node[arrowlabel] {Génère} (exeC);
\draw[arrow] (exeC) -- node[arrowlabel] {Lance} (runC);
% Interpretation path (below)
\node[box, below=2.5cm of srcC, align=center] (srcPy){Code Source JS / Python\\\texttt{script.js / script.py}};
\node[box, right=of srcPy, align=center] (interp){Interpréteur / Moteur JS\\\texttt{node / python3 / Navigateur}};
\node[box, right=of interp, align=center] (runPy){Exécution pas à pas\\Résultats / Erreurs};
\draw[arrow] (srcPy) -- node[arrowlabel] {Lit, traduit, exécute \textbf{ligne par ligne}} (interp);
\draw[arrow] (interp) -- node[arrowlabel] {Pilote} (runPy);
% Labels
\node[font=\bfseries\small, text=popblue, above=0.2cm of srcC] {\textbf{COMPILATION (C, C++, Java*)}};
\node[font=\bfseries\small, text=poporange, above=0.2cm of srcPy] {\textbf{INTERPRÉTATION (JavaScript, Python, Bash)}};
\end{tikzpicture}
\legende{Schéma comparatif : Compilation (bloc unique, fichier exécutable intermédiaire) vs Interprétation (flux continu, pas de fichier exécutable persistant). *Java compile vers du bytecode pour machine virtuelle (JVM), approche hybride.}
\end{popfigure}

\begin{definition}[Compilateur]
Programme (ex: \texttt{gcc} pour C/C++, \texttt{javac} pour Java) qui analyse \textbf{la totalité} du code source, vérifie la syntaxe, optimise, et produit un \textbf{fichier exécutable} (code machine natif ou bytecode) distinct du source. L'exécution a lieu \emph{après}, séparément. \textbf{Avantage :} rapidité d'exécution, détection précoce des erreurs de syntaxe. \textbf{Inconvénient :} cycle \emph{Éditer $\to$ Compiler $\to$ Lancer} plus long ; l'exécutable est souvent spécifique à un OS/architecture.
\end{definition}

\begin{definition}[Interpréteur]
Programme (ex: \texttt{python3}, \texttt{node} pour JS, moteur JS du navigateur) qui lit le code source \textbf{instruction par instruction}, la traduit à la volée en actions machine, et l'exécute immédiatement. Pas de fichier exécutable intermédiaire persistant. \textbf{Avantage :} cycle \emph{Éditer $\to$ Lancer} immédiat, portabilité totale (même source partout), idéal pour scripts, tests, apprentissage. \textbf{Inconvénient :} exécution plus lente (analyse répétée), erreurs de syntaxe découvertes seulement quand la ligne est atteinte.
\end{definition}

\begin{exemplebox}[Exemple chiffré : Même algorithme, trois syntaxes (Pseudo-code, C, JavaScript)]
Algorithme : \emph{« Lire un entier x. Si x est strictement supérieur à 5, afficher "Grand", sinon afficher "Petit ou égal". »}

\begin{lstlisting}[language=C, caption=C (Compilé avec gcc), label=lst:c]
/* Fichier : test_if.c */
#include <stdio.h>             // Bibliothèque entrée/sortie

int main(void) {               // Point d'entrée obligatoire
    int x;                     // Déclaration + TYPE obligatoire
    printf("Entier x : ");
    scanf("%d", &x);           // Lecture formatée (adresse &x)
    if (x > 5) {               // Parenthèses OBLIGATOIRES, accolades pour bloc
        printf("Grand\n");
    } else {
        printf("Petit ou egal\n");
    }
    return 0;                  // Code retour au système
}
\end{lstlisting}

\begin{lstlisting}[language=JavaScript, caption=JavaScript (Interprété par navigateur / Node.js), label=lst:js]
// Fichier : test_if.js
const readline = require('readline').createInterface({
    input: process.stdin,
    output: process.stdout
});

readline.question("Entier x : ", (input) => {
    const x = parseInt(input); // Conversion chaîne -> entier
    if (x > 5) {               // Parenthèses OBLIGATOIRES, accolades recommandées
        console.log("Grand");
    } else {
        console.log("Petit ou egal");
    }
    readline.close();
});
\end{lstlisting}

\begin{lstlisting}[language=Python, caption=Python 3 (Interprété), label=lst:py]
# Fichier : test_if.py
x = int(input("Entier x : "))  # Lecture clavier (str -> int)
if x > 5:                      // Deux-points OBLIGATOIRE, indentation OBLIGATOIRE
    print("Grand")             // Bloc SI (4 espaces)
else:
    print("Petit ou egal")     // Bloc SINON
\end{lstlisting}

\begin{lstlisting}[caption=Pseudo-code (Français Structuré), label=lst:pseudo]
Algorithme ComparerX
Variable x : Entier
Debut
    Afficher "Entier x : "
    Lire x
    Si x > 5 Alors
        Afficher "Grand"
    Sinon
        Afficher "Petit ou egal"
    Fin Si
Fin
\end{lstlisting}

\textbf{Observations clés :}
\begin{itemize}
    \item \textbf{Types :} Python/JS infèrent le type (\texttt{x} devient \texttt{number/int} à l'affectation). C l'exige \textbf{avant} (\texttt{int x;}).
    \item \textbf{Blocs :} Python utilise l'\textbf{indentation} (recul). C et JavaScript utilisent les \textbf{accolades \{\}} (l'indentation est libre mais \emph{recommandée} pour la lisibilité).
    \item \textbf{Entrée/Sortie :} Python : \texttt{input()}/\texttt{print()}. C : \texttt{scanf()}/\texttt{printf()} (format \texttt{\%d}, adresse \texttt{\&x}). JS (Node) : \texttt{readline}/\texttt{console.log} ; JS (Navigateur) : \texttt{prompt()}/\texttt{alert()} ou DOM.
    \item \textbf{Structure :} C impose une fonction \texttt{main()}. JS et Python exécutent séquentiellement le script (ou via fonctions fléchées/callbacks pour JS asynchrone).
\end{itemize}
\end{exemplebox}

\courssub{Code source, code objet, exécutable : le vocabulaire précis}

\begin{center}
\begin{tabularx}{\linewidth}{|l|X|X|}
\hline
\rowcolor{popblueL}
\textbf{Terme} & \textbf{Description} & \textbf{Exemples d'extensions} \\
\hline
\cle{Code source} & Texte écrit par le programmeur, lisible par l'humain. & \texttt{.c}, \texttt{.cpp}, \texttt{.js}, \texttt{.py}, \texttt{.java}, \texttt{.cs} \\
\hline
\cle{Code objet} & Résultat intermédiaire de la compilation (binaire non lié, symboles non résolus). & \texttt{.o} (Linux), \texttt{.obj} (Windows) \\
\hline
\cle{Édition de liens (Linking)} & Fusion du code objet avec les bibliothèques (ex: \texttt{printf}, \texttt{math.h}) pour résoudre les adresses. & — \\
\hline
\cle{Exécutable} & Fichier final prêt à être chargé en mémoire par l'OS et exécuté par le CPU. & \texttt{a.out}, \texttt{.exe}, \texttt{.elf} (sans extension souvent sous Linux) \\
\hline
\cle{Bytecode} & Code intermédiaire portable (machine virtuelle). & \texttt{.class} (Java), \texttt{.pyc} (Python cache) \\
\hline
\end{tabularx}
\end{center}

\begin{attention}[Ne pas confondre]
\begin{itemize}
    \item \textbf{Fichier source} $\neq$ \textbf{Fichier exécutable}. On ne modifie \textbf{jamais} l'exécutable directement (binaire). On modifie le source, puis on recompile/relance.
    \item En Python, les fichiers \texttt{.pyc} (dans \texttt{\_\_pycache\_\_}) sont du \emph{bytecode mis en cache} pour accélérer le démarrage suivant. Ce sont des fichiers temporaires, on ne les édite pas.
    \item \texttt{gcc programme.c} fait \textbf{les deux étapes} : compilation $\to$ édition de liens $\to$ produit \texttt{a.out} (ou \texttt{programme.exe}).
\end{itemize}
\end{attention}

\courssub{Les trois familles d'erreurs : les identifier pour les corriger}

Quand « ça ne marche pas », il faut savoir \textbf{où} et \textbf{quand} l'erreur survient. C'est la compétence n°1 du développeur débutant.

\begin{propriete}[Classification des erreurs]
\begin{enumerate}
    \item \textbf{Erreurs de compilation (Syntaxiques) :} Détectées par le \textbf{compilateur} (ou l'interpréteur \emph{avant} d'exécuter la ligne fautive en Python/JS). Le programme \textbf{ne démarre pas}.
        \begin{itemize}
            \item Oubli de \texttt{;} (C, JS), \texttt{:} (Python), parenthèses/accolades déséquilibrées.
            \item Variable non déclarée (C), mot-clé mal orthographié (\texttt{whille}, \texttt{pritn}, \texttt{consloe.log}).
            \item Indentation incorrecte (Python \texttt{IndentationError}).
        \end{itemize}
        \emph{Message typique :} \texttt{error: expected `;' before `\}' token} (C) ou \texttt{SyntaxError: invalid syntax} (Python) ou \texttt{Uncaught SyntaxError} (JS).
    \item \textbf{Erreurs d'exécution (Runtime / Exceptions) :} Le programme démarre, mais \textbf{plante en cours de route} sur une instruction impossible.
        \begin{itemize}
            \item Division par zéro (\texttt{ZeroDivisionError}, \texttt{Floating point exception} / \texttt{Infinity} en JS).
            \item Accès index hors borne tableau/liste (\texttt{IndexError}, \texttt{Segmentation fault} / \texttt{undefined} en JS).
            \item Conversion invalide : \texttt{int("abc")} $\to$ \texttt{ValueError} / \texttt{NaN} en JS.
            \item Fichier introuvable (\texttt{FileNotFoundError}).
        \end{itemize}
        \emph{Message typique :} \texttt{Traceback (most recent call last): ... ZeroDivisionError: division by zero}.
    \item \textbf{Erreurs logiques (Sémantiques) :} Le programme \textbf{tourne sans planter}, mais produit un \textbf{mauvais résultat}. Le code est syntaxiquement correct, l'algorithme est faux.
        \begin{itemize}
            \item Condition inversée : \texttt{if (x < 5)} au lieu de \texttt{>}.
            \item Boucle \texttt{while} qui ne s'arrête jamais (condition de sortie jamais vraie).
            \item Variable mal initialisée (somme = 1 au lieu de 0).
            \item Priorité d'opérateurs oubliée : \texttt{a + b * c} vs \texttt{(a + b) * c}.
        \end{itemize}
        \emph{Symptôme :} « Mon programme affiche 15 au lieu de 20. »
\end{enumerate}
\end{propriete}

\begin{methode}[Débogage basique : la démarche en 3 étapes]
\begin{enumerate}
    \item \textbf{Lire le message d'erreur (la trace).} Ne pas paniquer. Repérer : \textbf{Numéro de ligne}, \textbf{Type d'erreur} (\texttt{SyntaxError}, \texttt{ZeroDivisionError}, \texttt{NameError}, \texttt{ReferenceError}), \textbf{Description}. En C, la première erreur cascade souvent : corrige-la, recompile, les suivantes disparaissent souvent.
    \item \textbf{Vérifier la syntaxe locale (autour de la ligne).} Deux-points \texttt{:} ? Parenthèses \texttt{() \{\}} équilibrées ? Indentation cohérente (Python) ? Point-virgule \texttt{;} (C, JS) ? Guillemets fermés ? Noms de variables exacts (casse respectée : \texttt{sommeTotale} $\neq$ \texttt{sommetotale}) ?
    \item \textbf{Afficher les valeurs intermédiaires (\texttt{print} / \texttt{printf} / \texttt{console.log}).} C'est la \textbf{technique n°1} : insère des \texttt{print("DEBUG: x =", x)} avant la ligne qui plante ou donne un mauvais résultat. Suit le flux d'exécution. Retire-les une fois le bug corrigé.
\end{enumerate}
\end{methode}

\begin{experience}[TP guidé : Chasser les 3 types d'erreurs]
\textbf{Objectif :} Provocer, observer, corriger une erreur de compilation, une erreur d'exécution, une erreur logique.
\textbf{Outils :} Éditeur de code (VS Code, Geany, Thonny, ou replit.com en ligne) + Python 3 \textbf{et} Compilateur C (gcc via terminal ou IDE Code::Blocks) \textbf{et} Navigateur/Node.js pour JS.

\textbf{Étape 1 — Erreur de compilation (Python)} \\
Copie ce code, sauvegarde \texttt{erreur\_syntaxe.py}, lance-le.
\begin{lstlisting}[language=Python]
x = 10
if x > 5      # Oublié les deux-points :
    print("Grand")
else
    print("Petit")
\end{lstlisting}
\textbf{Observation :} \texttt{SyntaxError: expected ':'} sur la ligne 2. \textbf{Correction :} Ajoute \texttt{:} après \texttt{if x > 5} et après \texttt{else}.

\textbf{Étape 2 — Erreur d'exécution (C)} \\
Fichier \texttt{div0.c} :
\begin{lstlisting}[language=C]
#include <stdio.h>
int main() {
    int a = 10, b = 0;
    int res = a / b;   // Division par zéro !
    printf("Resultat : %d\n", res);
    return 0;
}
\end{lstlisting}
Compile : \texttt{gcc div0.c -o div0}. Lance : \texttt{./div0}.
\textbf{Observation :} \texttt{Erreur arithmétique (core dumped)} ou \texttt{Floating point exception}. Pas de message Python-style, le programme s'arrête brutalement. \textbf{Correction :} Tester \texttt{if (b != 0)} avant de diviser.

\textbf{Étape 3 — Erreur logique (JavaScript)} \\
Fichier \texttt{moyenne\_bug.js} : Calculer la moyenne de 3 notes.
\begin{lstlisting}[language=JavaScript]
let n1 = 12, n2 = 14, n3 = 16;
let moyenne = n1 + n2 + n3 / 3;   // Piège priorité : division avant addition !
console.log("Moyenne :", moyenne);  // Affiche 24.66... au lieu de 14
\end{lstlisting}
\textbf{Observation :} Programme tourne, affiche un nombre, mais \textbf{faux}. \textbf{Débogage (Méthode « Débogage basique » étape 3) :}
\begin{lstlisting}[language=JavaScript]
console.log("DEBUG: n1+n2+n3 =", n1+n2+n3);  // 42
console.log("DEBUG: n3/3 =", n3/3);          // 5.33
let moyenne = (n1 + n2 + n3) / 3;            // Parentheses obligatoires
console.log("Moyenne corrigée :", moyenne);  // 14
\end{lstlisting}
\textbf{Interprétation :} L'erreur logique est la plus sournoise : pas de message d'erreur, il faut \emph{tester avec des valeurs connues} (jeu de test) et \emph{afficher les étapes}.
\end{experience}

\courssub{Bonnes pratiques : écrire du code lisible et maintenable}

Un programme se lit \textbf{beaucoup plus souvent qu'il ne s'écrit} (relecture, correction, travail en équipe, reprise dans 6 mois). Le \emph{style} n'est pas cosmétique : c'est de la \textbf{prévention de bugs}.

\begin{aretenir}[Règles d'or du style (Clean Code niveau 3\up{me})]
\begin{enumerate}
    \item \textbf{Indentation rigoureuse :} 4 espaces (Python, standard PEP 8 ; JS standard) ou 1 tabulation (C). \textbf{Jamais} mélange espaces/tabs. Un bloc \texttt{SI}/\texttt{POUR}/\texttt{FONCTION} \emph{doit} être visuellement décalé.
    \item \textbf{Noms explicites (self-documenting code) :} \texttt{somme\_totale}, \texttt{note\_max}, \texttt{calculer\_moyenne()} au lieu de \texttt{x}, \texttt{n}, \texttt{f()}. \textbf{Convention :} \texttt{snake\_case} (Python, C) : \texttt{ma\_variable}, \texttt{ma\_fonction()}. \texttt{camelCase} (Java, JS) : \texttt{maVariable}, \texttt{maFonction()}. \textbf{Anglais recommandé} pour les noms (standard international), français accepté au BEPC si cohérent.
    \item \textbf{Commentaires utiles :} Expliquer \textbf{le « pourquoi »} (intention, contrainte métier, astuce), pas le « quoi » (le code le dit).
        \begin{lstlisting}[language=Python]
# MAUVAIS : x = x + 1  # On ajoute 1 à x
# BON :  x = x + 1     // Passage au candidat suivant (tour de boucle)
        \end{lstlisting}
    \item \textbf{Constantes nommées :} \texttt{TAUX\_TVA = 19.25} (majuscules + underscore) au lieu de \texttt{0.1925} « magique » au milieu du code.
    \item \textbf{Une instruction par ligne :} Pas de \texttt{x=1; y=2; if x>y: print("ok")}.
    \item \textbf{Limite de ligne :} $\le$ 80-100 caractères (lisibilité sans défilement horizontal).
\end{enumerate}
\end{aretenir}

\begin{attention}[Le piège de l'indentation en Python vs C / JavaScript]
En \textbf{Python}, l'indentation \textbf{définit la structure} du programme.
\begin{lstlisting}[language=Python]
if x > 0:
    print("Positif")
    print("Traitement suite")  # Dans le SI
print("Fin")                   # Hors du SI
\end{lstlisting}
Si tu décalages mal \texttt{print("Traitement suite")}, il sort du \texttt{if} $\to$ logique changée $\to$ \textbf{bug silencieux (erreur logique)}.

En \textbf{C} et \textbf{JavaScript}, les accolades \texttt{\{\}} définissent les blocs. L'indentation est \emph{ignorée} par le compilateur/interpréteur.
\begin{lstlisting}[language=C]
if (x > 0) {
    printf("Positif\n");
    printf("Traitement suite\n"); // Dans le SI grâce aux accolades
}
printf("Fin\n");                  // Hors du SI
\end{lstlisting}
\textbf{Conséquence :} En C/JS, on \textbf{peut} écrire tout sur une ligne (compilé/exécuté), mais c'est illisible. En Python, on \textbf{ne peut pas} : l'indentation \emph{est} la syntaxe. \textbf{Adopte l'habitude universelle :} \textbf{toujours mettre les accolades \{\}} même pour une seule instruction, et \textbf{toujours indenter} correctement.
\end{attention}

\begin{exoresolu}[De l'algorithme au code C complet : Calcul de factorielle]
\textbf{Énoncé :} Écrire l'algorithme (pseudo-code) puis le programme C qui demande un entier $n \ge 0$ et affiche $n!$ (factorielle). Rappel : $0! = 1$, $n! = 1 \times 2 \times \dots \times n$. Gérer le cas $n < 0$ (erreur).

\tcblower

\textbf{1. Algorithme (Pseudo-code)}
\begin{lstlisting}[caption=Algorithme Factorielle]
Algorithme Factorielle
Variables n, i, fact : Entier
Debut
    Afficher "Entier n (>= 0) : "
    Lire n
    Si n < 0 Alors
        Afficher "Erreur : n doit etre >= 0"
    Sinon
        fact <- 1
        Pour i de 1 a n Faire
            fact <- fact * i
        Fin Pour
        Afficher "Factorielle de ", n, " = ", fact
    Fin Si
Fin
\end{lstlisting}

\textbf{2. Programme C (avec commentaires)}
\begin{lstlisting}[language=C, caption=factorielle.c]
#include <stdio.h>
#include <stdlib.h>   // Pour EXIT_SUCCESS / EXIT_FAILURE

int main(void) {
    int n;           // Entier saisi
    long long fact = 1; // long long pour éviter débordement rapide (13! > 2^31)
    int i;

    printf("Entier n (>= 0) : ");
    if (scanf("%d", &n) != 1) { // Vérification lecture correcte
        fprintf(stderr, "Erreur de lecture.\n");
        return EXIT_FAILURE;
    }

    if (n < 0) {
        printf("Erreur : n doit etre >= 0\n");
        return EXIT_FAILURE; // Code retour non-nul = erreur
    }

    // Boucle factorielle : i de 1 à n inclus
    for (i = 1; i <= n; i++) {
        fact = fact * i; // ou fact *= i;
    }

    printf("Factorielle de %d = %lld\n", n, fact);
    return EXIT_SUCCESS; // 0 = succès
}
\end{lstlisting}

\textbf{3. Compilation et tests (Terminal Linux/Mac/WSL)}
\begin{lstlisting}[language=bash, caption=Session terminal]
$ gcc -Wall -Wextra -std=c11 factorielle.c -o factorielle
$ ./factorielle
Entier n (>= 0) : 5
Factorielle de 5 = 120
$ ./factorielle
Entier n (>= 0) : 0
Factorielle de 0 = 1
$ ./factorielle
Entier n (>= 0) : -3
Erreur : n doit etre >= 0
$ ./factorielle
Entier n (>= 0) : 20
Factorielle de 20 = 2432902008176640000
\end{lstlisting}
\textbf{Analyse :} Options \texttt{-Wall -Wextra} activent \textbf{tous les avertissements} (variables non utilisées, comparaisons signées/non-signées, etc.) : \textbf{habitude pro à prendre dès maintenant}. \texttt{-std=c11} fixe la norme C11. \texttt{long long} (entier 64 bits) permet d'aller jusqu'à $20!$ ; au-delà, débordement (besoin bibliothèques grands entiers).
\end{exoresolu}

\begin{savaistu}[Histoire : Grace Hopper et le premier compilateur (1952)]
Avant 1952, programmer signifiait câbler des fils ou écrire en code machine/octal. \textbf{Grace Hopper} (amirale US Navy, pionnière informatique) a développé le \textbf{A-0 System}, le premier \textbf{compilateur} : il traduisait des notations mathématiques symboliques en code machine pour l'UNIVAC I. On lui a dit : « Les ordinateurs ne font que de l'arithmétique, ils ne peuvent pas "programmer" ». Elle a répondu en créant \textbf{FLOW-MATIC}, premier langage utilisant des mots anglais (\texttt{READ}, \texttt{WRITE}, \texttt{IF}), ancêtre direct du \textbf{COBOL} (1959), encore utilisé aujourd'hui dans la banque mondiale. Elle a aussi popularisé le terme \textbf{« bug »} (papillon de nuit coincé dans un relais du Harvard Mark II, 1947). \emph{Moralité :} La compilation, c'est l'invention qui a permis à l'informatique de sortir des laboratoires pour devenir un outil universel.
\end{savaistu}

\begin{exercice}[Entraînement : Traduction, Erreurs, Style]
\begin{enumerate}
    \item \textbf{Traduction pseudo-code $\to$ JavaScript :} Écris le programme JavaScript (Node.js) correspondant à cet algorithme :
    \begin{lstlisting}
    Algorithme Signe
    Variable x : Reel
    Debut
        Afficher "Nombre : "
        Lire x
        Si x > 0 Alors Afficher "Positif"
        Sinon Si x < 0 Alors Afficher "Negatif"
        Sinon Afficher "Nul"
        Fin Si
    Fin
    \end{lstlisting}
    \item \textbf{Chasse aux erreurs (C) :} Ce code C contient \textbf{4 anomalies} (1 avertissement compilation / logique, 1 exécution, 1 logique, 1 style). Identifie-les et corrige-les.
    \begin{lstlisting}[language=C]
    #include <stdio.h>
    int main() {
        int a = 10, b = 0;
        printf("Division : %d / %d = ", a, b);
        printf("%d\n", a / b);      // Ligne 5
        if (a = 5) {                // Ligne 6
            printf("a vaut 5\n");
        }
        return 0;
    }
    \end{lstlisting}
    \item \textbf{Style :} Réécris ce fragment Python en respectant les conventions (noms, indentation, commentaire utile, constante).
    \begin{lstlisting}[language=Python]
    p=19.25;ht=150;ttc=ht*(1+p/100);print(ttc)
    \end{lstlisting}
\end{enumerate}
\end{exercice}

\begin{corrige}[Exercice n]
\begin{enumerate}
    \item \textbf{JavaScript (Node.js) :}
    \begin{lstlisting}[language=JavaScript]
    const readline = require('readline').createInterface({
        input: process.stdin,
        output: process.stdout
    });

    readline.question("Nombre : ", (input) => {
        const x = parseFloat(input);
        if (x > 0) {
            console.log("Positif");
        } else if (x < 0) {          // else if = Sinon Si
            console.log("Negatif");
        } else {
            console.log("Nul");
        }
        readline.close();
    });
    \end{lstlisting}
    \item \textbf{Anomalies C :}
    \begin{enumerate}
        \item [Ligne 5 - Exécution] \texttt{a / b} : \textbf{Division par zéro} ($b=0$). Le programme plante (\texttt{Floating point exception}). \textbf{Correction :} Tester \texttt{if (b != 0)} avant de diviser, sinon afficher erreur.
        \item [Ligne 6 - Compilation/Logique] \texttt{if (a = 5)} : \texttt{=} est une \textbf{affectation} (met 5 dans \texttt{a}, l'expression vaut 5 $\to$ vrai). \texttt{gcc -Wall} émet un avertissement \texttt{suggest parentheses around assignment used as truth value}. Il faut \texttt{==} (comparaison). \textbf{Correction :} \texttt{if (a == 5)}.
        \item [Logique] Le programme affiche « a vaut 5 » car l'affectation \texttt{a=5} a changé la valeur de \texttt{a} (était 10).
        \item [Style] \texttt{int main()} $\to$ préférer \texttt{int main(void)} pour prototypage strict. Ajouter \texttt{\#include <stdlib.h>} si utilisation de \texttt{EXIT\_SUCCESS/FAILURE}.
    \end{enumerate}
    Code corrigé :
    \begin{lstlisting}[language=C]
    #include <stdio.h>
    #include <stdlib.h>

    int main(void) {
        int a = 10, b = 2; // b != 0 pour test division
        if (b != 0) {
            printf("Division : %d / %d = %d\n", a, b, a / b);
        } else {
            printf("Erreur : division par zero.\n");
        }
        if (a == 5) { // Comparaison ==
            printf("a vaut 5\n");
        }
        return EXIT_SUCCESS;
    }
    \end{lstlisting}
    \item \textbf{Style Python (PEP 8) :}
    \begin{lstlisting}[language=Python]
    # Calcul TTC à partir du HT et taux TVA
    TAUX_TVA = 19.25      # Constante : majuscules + underscore
    prix_ht = 150.0       # Noms explicites, snake_case, type float explicite
    prix_ttc = prix_ht * (1 + TAUX_TVA / 100)
    print(f"Prix TTC : {prix_ttc:.2f}") # f-string formatée (2 décimales)
    \end{lstlisting}
\end{enumerate}
\end{corrige}

\coursec{Les briques de base : Variables, Types et Structures de contrôle}

\courssub{Introduction : le vocabulaire de l'algorithme}

Dans la section précédente, nous avons vu qu'un algorithme s'écrit dans un langage précis (pseudo-code) avant d'être traduit en langage machine. Mais quel que soit le langage, \textbf{tout programme repose sur les mêmes trois piliers} : stocker des données (variables), décider quoi faire selon des conditions (sélection), et répéter des actions (itération). C'est le « B-A-BA » de la programmation, valable du plus petit script au plus gros système d'exploitation.

Imagine que tu gères la caisse d'une tontine scolaire à Yaoundé. Tu dois : \textbf{mémoriser} le total des cotisations (variable), \textbf{tester} si un membre a payé (condition), \textbf{répéter} la lecture des 20 reçus (boucle). Sans ces trois briques, impossible d'automatiser la tâche. Maîtriser ces concepts, c'est acquérir le pouvoir de décomposer n'importe quel problème en instructions élémentaires que l'ordinateur saura exécuter.

\courssub{Les variables : des boîtes étiquetées en mémoire}

\begin{definition}[Variable]
Une \cle{variable} est une zone de la mémoire vive (RAM) identifiée par un \textbf{nom} (identificateur), capable de contenir \textbf{une seule valeur à la fois} d'un \textbf{type} donné. Sa valeur peut changer au cours de l'exécution du programme (d'où le nom « variable »).
\end{definition}

Contrairement aux mathématiques où $x = 5$ est une vérité immuable, en algorithmique \texttt{x <- 5} signifie : « réserve une case mémoire nommée \texttt{x} et dépose-y la valeur 5 ». Plus tard, \texttt{x <- x + 1} remplace l'ancienne valeur par la nouvelle.

\begin{propriete}[Règles de nommage et déclaration]
\begin{itemize}
    \item \textbf{Nom} : commence par une lettre (pas de chiffre), sans espace, sans accent (ex. \texttt{sommeTotale}, \texttt{nbEleves}, \texttt{\_note}). Sensible à la casse en C/Python (\texttt{Note} $\neq$ \texttt{note}), mais généralement insensible en pseudo-code MINESEC (on évite toutefois les ambiguïtés).
    \item \textbf{Déclaration} (pseudo-code) : \texttt{Variable x : ENTIER} — réserve l'espace et impose le type.
    \item \textbf{Affectation} : \texttt{x <- 10} (flèche vers la variable). La lecture se fait de droite à gauche : on évalue l'expression à droite, puis on stocke dans la variable à gauche.
\end{itemize}
\end{propriete}

\begin{propriete}[Initialisation obligatoire]
Toute variable doit recevoir une \textbf{première valeur} (initialisation) \emph{avant} d'être utilisée dans un calcul. Sinon, elle contient une « valeur poubelle » (résidu mémoire) et le résultat devient imprévisible.
\end{propriete}

\begin{attention}[Affectation vs Égalité — Piège n°1]
En algorithmique et en C : \texttt{=} est l'affectation (\texttt{x = 5} stocke 5), \texttt{==} est le test d'égalité (\texttt{if (x == 5)}). En pseudo-code MINESEC : \texttt{x <- 5} (affectation) et \texttt{x = 5} (test). \textbf{Ne jamais confondre les deux} : \texttt{x = x + 1} en C incrémente ; en maths c'est faux ($0=1$).
\end{attention}

\courssub{Les types de base : quel genre de valeur dans la boîte ?}

Le type détermine : l'espace mémoire réservé, les valeurs autorisées, et les opérations possibles.

\begin{center}
\begin{tabularx}{\linewidth}{|l|X|X|l|}
\hline
\rowcolor{popblueL}
\textbf{Type} & \textbf{Valeurs typiques} & \textbf{Opérations principales} & \textbf{Exemple déclaration} \\
\hline
ENTIER & Nombres sans virgule : \texttt{-12, 0, 42, 2025} & \texttt{+, -, *, / (division entière), MOD, DIV} & \texttt{Variable age : ENTIER} \\
\hline
REEL & Nombres à virgule : \texttt{3.14, -0.5, 12.0} & \texttt{+, -, *, / (division réelle)} & \texttt{Variable moyenne : REEL} \\
\hline
CARACTERE & Un seul symbole : \texttt{'A', '7', '@', 'é'} & Concaténation, comparaison & \texttt{Variable sexe : CARACTERE} \\
\hline
CHAINE & Suite de caractères : \texttt{"Yaoundé", "Classe 3ème"} & Concaténation (+), longueur, extraction & \texttt{Variable nom : CHAINE} \\
\hline
BOOLEEN & \texttt{VRAI} ou \texttt{FAUX} & \texttt{ET, OU, NON} & \texttt{Variable aPaye : BOOLEEN} \\
\hline
\end{tabularx}
\end{center}

\begin{exemplebox}[Choisir le bon type pour la tontine]
\begin{itemize}
    \item Nombre de membres (20) $\rightarrow$ \texttt{ENTIER}
    \item Montant d'une cotisation (1500 FCFA) $\rightarrow$ \texttt{ENTIER} (pas de centimes) ou \texttt{REEL} si précision décimale
    \item Nom du membre $\rightarrow$ \texttt{CHAINE}
    \item A-t-il payé ce mois-ci ? $\rightarrow$ \texttt{BOOLEEN}
    \item Moyenne des cotisations $\rightarrow$ \texttt{REEL} (division peut donner décimal)
\end{itemize}
\end{exemplebox}

\courssub{Structure SÉQUENCE : l'ordre naturel}

Par défaut, les instructions s'exécutent \textbf{une après l'autre}, de la première à la dernière. C'est la structure la plus simple : pas de choix, pas de retour en arrière.

\begin{popfigure}
\begin{tikzpicture}[
    node distance=1.2cm and 1.5cm,
    rect/.style={rectangle, draw=popblue, thick, fill=popblueL, minimum width=2.8cm, minimum height=0.9cm, rounded corners=4pt},
    arrow/.style={-Stealth, thick, popdark},
    startend/.style={rectangle, draw=popgreen, thick, fill=popgreenL, minimum width=2.2cm, minimum height=0.8cm, rounded corners=20pt}
]
\node[startend] (debut) {DÉBUT};
\node[rect, below=of debut] (i1) {Instruction 1};
\node[rect, below=of i1] (i2) {Instruction 2};
\node[rect, below=of i2] (i3) {Instruction 3};
\node[startend, below=of i3] (fin) {FIN};
\draw[arrow] (debut) -- (i1);
\draw[arrow] (i1) -- (i2);
\draw[arrow] (i2) -- (i3);
\draw[arrow] (i3) -- (fin);
\end{tikzpicture}
\legende{Diagramme de flux (flowchart) d'une structure SÉQUENCE : exécution linéaire.}
\end{popfigure}

\begin{lstlisting}[caption=Exemple de séquence en pseudo-code]
// Calcul du périmètre d'un rectangle
Variable longueur, largeur, perimetre : ENTIER
DEBUT
    longueur <- 12      // Instruction 1
    largeur <- 7        // Instruction 2
    perimetre <- 2 * (longueur + largeur)  // Instruction 3
    AFFICHER "Périmètre = ", perimetre     // Instruction 4
FIN
\end{lstlisting}

\courssub{Structure SÉLECTION (Conditionnelle) : SI … ALORS … SINON}

La sélection permet d'exécuter \textbf{un bloc d'instructions seulement si une condition est vraie}, et éventuellement un autre bloc si elle est fausse. C'est la prise de décision.

\begin{popfigure}
\begin{tikzpicture}[
    node distance=1.1cm and 1.8cm,
    rect/.style={rectangle, draw=popblue, thick, fill=popblueL, minimum width=2.6cm, minimum height=0.85cm, rounded corners=4pt},
    diam/.style={diamond, draw=poporange, thick, fill=poporangeL, minimum width=2.8cm, minimum height=1.3cm, aspect=2},
    arrow/.style={-Stealth, thick, popdark},
    startend/.style={rectangle, draw=popgreen, thick, fill=popgreenL, minimum width=2.2cm, minimum height=0.8cm, rounded corners=20pt},
    label/.style={font=\small, text=popdark}
]
\node[startend] (debut) {DÉBUT};
\node[diam, below=of debut] (cond) {Condition ?};
\node[rect, below left=of cond] (alors) {Bloc ALORS};
\node[rect, below right=of cond] (sinon) {Bloc SINON};
\node[rect, below=1.5cm of cond] (suite) {Instruction suivante};
\node[startend, below=of suite] (fin) {FIN};
\draw[arrow] (debut) -- (cond);
\draw[arrow] (cond) -- node[label, pos=0.6, left] {VRAI} (alors);
\draw[arrow] (cond) -- node[label, pos=0.6, right] {FAUX} (sinon);
\draw[arrow] (alors) -- (suite);
\draw[arrow] (sinon) -- (suite);
\draw[arrow] (suite) -- (fin);
\end{tikzpicture}
\legende{Diagramme de flux d'une structure SI … ALORS … SINON. Le losange = test ; deux branches VRAI/FAUX.}
\end{popfigure}

\begin{propriete}[Syntaxe pseudo-code MINESEC]
\begin{lstlisting}
SI <condition> ALORS
    <instructions si VRAI>
SINON
    <instructions si FAUX>
FIN SI
\end{lstlisting}
La partie \texttt{SINON ...} est \textbf{optionnelle} (SI simple). Les blocs peuvent contenir d'autres SI (imbrication).
\end{propriete}

\begin{exemplebox}[Admission au BEPC — SI imbriqué]
\begin{lstlisting}
SI moyenne >= 10 ALORS
    SI noteFrancais >= 10 ET noteMaths >= 10 ALORS
        AFFICHER "Admis au BEPC"
    SINON
        AFFICHER "Ajourné (épreuve éliminatoire)"
    FIN SI
SINON
    AFFICHER "Ajourné (moyenne < 10)"
FIN SI
\end{lstlisting}
\end{exemplebox}

\begin{attention}[Pièges classiques de la sélection]
\begin{itemize}
    \item Oublier \texttt{FIN SI} $\rightarrow$ erreur de structure.
    \item Confondre \texttt{=} (affectation) et \texttt{==} (test) dans la condition (en C) ; en pseudo-code, confondre \texttt{<-} et \texttt{=}.
    \item Écrire \texttt{SI x > 5 ET < 10} au lieu de \texttt{SI x > 5 ET x < 10}.
    \item Ordonner les tests du plus restrictif au plus large (sinon le premier « attrape » tout).
\end{itemize}
\end{attention}

\courssub{Structure ITÉRATION (Boucles) : répéter sans recopier}

Une boucle exécute un bloc d'instructions \textbf{plusieurs fois}. Deux grandes familles selon qu'on connaît le nombre de tours à l'avance.

\begin{popfigure}
\begin{tikzpicture}[
    node distance=1.1cm and 1.8cm,
    rect/.style={rectangle, draw=popblue, thick, fill=popblueL, minimum width=2.6cm, minimum height=0.85cm, rounded corners=4pt},
    diam/.style={diamond, draw=poppurple, thick, fill=poppurpleL, minimum width=2.8cm, minimum height=1.3cm, aspect=2},
    arrow/.style={-Stealth, thick, popdark},
    startend/.style={rectangle, draw=popgreen, thick, fill=popgreenL, minimum width=2.2cm, minimum height=0.8cm, rounded corners=20pt},
    label/.style={font=\small, text=popdark}
]
\node[startend] (debut) {DÉBUT};
\node[diam, below=of debut] (cond) {Condition ?};
\node[rect, below=of cond] (corps) {Corps de boucle};
\node[startend, below=of corps] (fin) {FIN};
\draw[arrow] (debut) -- (cond);
\draw[arrow] (cond) -- node[label, left] {VRAI} (corps);
\draw[arrow] (corps) |- ++(-2.5,0) |- (cond);
\draw[arrow] (cond) -- node[label, right] {FAUX} (fin);
\end{tikzpicture}
\legende{Diagramme de flux d'une boucle TANT QUE : test \emph{avant} le corps. Flèche de retour (VRAI) vers le test.}
\end{popfigure}

\begin{methode}[Choisir sa boucle : POUR vs TANT QUE]
\begin{enumerate}
    \item \textbf{Je connais le nombre exact de tours à l'avance} (ex. 20 membres, 12 mois, 100 lignes) $\rightarrow$ \cle{POUR} (borne fixe).
    \item \textbf{Je m'arrête quand une condition devient fausse} (ex. « tant que caisse < objectif », « tant que l'utilisateur ne tape pas 'Q' ») $\rightarrow$ \cle{TANT QUE} (borne inconnue).
    \item \textbf{Attention} : en TANT QUE, la condition doit \textbf{éventuellement devenir fausse} (variable modifiée dans le corps), sinon \textbf{boucle infinie} $\rightarrow$ plantage.
\end{enumerate}
\end{methode}

\begin{propriete}[Syntaxe pseudo-code — Boucle POUR]
\begin{lstlisting}
POUR i DE <début> A <fin> [PAS <pas>] FAIRE
    <corps de boucle>
FIN POUR
\end{lstlisting}
\begin{itemize}
    \item \texttt{i} : variable de boucle (compteur), souvent \texttt{ENTIER}.
    \item \texttt{PAS} optionnel (par défaut +1). Peut être négatif pour décompter.
\end{itemize}
\end{propriete}

\begin{propriete}[Syntaxe pseudo-code — Boucle TANT QUE]
\begin{lstlisting}
TANT QUE <condition> FAIRE
    <corps de boucle>
FIN TANT QUE
\end{lstlisting}
Test \emph{avant} chaque tour. Si condition fausse au départ $\rightarrow$ 0 tour.
\end{propriete}

\begin{exemplebox}[Calcul de la moyenne d'une tontine — 20 membres]
\begin{lstlisting}
Variable somme, montant, moyenne : REEL
Variable i : ENTIER

DEBUT
    somme <- 0
    POUR i DE 1 A 20 FAIRE
        AFFICHER "Cotisation membre ", i, " : "
        LIRE montant
        somme <- somme + montant
    FIN POUR

    moyenne <- somme / 20
    AFFICHER "Moyenne = ", moyenne
FIN
\end{lstlisting}
\textbf{Trace} (premiers tours) : $i=1$, lu 1500 $\rightarrow$ somme=1500 ; $i=2$, lu 2000 $\rightarrow$ somme=3500 ; … $i=20$, lu 1800 $\rightarrow$ somme=38500 ; moyenne = 1925.
\end{exemplebox}

\begin{exemplebox}[Saisie jusqu'à mot de passe correct — TANT QUE]
\begin{lstlisting}
Variable mdpSaisi : CHAINE
Variable ok : BOOLEEN

DEBUT
    ok <- FAUX
    TANT QUE ok = FAUX FAIRE
        AFFICHER "Mot de passe : "
        LIRE mdpSaisi
        SI mdpSaisi = "OnBuch2025" ALORS
            ok <- VRAI
        SINON
            AFFICHER "Incorrect, réessayez."
        FIN SI
    FIN TANT QUE
    AFFICHER "Accès autorisé."
FIN
\end{lstlisting}
Ici, le nombre de tours est \textbf{inconnu} à l'avance (dépend de l'utilisateur).
\end{exemplebox}

\begin{attention}[Boucle infinie — Danger n°1]
\begin{lstlisting}
Variable i : ENTIER
DEBUT
    i <- 1
    TANT QUE i <= 10 FAIRE
        AFFICHER i
        // OUBLI : i <- i + 1
    FIN TANT QUE
FIN
\end{lstlisting}
$\rightarrow$ Affiche « 1 » indéfiniment. \textbf{Toujours} modifier la variable de condition dans le corps (incrémenter compteur, lire nouvelle saisie, etc.).
\end{attention}

\courssub{Les opérateurs : calculer, comparer, combiner}

\begin{center}
\begin{tabularx}{\linewidth}{|l|l|X|l|}
\hline
\rowcolor{popblueL}
\textbf{Famille} & \textbf{Opérateurs} & \textbf{Signification} & \textbf{Exemple (VRAI/FAUX)} \\
\hline
Arithmétiques & \texttt{+ - * /} & Addition, soustraction, multiplication, division & \texttt{7 / 2 = 3.5} (REEL) \\
 & \texttt{DIV} & Division entière (quotient) & \texttt{7 DIV 2 = 3} \\
 & \texttt{MOD} & Modulo (reste division entière) & \texttt{7 MOD 2 = 1} \\
\hline
Relationnels & \texttt{=} & Égal à & \texttt{5 = 5 $\rightarrow$ VRAI} \\
 & \texttt{<>} & Différent de & \texttt{5 <> 3 $\rightarrow$ VRAI} \\
 & \texttt{< > <= >=} & Inférieur, supérieur, $\le$, $\ge$ & \texttt{10 >= 10 $\rightarrow$ VRAI} \\
\hline
Logiques & \texttt{ET} & Conjonction : VRAI si \textbf{les deux} VRAI & \texttt{VRAI ET FAUX $\rightarrow$ FAUX} \\
 & \texttt{OU} & Disjonction : VRAI si \textbf{au moins un} VRAI & \texttt{VRAI OU FAUX $\rightarrow$ VRAI} \\
 & \texttt{NON} & Négation : inverse & \texttt{NON VRAI $\rightarrow$ FAUX} \\
\hline
\end{tabularx}
\end{center}

\begin{propriete}[Priorité des opérateurs (du plus fort au plus faible)]
\[
\texttt{()} \;\rightarrow\; \texttt{NON} \;\rightarrow\; \texttt{* / DIV MOD} \;\rightarrow\; \texttt{+ -} \;\rightarrow\; \texttt{< > <= >= = <>} \;\rightarrow\; \texttt{ET} \;\rightarrow\; \texttt{OU}
\]
\begin{itemize}
    \item Les parenthèses forcent l'évaluation (comme en maths).
    \item \texttt{NON} s'applique à l'opérande immédiat : \texttt{NON a ET b} = \texttt{(NON a) ET b}.
    \item En cas de doute : \textbf{parenthéser} pour la lisibilité.
\end{itemize}
\end{propriete}

\begin{savaistu}[Le modulo (MOD) : le couteau suisse des programmeurs]
L'opérateur \texttt{MOD} (modulo) donne le \textbf{reste de la division entière}. Applications courantes :
\begin{itemize}
    \item \textbf{Pair / Impair} : \texttt{x MOD 2 = 0} $\iff$ pair.
    \item \textbf{Dernier chiffre} : \texttt{x MOD 10} donne l'unité (ex. 2025 MOD 10 = 5).
    \item \textbf{Cycles} : jour de la semaine = \texttt{(jourDepart + nbJours) MOD 7}.
    \item \textbf{Conversion base} : extraire bits (binaire), chiffres (hexadécimal).
\end{itemize}
C'est l'outil de base pour tout ce qui est « cyclique » ou « périodique ».
\end{savaistu}

\courssub{Exemple résolu complet : Gestion d'un mini-stock de cantine}

\begin{exoresolu}[Stock cantine — Variables, SI, POUR, TANT QUE]
\textbf{Énoncé.} La cantine du lycée gère 5 produits. Pour chaque produit, on connaît le stock initial. Le programme doit : (1) afficher le stock initial de chaque produit ; (2) simuler 3 jours de vente : chaque jour, on saisit le produit vendu (1 à 5) et la quantité ; on met à jour le stock si quantité $\le$ stock, sinon message « Rupture » ; (3) afficher le stock final.

\textbf{Analyse.}
\begin{itemize}
    \item 5 produits $\rightarrow$ 5 variables distinctes \texttt{stock1..5} (ENTIER) (les tableaux ne sont pas au programme de 3ème).
    \item 3 jours fixes $\rightarrow$ boucle \texttt{POUR jour DE 1 A 3}.
    \item Choix produit + quantité $\rightarrow$ \texttt{LIRE} dans la boucle.
    \item Test stock suffisant $\rightarrow$ \texttt{SI qte <= stock ALORS ... SINON ...}.
\end{itemize}

\textbf{Algorithme (pseudo-code MINESEC).}
\begin{lstlisting}
ALGORITHME StockCantine
VARIABLES
    stock1, stock2, stock3, stock4, stock5 : ENTIER
    jour, choix, qte : ENTIER
DEBUT
    // 1. Initialisation (valeurs exemple)
    stock1 <- 50; stock2 <- 30; stock3 <- 40; stock4 <- 25; stock5 <- 60

    // Affichage initial
    AFFICHER "Stock initial :"
    AFFICHER "Produit 1 : ", stock1
    AFFICHER "Produit 2 : ", stock2
    AFFICHER "Produit 3 : ", stock3
    AFFICHER "Produit 4 : ", stock4
    AFFICHER "Produit 5 : ", stock5

    // 2. Simulation 3 jours
    POUR jour DE 1 A 3 FAIRE
        AFFICHER "--- Jour ", jour, " ---"
        AFFICHER "Produit vendu (1-5) : "; LIRE choix
        AFFICHER "Quantité : "; LIRE qte

        // Mise à jour selon choix
        SI choix = 1 ALORS
            SI qte <= stock1 ALORS stock1 <- stock1 - qte
            SINON AFFICHER "Rupture produit 1" FIN SI
        SINON SI choix = 2 ALORS
            SI qte <= stock2 ALORS stock2 <- stock2 - qte
            SINON AFFICHER "Rupture produit 2" FIN SI
        SINON SI choix = 3 ALORS
            SI qte <= stock3 ALORS stock3 <- stock3 - qte
            SINON AFFICHER "Rupture produit 3" FIN SI
        SINON SI choix = 4 ALORS
            SI qte <= stock4 ALORS stock4 <- stock4 - qte
            SINON AFFICHER "Rupture produit 4" FIN SI
        SINON SI choix = 5 ALORS
            SI qte <= stock5 ALORS stock5 <- stock5 - qte
            SINON AFFICHER "Rupture produit 5" FIN SI
        SINON
            AFFICHER "Produit inexistant"
        FIN SI
    FIN POUR

    // 3. Stock final
    AFFICHER "=== Stock final ==="
    AFFICHER "P1:", stock1, " P2:", stock2, " P3:", stock3, " P4:", stock4, " P5:", stock5
FIN
\end{lstlisting}

\textbf{Test manuel (trace partielle).}
\begin{itemize}
    \item Initial : [50, 30, 40, 25, 60]
    \item Jour 1 : choix=2, qte=10 $\rightarrow$ stock2=20 (OK)
    \item Jour 2 : choix=4, qte=30 $\rightarrow$ 30 > 25 $\rightarrow$ « Rupture produit 4 », stock4 inchangé
    \item Jour 3 : choix=1, qte=50 $\rightarrow$ stock1=0 (OK)
    \item Final : [0, 20, 40, 25, 60]
\end{itemize}

\textbf{Remarque.} Ce code répète la logique pour chaque produit. Dans la prochaine section (modularité), nous verrons comment factoriser avec une \textbf{fonction} pour éviter la duplication.
\tcblower
\textbf{Solution détaillée.} Voir l'algorithme complet ci-dessus. Points clés validés :
\begin{itemize}
    \item Déclaration explicite de toutes variables avec types.
    \item Initialisation avant utilisation (propriété respectée).
    \item Boucle POUR pour nombre de jours connu (3).
    \item SI / SINON SI / SINON pour dispatcher selon le produit.
    \item Test de garde \texttt{qte <= stock} avant déduction (évite stock négatif).
    \item Affichages clairs pour traçabilité.
\end{itemize}
\end{exoresolu}

\courssub{Synthèse : les 3 piliers en une image mentale}

\begin{center}
\begin{tikzpicture}[
    node distance=1.5cm and 2cm,
    pillar/.style={rectangle, draw=popblue, very thick, fill=popblueL, minimum width=3.2cm, minimum height=2.2cm, rounded corners=6pt},
    title/.style={font=\bfseries, text=popblue},
    item/.style={font=\small, text=popdark}
]
\node[pillar] (var) at (0,0) {};
\node[title] at (var.north) [yshift=-0.3cm] {VARIABLES};
\node[item, align=center] at (var.center) {Nom + Type + Valeur\\Déclaration, Affectation\\Initialisation obligatoire};
\node[pillar] (sel) at (5,0) {};
\node[title] at (sel.north) [yshift=-0.3cm] {SÉLECTION};
\node[item, align=center] at (sel.center) {SI condition ALORS\\... SINON ... FIN SI\\Imbrication possible\\Opérateurs relationnels/logiques};
\node[pillar] (it) at (10,0) {};
\node[title] at (it.north) [yshift=-0.3cm] {ITÉRATION};
\node[item, align=center] at (it.center) {POUR (nb tours connu)\\TANT QUE (condition)\\Risque : boucle infinie\\Modulo (MOD) pour cycles};
\end{tikzpicture}
\end{center}

\begin{aretenir}[Ce qu'il faut retenir pour le BEPC]
\begin{enumerate}
    \item Une \textbf{variable} = nom + type + valeur changeante. Toujours l'initialiser.
    \item Quatre types de base : \texttt{ENTIER, REEL, CHAINE/CARACTERE, BOOLEEN}.
    \item \textbf{SÉQUENCE} = défaut (linéaire). \textbf{SÉLECTION} = \texttt{SI ... ALORS ... SINON ... FIN SI}. \textbf{ITÉRATION} = \texttt{POUR} (compteur) ou \texttt{TANT QUE} (condition).
    \item Opérateurs : arithmétiques (+, -, *, /, DIV, MOD), relationnels (=, <>, <, >, <=, >=), logiques (ET, OU, NON).
    \item Priorité : () $\rightarrow$ NON $\rightarrow$ * / DIV MOD $\rightarrow$ + - $\rightarrow$ relationnels $\rightarrow$ ET $\rightarrow$ OU.
    \item \texttt{MOD} : reste division entière $\rightarrow$ pair/impair, cycles, dernier chiffre.
    \item Pièges : \texttt{<-} vs \texttt{=}, boucle infinie, oubli \texttt{FIN SI}, condition mal parenthésée.
\end{enumerate}
\end{aretenir}

\begin{exercice}[Entraînement — Variables, Conditions, Boucles]
\begin{enumerate}
    \item \textbf{Type et initialisation.} Pour chaque donnée, donne le type le plus adapté et une initialisation correcte :
    \begin{itemize}
        \item (a) Nombre d'élèves dans la classe.
        \item (b) Prix unitaire d'un cahier (250 FCFA).
        \item (c) Nom de l'élève.
        \item (d) L'élève a-t-il rendu son livre ? (Oui/Non).
        \item (e) Moyenne de la classe (12,75).
    \end{itemize}
    \item \textbf{Trace d'algorithme.} Que affiche cet algorithme ?
    \begin{lstlisting}
    Variable x, y : ENTIER
    DEBUT
    x <- 3; y <- 5
    SI x + y > 7 ALORS
        x <- x * 2
    SINON
        y <- y - 1
    FIN SI
    AFFICHER x, y
    FIN
    \end{lstlisting}
    \item \textbf{Boucle POUR.} Écrire un algorithme qui calcule la factorielle $n! = 1 \times 2 \times \dots \times n$ pour un entier $n \ge 0$ saisi au clavier (rappel : $0! = 1$).
    \item \textbf{Boucle TANT QUE.} Écrire un algorithme qui demande des notes (sur 20) jusqu'à ce que l'utilisateur saisisse une note négative, puis affiche la moyenne des notes valides (si au moins une note).
    \item \textbf{Modulo.} Expliquer pourquoi \texttt{(annee MOD 4 = 0 ET annee MOD 100 <> 0) OU (annee MOD 400 = 0)} détecte une année bissextile.
\end{enumerate}
\end{exercice}

\begin{corrige}[Exercice 3.1]
\begin{enumerate}
    \item (a) ENTIER, \texttt{nbEleves <- 0} \quad (b) ENTIER ou REEL, \texttt{prix <- 250} \quad (c) CHAINE, \texttt{nom <- ""} \quad (d) BOOLEEN, \texttt{aRendu <- FAUX} \quad (e) REEL, \texttt{moyenne <- 0.0}
    \item $x=3, y=5 \rightarrow x+y=8 > 7$ VRAI $\rightarrow$ branche ALORS : $x \leftarrow 6$. Affichage : \texttt{6 5}.
    \item \textbf{Factorielle (POUR) :}
    \begin{lstlisting}
    Variable n, i, fact : ENTIER
    DEBUT
    AFFICHER "n = "; LIRE n
    fact <- 1
    POUR i DE 1 A n FAIRE
        fact <- fact * i
    FIN POUR
    AFFICHER n, "! = ", fact
    FIN
    \end{lstlisting}
    \item \textbf{Moyenne notes (TANT QUE) :}
    \begin{lstlisting}
    Variable note, somme : REEL
    Variable nb : ENTIER
    DEBUT
    somme <- 0; nb <- 0
    AFFICHER "Note (négative pour finir) : "; LIRE note
    TANT QUE note >= 0 FAIRE
        somme <- somme + note
        nb <- nb + 1
        AFFICHER "Note suivante : "; LIRE note
    FIN TANT QUE
    SI nb > 0 ALORS AFFICHER "Moyenne = ", somme / nb
    SINON AFFICHER "Aucune note saisie" FIN SI
    FIN
    \end{lstlisting}
    \item \textbf{Bissextile :} Règle grégorienne. Divisible par 4 SAUF siècles non divisibles par 400. \texttt{MOD 4 = 0} = multiple de 4. \texttt{MOD 100 <> 0} = pas un siècle. \texttt{MOD 400 = 0} = siècle multiple de 400 (ex. 2000 bissextile, 1900 non).
\end{enumerate}
\end{corrige}

\courssub{Transition vers la section 4}

Nous savons maintenant stocker, décider, répéter. Mais un programme de 500 lignes tout dans le \texttt{DEBUT...FIN} devient illisible. La section 4 — \textbf{Modularité : Fonctions et Procédures} — nous apprendra à \textbf{découper} le problème en sous-problèmes indépendants (diviser pour régner), à réutiliser du code, et à rendre l'algorithme clair comme de l'eau de source. C'est le passage du « script » au « logiciel structuré ».

\coursec{Modularité : Fonctions et Procédures (Diviser pour régner)}

Dans la section précédente, nous avons découvert les \cle{briques de base} de tout algorithme : les variables, les types de données et les structures de contrôle (séquence, alternative, boucles). Avec ces outils, tu es déjà capable d'écrire des algorithmes complets. Mais imagine maintenant que tu doives écrire le programme de gestion complet d'une boutique à Yaoundé : saisie des produits, calcul des prix, gestion du stock, impression des reçus, statistiques de vente... Ton algorithme ferait des centaines de lignes ! Comment s'y retrouver ? Comment corriger une erreur sans tout casser ? C'est exactement le problème que résout la \cle{modularité}.

\courssub{Le problème : un algorithme unique devient illisible}

\begin{exemplebox}[Une situation de la vie courante]
La cantine du collège veut un programme qui calcule la facture de chaque élève : prix du repas, réduction pour les boursiers, taxe, puis affichage du reçu. Si l'on écrit tout dans un seul bloc, on obtient un long algorithme où se mélangent calculs et affichages. Le jour où le taux de réduction change, il faut relire tout le programme pour trouver la bonne ligne... et corriger sans casser le reste !
\end{exemplebox}

Face à un problème complexe, l'être humain applique naturellement une stratégie très ancienne : \emph{diviser pour régner}. Au lieu d'attaquer le gros problème d'un coup, on le découpe en petits sous-problèmes simples, que l'on résout un par un. En algorithmique, chaque sous-problème est confié à un \cle{module}.

\begin{definition}[Module (ou sous-algorithme)]
Un \cle{module} (ou \cle{sous-algorithme}) est un morceau de programme nommé, qui réalise une tâche précise et qui peut être \emph{appelé} (utilisé) par le programme principal ou par un autre module. Il existe deux sortes de modules : les \cle{fonctions} et les \cle{procédures}.
\end{definition}

\begin{popfigure}
\begin{tikzpicture}[
  font=\small,
  main/.style={draw=popdark, fill=popblueL, rounded corners, thick, minimum width=6.2cm, minimum height=1.1cm, align=center},
  module/.style={draw=popdark, fill=popgreenL, rounded corners, thick, minimum width=4.6cm, minimum height=1.6cm, align=center},
  local/.style={draw=popmuted, dashed, fill=popcream, rounded corners, align=center, font=\scriptsize},
  fleche/.style={-{Stealth[length=3mm]}, thick, popdark},
  retour/.style={-{Stealth[length=3mm]}, thick, poporange, dashed}
]
\node[main, align=center] (pp){\textbf{Programme Principal}\\ \scriptsize (algorithme FactureCantine)};
\node[module, below left=1.9cm and -1.2cm of pp, align=center] (fonc){\textbf{FONCTION CalculerReduction()}\\ \scriptsize calcule et renvoie une valeur};
\node[module, below right=1.9cm and -1.2cm of pp, align=center] (proc){\textbf{PROCEDURE AfficherRecu()}\\ \scriptsize effectue une action (affichage)};
\node[local, below=0.15cm of fonc, align=center] (loc1){Variables locales :\\ reduction, tauxApplique};
\node[local, below=0.15cm of proc, align=center] (loc2){Variables locales :\\ ligne, totalFormate};
\draw[fleche] (pp.south) ++(-1.1,0) -- (fonc.north) node[midway, above left, font=\scriptsize, align=center]{Paramètres\\ d'entrée};
\draw[retour] (fonc.north) ++(0.9,0) -- (pp.south) node[midway, above right, font=\scriptsize, align=center]{Valeur de\\ retour};
\draw[fleche] (pp.south) ++(1.1,0) -- (proc.north) node[midway, above right, font=\scriptsize, align=center]{Paramètres\\ d'entrée};
\end{tikzpicture}
\legende{Le programme principal délègue le travail à des modules : il envoie des paramètres d'entrée ; la fonction renvoie une valeur, la procédure agit sans rien renvoyer. Chaque module possède ses propres variables locales.}
\end{popfigure}

\courssub{La fonction : un module qui calcule et renvoie une valeur}

\begin{definition}[Fonction]
Une \cle{fonction} est un module qui reçoit des \cle{paramètres} (les entrées), effectue un traitement et \textbf{renvoie obligatoirement une unique valeur résultat} (le retour), grâce au mot-clé \texttt{RETOURNER}.
\end{definition}

L'intuition est simple : une fonction est comme une machine à calculer. Tu lui donnes des ingrédients (les paramètres), elle te rend un plat (le résultat). Exemple : \texttt{CalculerInteret(capital, taux)} reçoit un capital et un taux, et renvoie l'intérêt calculé.

\begin{propriete}[Syntaxe générale d'une fonction en pseudo-code]
\begin{lstlisting}
FONCTION NomFonction(param1 : TYPE, param2 : TYPE) : TYPE_RETOUR
    VARIABLES
        variableLocale : TYPE
    DEBUT
        // traitements, calculs
        RETOURNER resultat
    FIN
\end{lstlisting}
Le type écrit après les deux-points finaux (ici \texttt{TYPE\_RETOUR}) indique la nature de la valeur renvoyée (ENTIER, REEL, CHAINE, BOOLEEN...).
\end{propriete}

\begin{methode}[Comment écrire une fonction, en 5 étapes]
\begin{enumerate}
  \item \textbf{Nommer} la fonction avec un verbe d'action : \texttt{Calculer}, \texttt{Trouver}, \texttt{Convertir}... Le nom doit annoncer clairement ce qu'elle fait.
  \item \textbf{Lister les paramètres entrants} avec leurs types : de quoi la fonction a-t-elle besoin pour travailler ?
  \item \textbf{Déclarer les variables locales} nécessaires aux calculs intermédiaires.
  \item \textbf{Écrire le corps} de la fonction : les calculs, tests ou boucles qui produisent le résultat.
  \item \textbf{Renvoyer le résultat} avec \texttt{RETOURNER}.
\end{enumerate}
\end{methode}

\begin{exemplebox}[Exemple chiffré : les frais d'un transfert d'argent]
Une agence de transfert à Douala applique le tarif suivant : pour un montant inférieur ou égal à \num{5000} FCFA, les frais sont fixes : 150 FCFA ; au-delà, les frais valent 2\,\% du montant. Écrivons la fonction :
\begin{lstlisting}
FONCTION FraisTransfert(montant : REEL) : REEL
    DEBUT
        SI montant <= 5000 ALORS
            RETOURNER 150
        SINON
            RETOURNER montant * 0.02
        FINSI
    FIN
\end{lstlisting}
Vérifions à la main : \texttt{FraisTransfert(3000)} renvoie $150$ (car $3000 \leq 5000$). \texttt{FraisTransfert(20000)} renvoie $20000 \times 0{,}02 = 400$ FCFA. La fonction fait \emph{un seul} travail : calculer les frais. Elle ne les affiche pas, elle les renvoie.
\end{exemplebox}

\begin{attention}[Principe de responsabilité unique]
Une fonction ne doit faire qu'\textbf{une seule chose}, et la faire bien. Erreur classique du débutant : mélanger le \emph{calcul} et l'\emph{affichage} dans la même fonction. Si \texttt{FraisTransfert} affiche le résultat au lieu de le renvoyer, on ne peut plus réutiliser ce calcul ailleurs (par exemple pour additionner les frais de plusieurs transferts). La règle d'or : \textbf{la fonction calcule et renvoie ; c'est le programme principal qui affiche}.
\end{attention}

\courssub{La procédure : un module qui agit sans rien renvoyer}

Parfois, on n'a pas besoin d'un résultat calculé : on veut simplement \emph{faire quelque chose} — afficher un menu à l'écran, saisir les informations d'un membre, imprimer un reçu. C'est le rôle de la procédure.

\begin{definition}[Procédure]
Une \cle{procédure} est un module qui reçoit des paramètres, effectue un traitement (par exemple un affichage, une saisie, ou la modification d'une variable passée en paramètre de sortie) mais \textbf{ne renvoie pas de valeur} par \texttt{RETOURNER}.
\end{definition}

\begin{exemplebox}[Deux procédures typiques]
\begin{lstlisting}
PROCEDURE AfficherMenu()
    DEBUT
        ECRIRE("1. Ajouter un membre")
        ECRIRE("2. Afficher les membres")
        ECRIRE("3. Quitter")
    FIN

PROCEDURE AfficherRecu(nom : CHAINE, total : REEL)
    DEBUT
        ECRIRE("---- RECU ----")
        ECRIRE("Client : ", nom)
        ECRIRE("Total a payer : ", total, " FCFA")
    FIN
\end{lstlisting}
On appelle ensuite ces procédures dans le programme principal simplement en écrivant leur nom : \texttt{AfficherMenu()} ou \texttt{AfficherRecu("Amina", 4500)}. Remarque la différence avec une fonction : l'appel d'une procédure est une \emph{instruction complète}, alors que l'appel d'une fonction produit une \emph{valeur} que l'on range dans une variable : \texttt{f <- FraisTransfert(20000)}.
\end{exemplebox}

\courssub{Les paramètres : entrées et sorties}

Les \cle{paramètres} sont les données échangées entre le programme appelant et le module. On distingue :

\begin{itemize}
  \item les \cle{paramètres d'entrée} : une valeur \emph{fournie} au module pour qu'il travaille (exemple : le \texttt{capital} et le \texttt{taux} de \texttt{CalculerInteret}) ;
  \item les \cle{paramètres de sortie} : une variable que le module \emph{modifie} pour communiquer un résultat (exemple : \texttt{nouveauSolde} mis à jour par une procédure \texttt{Deposer}).
\end{itemize}

\begin{propriete}[Passage par valeur et passage par référence (notion)]
Lors de l'appel d'un module, les paramètres peuvent être transmis de deux façons :
\begin{itemize}
  \item \textbf{Passage par valeur} : le module reçoit une \emph{copie} de la donnée. S'il la modifie, la variable d'origine du programme principal \emph{ne change pas}.
  \item \textbf{Passage par référence} : le module reçoit l'\emph{adresse} de la variable d'origine. Toute modification faite dans le module est \emph{répercutée} dans le programme principal. C'est le mécanisme utilisé pour les paramètres de sortie.
\end{itemize}
Au niveau de la classe de 3ème, retiens simplement : une \textbf{fonction} communique son résultat par \texttt{RETOURNER} ; une \textbf{procédure} peut modifier des variables grâce aux paramètres de sortie.
\end{propriete}

\courssub{La portée des variables : locales et globales}

\begin{definition}[Portée d'une variable]
La \cle{portée} d'une variable est la zone du programme où elle est visible et utilisable.
\begin{itemize}
  \item Une variable \cle{locale} est déclarée \emph{à l'intérieur} d'un module : elle n'existe que pendant l'exécution de ce module et n'est visible nulle part ailleurs.
  \item Une variable \cle{globale} est déclarée \emph{en dehors de tout module} : elle est visible partout, par le programme principal et par tous les modules.
\end{itemize}
\end{definition}

\begin{attention}[Pourquoi éviter les variables globales ?]
Les variables globales semblent pratiques, mais elles sont dangereuses : n'importe quel module peut les modifier, et quand une erreur apparaît, il est très difficile de savoir \emph{qui} a changé la valeur. Règle de bonne conduite : \textbf{on privilégie les variables locales et les paramètres} ; on n'utilise une variable globale qu'en dernier recours, pour une information vraiment partagée par tout le programme.
\end{attention}

\courssub{Les avantages de la modularité}

\begin{aretenir}[Pourquoi découper un programme en modules ?]
\begin{itemize}
  \item \textbf{Lisibilité} : chaque module est court et porte un nom clair ; on comprend le programme en lisant les noms des appels.
  \item \textbf{Réutilisabilité} : une fonction écrite une fois peut être appelée autant de fois qu'on veut, dans le même programme ou dans un autre.
  \item \textbf{Test unitaire} : on peut tester chaque module \emph{indépendamment} des autres, ce qui rend la chasse aux erreurs bien plus facile.
  \item \textbf{Travail d'équipe} : chaque élève du groupe peut écrire un module différent, puis on assemble le tout.
  \item \textbf{Maintenance} : si une règle change (par exemple le tarif des frais), on ne modifie qu'un seul module.
\end{itemize}
\end{aretenir}

\begin{exoresolu}[Gestion d'une caisse de tontine]
Énoncé. Une tontine scolaire applique une cotisation fixe de \num{1000} FCFA par membre, plus une commission de 5\,\% sur le montant épargné. 1) Écris une fonction \texttt{Commission(montant : REEL) : REEL} qui renvoie la commission. 2) Écris une procédure \texttt{AfficherBilan(nom : CHAINE, montant : REEL)} qui affiche le nom du membre, la commission et le total à payer. 3) Écris le programme principal pour un membre « Ngo Bell » ayant épargné \num{20000} FCFA.
\tcblower
Solution détaillée. 1) La fonction calcule et renvoie une seule valeur :
\begin{lstlisting}
FONCTION Commission(montant : REEL) : REEL
    DEBUT
        RETOURNER montant * 0.05
    FIN
\end{lstlisting}
2) La procédure effectue des affichages, elle ne renvoie rien. Elle \emph{appelle} elle-même la fonction \texttt{Commission} : c'est la réutilisabilité en action.
\begin{lstlisting}
PROCEDURE AfficherBilan(nom : CHAINE, montant : REEL)
    VARIABLES
        com, total : REEL   // variables locales
    DEBUT
        com <- Commission(montant)
        total <- 1000 + com
        ECRIRE("Membre : ", nom)
        ECRIRE("Commission : ", com, " FCFA")
        ECRIRE("Total a payer : ", total, " FCFA")
    FIN
\end{lstlisting}
3) Le programme principal devient très court et très lisible :
\begin{lstlisting}
ALGORITHME Tontine
DEBUT
    AfficherBilan("Ngo Bell", 20000)
FIN
\end{lstlisting}
Vérification : commission $= 20000 \times 0{,}05 = 1000$ FCFA ; total $= 1000 + 1000 = 2000$ FCFA. Le programme affiche donc : Membre : Ngo Bell — Commission : 1000 FCFA — Total à payer : 2000 FCFA. Remarque que \texttt{com} et \texttt{total} sont des variables \emph{locales} : elles n'existent que pendant l'exécution de la procédure.
\end{exoresolu}

\begin{experience}[TP guidé : construire et tester tes premiers modules]
\textbf{Objectif} : découper un problème en modules et tester chaque module séparément. \textbf{Matériel} : papier et crayon, ou un logiciel d'algorithmique (ou Python si disponible en salle informatique).
\begin{enumerate}
  \item Écris sur papier la fonction \texttt{FraisTransfert(montant : REEL) : REEL} vue plus haut.
  \item \textbf{Test unitaire} : construis une table d'exécution avec trois montants : 2000, 5000 et 50000. Calcule à la main le résultat attendu pour chacun (150 ; 150 ; 1000).
  \item Écris une procédure \texttt{AfficherFacture(montant : REEL)} qui appelle \texttt{FraisTransfert} puis affiche le montant, les frais et le total.
  \item Écris le programme principal qui appelle \texttt{AfficherFacture(50000)}.
  \item \textbf{Observation} : modifie le tarif (par exemple 3\,\% au lieu de 2\,\%). Que dois-tu changer ? Uniquement le corps de la fonction \texttt{FraisTransfert} : c'est tout l'intérêt de la modularité !
\end{enumerate}
\textbf{Interprétation} : tester chaque module isolément permet de localiser immédiatement l'origine d'une erreur, au lieu de chercher dans tout le programme.
\end{experience}

\begin{savaistu}[Ne réinvente pas la roue !]
Les langages de programmation modernes sont livrés avec des \cle{bibliothèques standard} : d'immenses collections de fonctions déjà écrites, testées et optimisées par des experts. En Python par exemple, la bibliothèque \texttt{math} fournit des fonctions comme \texttt{sqrt} (racine carrée), \texttt{random} permet de tirer des nombres au hasard, et \texttt{datetime} gère les dates et les heures. Un bon programmeur n'écrit que ce qui n'existe pas encore : pour tout le reste, il réutilise les modules existants. C'est la modularité poussée à l'échelle mondiale !
\end{savaistu}

\begin{exercice}[À toi de jouer]
1) Écris une fonction \texttt{Moyenne(n1 : REEL, n2 : REEL, n3 : REEL) : REEL} qui renvoie la moyenne de trois notes. 2) Écris une procédure \texttt{AfficherMention(moy : REEL)} qui affiche « Admis » si la moyenne est supérieure ou égale à 10, « Ajourné » sinon. 3) Écris le programme principal qui calcule puis affiche la mention d'un élève ayant obtenu 12, 8 et 14. 4) Indique, pour chaque variable que tu as utilisée, si elle est locale ou globale.
\end{exercice}

\begin{corrige}[Exercice « À toi de jouer »]
1) La fonction reçoit trois notes et renvoie leur moyenne :
\begin{lstlisting}
FONCTION Moyenne(n1 : REEL, n2 : REEL, n3 : REEL) : REEL
    DEBUT
        RETOURNER (n1 + n2 + n3) / 3
    FIN
\end{lstlisting}
2) La procédure affiche la mention, sans rien renvoyer :
\begin{lstlisting}
PROCEDURE AfficherMention(moy : REEL)
    DEBUT
        SI moy >= 10 ALORS
            ECRIRE("Admis")
        SINON
            ECRIRE("Ajourne")
        FINSI
    FIN
\end{lstlisting}
3) Le programme principal enchaîne les appels :
\begin{lstlisting}
ALGORITHME Bulletin
VARIABLES
    m : REEL   // variable du programme principal
DEBUT
    m <- Moyenne(12, 8, 14)
    ECRIRE("Moyenne : ", m)
    AfficherMention(m)
FIN
\end{lstlisting}
Vérification : $m = (12 + 8 + 14) / 3 = 34 / 3 \approx 11{,}33$, donc le programme affiche « Admis ». 4) La variable \texttt{m} est déclarée dans le programme principal : elle est \emph{locale} au programme principal (elle n'est pas visible à l'intérieur des modules). Les paramètres \texttt{n1}, \texttt{n2}, \texttt{n3} et \texttt{moy} sont des variables \emph{locales} à leurs modules respectifs. Aucune variable globale n'a été utilisée : c'est la bonne pratique.
\end{corrige}

\coursec{Valider et documenter : Tests, Débogage et Qualité}

Dans la section précédente, nous avons appris à \emph{diviser pour régner} : un gros problème se découpe en petites fonctions et procédures, chacune facile à écrire et à comprendre. Mais une question essentielle reste en suspens : \textbf{comment être sûr que notre programme fait bien ce qu'on lui demande ?} Un programme qui s'exécute sans planter n'est pas forcément un programme \emph{correct}. C'est toute la mission de cette dernière section : apprendre à \cle{tester}, à \cle{déboguer}, à \cle{documenter} et à \cle{sauvegarder} notre travail comme le font les vrais développeurs.

\courssub{Pourquoi tester ? Le code sans test est du code cassé}

\begin{definition}[Test d'un programme]
\cle{Tester} un programme, c'est l'exécuter avec des \textbf{données choisies à l'avance} (appelées \cle{jeu de test}) dont on connaît le résultat attendu, puis comparer le \textbf{résultat obtenu} au \textbf{résultat attendu}. Si les deux coïncident, le test \emph{passe} (OK) ; sinon, il \emph{échoue} (KO) et révèle un \cle{bug} (une erreur dans le programme).
\end{definition}

L'intuition est simple : imagine un menuisier de quartier à Douala qui fabrique une porte. Avant de la livrer, il vérifie qu'elle ferme, que la serrure tourne, qu'elle ne frotte pas le sol. S'il livre sans vérifier, le client découvrira les défauts à sa place — et ne le rappellera plus. En informatique, c'est pareil : \textbf{un code qui n'a jamais été testé doit être considéré comme cassé}, car on n'a aucune preuve qu'il fonctionne. Les erreurs ne se voient pas à l'œil nu dans le texte du programme : elles se révèlent à l'exécution.

\begin{definition}[Régression]
Une \cle{régression} est l'apparition d'un bug dans une fonctionnalité qui \textbf{marchait correctement avant une modification} du programme. Par exemple, on corrige le calcul des frais et, sans s'en rendre compte, on casse l'affichage du solde. Rejouer les tests après chaque modification permet de détecter immédiatement les régressions.
\end{definition}

\courssub{Les types de tests : du petit morceau au logiciel complet}

On ne teste pas tout de la même manière. On procède par couches, comme un mécanicien qui vérifie d'abord chaque pièce, puis le moteur assemblé, puis la voiture entière sur la route.

\begin{center}
\begin{tabularx}{\linewidth}{|l|X|X|}
\hline
\rowcolor{popblueL} \textbf{Type de test} & \textbf{Ce qu'on vérifie} & \textbf{Exemple (application de transfert)} \\
\hline
\cle{Test unitaire} & Chaque fonction ou procédure \textbf{isolément} & La fonction \texttt{FraisTransfert(montant)} renvoie-t-elle le bon montant de frais ? \\
\hline
\cle{Test d'intégration} & Le bon fonctionnement des morceaux \textbf{assemblés ensemble} & La fonction \texttt{FraisTransfert} est-elle bien appelée par la procédure \texttt{EffectuerTransfert}, et le solde est-il bien diminué du montant + frais ? \\
\hline
\cle{Test système} & Le logiciel complet répond-il au \textbf{besoin de l'utilisateur} ? & L'élève qui utilise l'application arrive-t-il à envoyer 5\,000 FCFA à un camarade et à voir son reçu ? \\
\hline
\end{tabularx}
\end{center}

\begin{aretenir}[Ordre logique des tests]
On teste d'abord \textbf{unitairement} chaque brique, puis l'\textbf{intégration} des briques entre elles, puis le \textbf{système} complet. C'est l'avantage direct de la modularité vue à la section 4 : une fonction isolée est beaucoup plus facile à tester qu'un programme d'un seul bloc.
\end{aretenir}

\courssub{Construire un bon jeu de test : la stratégie des cas de test}

La qualité d'un test ne dépend pas du \emph{nombre} de cas, mais de leur \emph{choix}. Un bon jeu de test explore trois familles de valeurs :

\begin{itemize}
\item les \cle{valeurs nominales} : les cas normaux, typiques, « au milieu » (un montant ordinaire) ;
\item les \cle{valeurs limites} : les bornes exactes des conditions (0, la valeur maximale, la valeur juste en dessous et juste au-dessus d'un seuil) ;
\item les \cle{valeurs invalides} : ce que le programme ne devrait jamais recevoir (un montant négatif, du texte à la place d'un nombre) pour vérifier qu'il réagit proprement (message d'erreur, refus) au lieu de produire un résultat absurde.
\end{itemize}

\begin{methode}[Construire un jeu de test complet]
Pour tester correctement une fonction ou un programme :
\begin{enumerate}
\item \textbf{Cas standard} : choisir une valeur moyenne, représentative de l'usage normal.
\item \textbf{Cas limites} : tester chaque borne des conditions, incluse et exclue (par exemple pour la condition \texttt{montant >= 5000}, tester 4999, 5000 et 5001).
\item \textbf{Cas d'erreur} : fournir des entrées invalides (négatif, texte, vide) et vérifier la gestion prévue (message d'erreur clair).
\item \textbf{Cas extrêmes} : tester un très grand nombre, une liste vide, une boucle qui s'exécute zéro fois.
\end{enumerate}
\end{methode}

\begin{exemplebox}[Tester un calcul de remise — exemple chiffré]
Une boutique applique la règle : \texttt{SI achat >= 10000 ALORS remise <- 5\% de achat SINON remise <- 0}. Voici le jeu de test minimal et les résultats attendus :
\begin{itemize}
\item \texttt{achat = 15000} (nominal) : remise attendue $= 15000 \times 0,05 = 750$ FCFA ;
\item \texttt{achat = 10000} (limite, seuil inclus) : remise attendue $= 500$ FCFA ;
\item \texttt{achat = 9999} (limite, juste sous le seuil) : remise attendue $= 0$ FCFA ;
\item \texttt{achat = -500} (invalide) : le programme doit signaler une erreur, pas calculer une remise négative.
\end{itemize}
Ces quatre cas suffisent à piéger l'erreur classique consistant à écrire \texttt{>} au lieu de \texttt{>=} : le test avec 10000 échouerait et révélerait le bug.
\end{exemplebox}

\begin{attention}[Le piège du « chemin heureux »]
Tester uniquement le cas où tout se passe bien (le « chemin heureux ») est \textbf{insuffisant}. Il faut aussi tester : les branches \texttt{SINON} des conditions, les boucles qui s'exécutent \textbf{zéro fois} (liste vide), les cas où une \textbf{division par zéro} pourrait se produire, et les entrées invalides. C'est presque toujours là que se cachent les bugs.
\end{attention}

\courssub{La table de test (table de traçabilité)}

Pour organiser les tests et prouver qu'ils ont été faits, on utilise une \cle{table de test} : chaque ligne est un cas de test, chaque colonne décrit une étape du raisonnement. Prenons la fonction \texttt{FraisTransfert} d'une application de transfert d'argent, dont le cahier des charges est :

\begin{itemize}
\item si \texttt{montant <= 0} : erreur, transfert refusé ;
\item si \texttt{0 < montant < 5000} : frais = 100 FCFA ;
\item si \texttt{5000 <= montant <= 10000} : frais = 250 FCFA ;
\item si \texttt{montant > 10000} : frais = 500 FCFA.
\end{itemize}

\begin{popfigure}
\begin{center}
\begin{tikzpicture}[x=1cm,y=1cm]
% En-tête
\fill[popblue] (0,0) rectangle (1.3,0.7);
\fill[popblue] (1.3,0) rectangle (3.6,0.7);
\fill[popblue] (3.6,0) rectangle (6.9,0.7);
\fill[popblue] (6.9,0) rectangle (9.6,0.7);
\fill[popblue] (9.6,0) rectangle (12.3,0.7);
\fill[popblue] (12.3,0) rectangle (14.3,0.7);
\node[white,font=\footnotesize\bfseries] at (0.65,0.35) {Test N°};
\node[white,font=\footnotesize\bfseries] at (2.45,0.35) {Montant};
\node[white,font=\footnotesize\bfseries] at (5.25,0.35) {Condition évaluée};
\node[white,font=\footnotesize\bfseries] at (8.25,0.35) {Frais attendus};
\node[white,font=\footnotesize\bfseries] at (10.95,0.35) {Résultat prog.};
\node[white,font=\footnotesize\bfseries] at (13.3,0.35) {Statut};
% Ligne 1
\fill[popblueL] (0,-0.7) rectangle (14.3,0);
\node[font=\footnotesize] at (0.65,-0.35) {1};
\node[font=\footnotesize] at (2.45,-0.35) {3000 (nominal)};
\node[font=\footnotesize] at (5.25,-0.35) {0 < m < 5000};
\node[font=\footnotesize] at (8.25,-0.35) {100};
\node[font=\footnotesize] at (10.95,-0.35) {100};
\node[font=\footnotesize\bfseries,popgreen] at (13.3,-0.35) {OK};
% Ligne 2
\fill[white] (0,-1.4) rectangle (14.3,-0.7);
\node[font=\footnotesize] at (0.65,-1.05) {2};
\node[font=\footnotesize] at (2.45,-1.05) {5000 (limite)};
\node[font=\footnotesize] at (5.25,-1.05) {5000 <= m <= 10000};
\node[font=\footnotesize] at (8.25,-1.05) {250};
\node[font=\footnotesize] at (10.95,-1.05) {250};
\node[font=\footnotesize\bfseries,popgreen] at (13.3,-1.05) {OK};
% Ligne 3
\fill[popblueL] (0,-2.1) rectangle (14.3,-1.4);
\node[font=\footnotesize] at (0.65,-1.75) {3};
\node[font=\footnotesize] at (2.45,-1.75) {10000 (limite)};
\node[font=\footnotesize] at (5.25,-1.75) {5000 <= m <= 10000};
\node[font=\footnotesize] at (8.25,-1.75) {250};
\node[font=\footnotesize] at (10.95,-1.75) {500};
\node[font=\footnotesize\bfseries,poporange] at (13.3,-1.75) {KO};
% Ligne 4
\fill[white] (0,-2.8) rectangle (14.3,-2.1);
\node[font=\footnotesize] at (0.65,-2.45) {4};
\node[font=\footnotesize] at (2.45,-2.45) {15000 (nominal)};
\node[font=\footnotesize] at (5.25,-2.45) {m > 10000};
\node[font=\footnotesize] at (8.25,-2.45) {500};
\node[font=\footnotesize] at (10.95,-2.45) {500};
\node[font=\footnotesize\bfseries,popgreen] at (13.3,-2.45) {OK};
% Ligne 5
\fill[popblueL] (0,-3.5) rectangle (14.3,-2.8);
\node[font=\footnotesize] at (0.65,-3.15) {5};
\node[font=\footnotesize] at (2.45,-3.15) {-100 (invalide)};
\node[font=\footnotesize] at (5.25,-3.15) {m <= 0};
\node[font=\footnotesize] at (8.25,-3.15) {Erreur};
\node[font=\footnotesize] at (10.95,-3.15) {Erreur};
\node[font=\footnotesize\bfseries,popgreen] at (13.3,-3.15) {OK};
% Cadre
\draw[popdark,thick] (0,0.7) rectangle (14.3,-3.5);
\draw[popline] (1.3,0.7) -- (1.3,-3.5);
\draw[popline] (3.6,0.7) -- (3.6,-3.5);
\draw[popline] (6.9,0.7) -- (6.9,-3.5);
\draw[popline] (9.6,0.7) -- (9.6,-3.5);
\draw[popline] (12.3,0.7) -- (12.3,-3.5);
\foreach \y in {0,-0.7,-1.4,-2.1,-2.8} {\draw[popline] (0,\y) -- (14.3,\y);};
\end{tikzpicture}
\end{center}
\legende{Table de test de la fonction \texttt{FraisTransfert}. Le test n°3 échoue (KO) : le programme renvoie 500 au lieu de 250 pour un montant de 10000, ce qui révèle un bug dans la borne de la deuxième condition (probablement \texttt{< 10000} écrit à la place de \texttt{<= 10000}).}
\end{popfigure}

Grâce à cette table, on voit immédiatement \emph{quel} cas échoue, \emph{quelle} condition est en cause, et on peut corriger précisément. C'est aussi un document à joindre à ton travail pour justifier ta démarche — une compétence exigée au BEPC comme en entreprise.

\courssub{Le débogage : traquer le bug}

Quand un test échoue, commence la phase de \cle{débogage} (en anglais \emph{debugging}) : retrouver et corriger l'erreur. Deux techniques de base, à la portée d'un élève de 3ème :

\begin{itemize}
\item la \cle{trace d'exécution} : on suit à la main, ligne par ligne, la valeur de chaque variable (comme la table d'exécution vue en algorithmique) ;
\item les \cle{affichages intermédiaires} : on insère temporairement des instructions \texttt{Afficher} pour suivre l'évolution des variables pendant l'exécution, puis on les retire une fois le bug corrigé.
\end{itemize}

\begin{exoresolu}[Déboguer la fonction FraisTransfert]
\textbf{Énoncé.} Le test n°3 de la table ci-dessus a échoué : pour \texttt{montant = 10000}, le programme affiche 500 au lieu de 250. Voici le pseudo-code fautif. Trouve le bug en traçant l'exécution, puis corrige-le.
\begin{lstlisting}
Fonction FraisTransfert(montant : reel) : reel
Debut
    Si montant <= 0 Alors
        Afficher("Erreur : montant invalide")
        Retourner -1
    FinSi
    Si montant < 5000 Alors
        Retourner 100
    SinonSi montant < 10000 Alors   // <- zone suspecte
        Retourner 250
    Sinon
        Retourner 500
    FinSi
Fin
\end{lstlisting}
\tcblower
\textbf{Solution.} Traçons l'exécution avec \texttt{montant = 10000} :
\begin{enumerate}
\item \texttt{10000 <= 0} ? Non.
\item \texttt{10000 < 5000} ? Non. On passe au \texttt{SinonSi}.
\item \texttt{10000 < 10000} ? \textbf{Non} (10000 n'est pas strictement inférieur à 10000). On tombe donc dans le \texttt{Sinon} : la fonction retourne 500. C'est bien le résultat fautif observé !
\end{enumerate}
Le bug est identifié : la deuxième condition exclut la borne 10000, alors que le cahier des charges exige \texttt{5000 <= montant <= 10000}. \textbf{Correction} : remplacer \texttt{SinonSi montant < 10000 Alors} par \texttt{SinonSi montant <= 10000 Alors}. On rejoue ensuite toute la table de test : les 5 cas passent au statut OK. \textbf{Leçon} : le débogage compare toujours le comportement \emph{réel} (tracé pas à pas) au comportement \emph{attendu} (cahier des charges), et les bornes des conditions (\texttt{<} contre \texttt{<=}) sont les premières suspectes.
\end{exoresolu}

\courssub{Documenter : écrire pour les humains, pas seulement pour la machine}

Un programme est lu bien plus souvent qu'il n'est écrit — par toi-même dans trois mois, par un camarade, par un correcteur. La \cle{documentation} rend le code compréhensible.

\begin{propriete}[Règles de documentation d'un programme]
\begin{itemize}
\item \textbf{Commentaire d'en-tête} : en début de fichier ou de fonction, indiquer l'\textbf{auteur}, la \textbf{date}, le \textbf{rôle} du programme, ses entrées et ses sorties.
\item \textbf{Commentaires en ligne} : ils expliquent le \textbf{pourquoi} d'une instruction délicate, pas le \emph{comment} (le code dit déjà ce qu'il fait). Inutile d'écrire \texttt{// on ajoute 1 à i} !
\item \textbf{Nommage clair} : des noms de variables explicites (\texttt{montantTransfert} plutôt que \texttt{x}), avec une convention constante : \cle{CamelCase} (\texttt{FraisTransfert}, chaque mot commence par une majuscule) ou \cle{snake\_case} (\texttt{frais\_transfert}, mots séparés par des tirets bas).
\end{itemize}
\end{propriete}

\begin{exemplebox}[Une fonction bien documentée]
\begin{lstlisting}
// ---------------------------------------------------------
// Fonction : FraisTransfert
// Auteur   : A. Mbarga - Date : 12/03/2026
// Role     : calcule les frais d'un transfert d'argent
// Entree   : montant (reel), montant du transfert en FCFA
// Sortie   : frais en FCFA, ou -1 si montant invalide
// ---------------------------------------------------------
Fonction FraisTransfert(montant : reel) : reel
Debut
    Si montant <= 0 Alors
        Retourner -1   // -1 signale une erreur a l'appelant
    FinSi
    Si montant < 5000 Alors
        Retourner 100
    SinonSi montant <= 10000 Alors
        Retourner 250  // borne haute incluse : cahier des charges
    Sinon
        Retourner 500
    FinSi
Fin
\end{lstlisting}
\end{exemplebox}

\courssub{Le versionning : garder l'historique avec Git (initiation)}

Dernier réflexe de professionnel : ne jamais perdre son travail et pouvoir revenir en arrière. Plutôt que d'accumuler des fichiers \texttt{prog\_final.txt}, \texttt{prog\_final2.txt}, \texttt{prog\_vraiment\_final.txt}, on utilise un \cle{gestionnaire de versions} comme \cle{Git}. Il enregistre l'\textbf{historique} des modifications et permet de \textbf{collaborer} à plusieurs sur le même projet.

\begin{methode}[Les trois commandes Git de base]
\begin{enumerate}
\item \texttt{git init} : crée le dépôt (le « dossier surveillé ») du projet, une seule fois au début ;
\item \texttt{git add nom\_fichier} : sélectionne les fichiers dont on veut enregistrer l'état actuel ;
\item \texttt{git commit -m "message"} : enregistre définitivement cet état dans l'historique, avec un message qui explique la modification (par exemple \texttt{"correction des bornes de FraisTransfert"}).
\end{enumerate}
Ainsi, chaque \texttt{commit} est une « photo » du projet : si une modification casse tout (régression !), on peut revenir à une photo précédente.
\end{methode}

\begin{experience}[TP guidé : tester, déboguer, documenter une fonction]
\textbf{Objectif :} appliquer toute la démarche qualité sur la fonction \texttt{FraisTransfert}. \textbf{Matériel :} un ordinateur avec un éditeur de texte (ou papier-crayon pour le pseudo-code).
\begin{enumerate}
\item \textbf{Écriture} : recopie la fonction documentée de l'exemple précédent, avec son en-tête complet (ton nom, la date du jour).
\item \textbf{Préparation} : reproduis la table de test à 5 lignes (3000, 5000, 10000, 15000, -100) et remplis la colonne « Frais attendus » \emph{avant} d'exécuter quoi que ce soit.
\item \textbf{Exécution} : fais tourner le programme (ou trace-le à la main) pour chaque cas ; note le « Résultat programme » et le statut OK/KO.
\item \textbf{Débogage} : introduis volontairement l'erreur \texttt{< 10000} à la place de \texttt{<= 10000}. Quel test échoue ? Retrouve le bug en traçant l'exécution pas à pas, puis corrige.
\item \textbf{Sauvegarde} : si Git est disponible, fais \texttt{git init}, \texttt{git add}, puis un \texttt{git commit} avant la correction et un après, avec des messages clairs.
\end{enumerate}
\textbf{Observations attendues :} le test avec le montant 10000 passe au statut KO avec la version fautive, puis redevient OK après correction. \textbf{Interprétation :} les valeurs limites sont les plus efficaces pour débusquer les erreurs de bornes.
\end{experience}

\begin{savaistu}[Les tests, une compétence professionnelle clé]
Dans les grandes entreprises du numérique, mais aussi dans les banques qui opèrent au Cameroun, on exige souvent qu'une grande partie du code (par exemple plus de 80\,\% des lignes) soit exécutée par au moins un test automatique avant toute mise en service du logiciel : c'est ce qu'on appelle la \cle{couverture de tests}. Un logiciel bancaire mal testé pourrait envoyer de l'argent au mauvais client ! Savoir écrire un jeu de test, c'est déjà posséder une compétence recherchée sur le marché du travail.
\end{savaistu}

\begin{aretenir}[L'essentiel de la section]
\begin{itemize}
\item Un code non testé est un code suspect : on compare toujours \textbf{résultat obtenu} et \textbf{résultat attendu}.
\item On teste par couches : \textbf{unitaire} (chaque fonction), \textbf{intégration} (l'assemblage), \textbf{système} (le besoin utilisateur).
\item Un bon jeu de test mélange valeurs \textbf{nominales}, \textbf{limites} (bornes incluses et exclues) et \textbf{invalides}, sans oublier les cas extrêmes.
\item La \textbf{table de test} trace chaque cas : entrée, condition, sortie attendue, résultat, statut OK/KO.
\item Le \textbf{débogage} utilise la trace d'exécution ou des affichages intermédiaires pour comparer l'exécution réelle au cahier des charges.
\item On \textbf{documente} (en-tête, commentaires utiles, nommage clair) et on \textbf{versionne} avec Git (\texttt{init}, \texttt{add}, \texttt{commit}) pour sécuriser son travail et collaborer.
\end{itemize}
\end{aretenir}

\begin{exercice}[À toi de jouer]
Une fonction \texttt{Mention(moyenne)} retourne « Ajourné » si $\text{moyenne} < 10$, « Passable » si $10 \le \text{moyenne} < 12$, « Bien » si $12 \le \text{moyenne} < 14$, et « Très Bien » si $\text{moyenne} \ge 14$. Construis une table de test complète (au moins 6 cas : nominals, limites 10, 12, 14, et un cas invalide comme $-3$ ou $25$), avec pour chaque ligne la condition évaluée, la sortie attendue et une colonne Statut à remplir.
\end{exercice}

\begin{corrige}[Exercice : table de test de la fonction Mention]
Cas proposés : moyenne $= 11$ (Passable, nominal) ; $10$ (Passable, limite incluse) ; $9{,}99$ (Ajourné, limite exclue) ; $12$ (Bien, limite incluse) ; $14$ (Très Bien, limite incluse) ; $-3$ (entrée invalide : le programme devrait signaler une erreur, car une moyenne est comprise entre 0 et 20) ; éventuellement $25$ (invalide également). Chaque ligne doit indiquer la condition évaluée (par exemple $10 \le m < 12$), la mention attendue, le résultat réel du programme et le statut OK/KO. Les bornes 10, 12 et 14 sont les cas les plus importants : elles détectent immédiatement toute confusion entre \texttt{<} et \texttt{<=}.
\end{corrige}

\newpage

\coursec{Activité d'intégration}

\coursec{Utilisation des concepts de base du développement logiciel}

\courssub{Objectifs de l'activité}

À l'issue de cette activité d'intégration, tu seras capable de mobiliser l'ensemble des notions de base du développement logiciel étudiées dans cette leçon pour résoudre un problème concret de la vie courante. Plus précisément, tu devras :

\begin{itemize}
\item \cle{Analyser} un énoncé de la vie réelle pour en extraire les \cle{entrées}, les \cle{sorties} et les \cle{règles de calcul} (étape d'analyse du cycle de vie d'un logiciel) ;
\item \cle{Concevoir} un algorithme en pseudo-code utilisant des \cle{variables} typées, des \cle{structures conditionnelles} et une \cle{fonction} avec paramètres et valeur de retour ;
\item \cle{Modulariser} la solution en séparant le calcul (fonction) de l'affichage (procédure), selon le principe \og diviser pour régner \fg ;
\item \cle{Valider} le programme par un \cle{jeu de tests} construit méthodiquement (cas normaux, cas limites, cas d'erreur) ;
\item \cle{Rédiger et justifier} ta démarche et tes choix, comme le ferait un vrai développeur qui documente son travail.
\end{itemize}

\courssub{Situation}

\textbf{Simulateur de Tarif Moto-Taxi Yaoundé.}\par
\medskip

À Yaoundé, les motos-taxis (appelées \og benskineurs \fg) sont un moyen de transport très utilisé, du quartier Nkolbisson jusqu'à Mokolo ou d'Ekounou au centre-ville. Mais les prix sont souvent discutés entre le client et le conducteur, ce qui crée des conflits quotidiens. Pour régler ce problème, la communauté urbaine a publié une \cle{grille tarifaire officielle} et souhaite qu'un petit logiciel, installé sur un simple téléphone ou affiché dans les agences, calcule automatiquement le prix d'une course.

La grille officielle est la suivante :

\begin{itemize}
\item \textbf{Prise en charge :} \num{300} FCFA, qui incluent les \num{2} premiers kilomètres ;
\item \textbf{Au-delà de \num{2} km :} chaque kilomètre supplémentaire coûte \num{100} FCFA ;
\item \textbf{Majoration de nuit :} entre \num{22} heures (inclus) et \num{5} heures (exclu), le prix total est majoré de \num{50} \% ;
\item \textbf{Distance invalide :} si la distance saisie est négative, le logiciel doit signaler une erreur au lieu de calculer un prix absurde.
\end{itemize}

L'heure est saisie comme un \cle{nombre entier} entre \num{0} et \num{23} (par exemple \num{14} pour 14 heures, \num{23} pour 23 heures). La distance est un \cle{nombre réel} exprimé en kilomètres (par exemple \num{3,5} km).

Ton établissement a été choisi pour proposer un prototype. Tu es le développeur : à toi d'analyser le problème, de concevoir l'algorithme, de le modulariser, puis de le valider par des tests avant de le présenter à la classe.

\courssub{Consignes}

\begin{enumerate}
\item \textbf{Analyse.} Lis attentivement la situation. Identifie et rédige dans un tableau : (a) les \cle{données d'entrée} du problème avec leur type (ENTIER ou REEL) ; (b) la \cle{donnée de sortie} ; (c) les \cle{constantes} du problème (valeurs fixées par la grille tarifaire) ; (d) les \cle{règles de calcul} exprimées en français.
\item \textbf{Conception de la fonction.} Écris en pseudo-code la fonction \texttt{CalculerTarif(distance : REEL, heure : ENTIER) : REEL} qui retourne le prix de la course. Elle devra : traiter d'abord le cas d'erreur (distance négative), puis calculer le prix de base selon que la distance dépasse ou non \num{2} km, puis appliquer la majoration de nuit si nécessaire. Attention : la nuit correspond à \texttt{heure >= 22} \textbf{ou} \texttt{heure < 5}.
\item \textbf{Modularisation.} Écris la procédure \texttt{AfficherFacture(prix : REEL)} qui affiche proprement la facture du client (par exemple : \og Prix de votre course : 800 FCFA \fg). Explique en deux phrases pourquoi il est préférable de séparer le calcul de l'affichage dans deux modules distincts.
\item \textbf{Validation.} Avant même d'exécuter quoi que ce soit sur machine, remplis la \cle{table de tests} suivante en calculant à la main le résultat attendu pour chaque cas, puis en vérifiant que ton algorithme produit bien ce résultat (trace d'exécution) :

\begin{center}
\begin{tabularx}{\linewidth}{|c|l|c|c|X|}
\hline
\rowcolor{popblueL}
\textbf{N°} & \textbf{Cas testé} & \textbf{distance} & \textbf{heure} & \textbf{Résultat attendu} \\
\hline
1 & Course courte de jour & \num{1} & \num{14} & \ldots \\
\hline
2 & Course moyenne de jour & \num{5} & \num{10} & \ldots \\
\hline
3 & Course courte de nuit & \num{2} & \num{23} & \ldots \\
\hline
4 & Longue course de nuit & \num{10} & \num{2} & \ldots \\
\hline
5 & Distance invalide & \num{-3} & \num{12} & \ldots \\
\hline
\end{tabularx}
\end{center}

\item \textbf{Justification.} Rédige un court paragraphe (5 à 8 lignes) qui justifie : (a) pourquoi la borne des \num{2} km est traitée par un test \texttt{distance <= 2} ; (b) pourquoi la nuit ne peut pas s'écrire avec un seul test du type \texttt{heure >= 22 ET heure < 5} ; (c) pourquoi il est indispensable de tester le cas de la distance négative.
\end{enumerate}

Travail individuel ou en binôme, sur papier ou au tableau. La correction sera faite collectivement au vidéoprojecteur.

\courssub{Ressources à mobiliser}

\begin{aretenir}[Boîte à outils de la leçon]
\begin{itemize}
\item \textbf{Cycle de vie d'un logiciel :} analyse $\rightarrow$ conception $\rightarrow$ codage $\rightarrow$ tests $\rightarrow$ mise en service. Ici, tu pratiques toutes les étapes sauf le codage en vrai langage.
\item \textbf{Variables et types :} une variable possède un \cle{nom}, un \cle{type} (ENTIER, REEL, CHAINE, BOOLEEN) et une \cle{valeur}. Les constantes de la grille (\num{300}, \num{100}, \num{2}, \num{22}, \num{5}) ne changent jamais.
\item \textbf{Structure conditionnelle :}
\begin{lstlisting}
Si condition Alors
    instructions
Sinon
    instructions
FinSi
\end{lstlisting}
\item \textbf{Fonction :} module qui reçoit des \cle{paramètres} et \cle{retourne} un résultat : \texttt{Fonction Nom(param : TYPE) : TYPE}.
\item \textbf{Procédure :} module qui effectue une action (ici un affichage) sans retourner de valeur.
\item \textbf{Test de logiciel :} on choisit des cas \cle{normaux}, des cas \cle{limites} (exactement \num{2} km, exactement \num{22} h) et des cas \cle{d'erreur} (distance négative).
\item \textbf{Formule de la majoration :} prix de nuit $=$ prix de jour $\times (1 + \dfrac{50}{100}) =$ prix de jour $\times 1{,}5$.
\end{itemize}
\end{aretenir}

\courssub{Démarche et corrigé commenté}

\begin{exoresolu}[Simulateur de Tarif Moto-Taxi Yaoundé]
\textbf{Énoncé.} Concevoir et valider un simulateur de tarif de moto-taxi : \num{300} FCFA les \num{2} premiers km, puis \num{100} FCFA par km supplémentaire, avec majoration de \num{50} \% la nuit (22 h à 5 h) et rejet des distances négatives.

\tcblower

\textbf{Étape 1 : Analyse du problème.}\par
On commence par le dictionnaire des données, comme dans toute analyse sérieuse :

\begin{center}
\begin{tabularx}{\linewidth}{|l|l|X|}
\hline
\rowcolor{popblueL}
\textbf{Élément} & \textbf{Rôle} & \textbf{Type / Valeur} \\
\hline
\texttt{distance} & Entrée & REEL, en km \\
\hline
\texttt{heure} & Entrée & ENTIER, entre 0 et 23 \\
\hline
\texttt{prix} & Sortie & REEL, en FCFA \\
\hline
Prise en charge & Constante & \num{300} FCFA (inclut \num{2} km) \\
\hline
Tarif kilométrique & Constante & \num{100} FCFA/km au-delà de \num{2} km \\
\hline
Majoration nuit & Constante & $+50$ \% si \texttt{heure >= 22} ou \texttt{heure < 5} \\
\hline
\end{tabularx}
\end{center}

Règles de calcul en français : si la distance est négative, c'est une erreur ; sinon, le prix de base vaut \num{300} FCFA si la distance est inférieure ou égale à \num{2} km, et \num{300} $+$ \num{100} $\times$ (distance $-$ \num{2}) au-delà ; enfin, si la course a lieu la nuit, on multiplie le total par \num{1,5}.

\medskip
\textbf{Étape 2 : Conception de la fonction.}\par
\begin{lstlisting}
Fonction CalculerTarif(distance : REEL, heure : ENTIER) : REEL
Variable
    prix : REEL
Debut
    // Cas d'erreur : distance impossible
    Si distance < 0 Alors
        Retourner -1   // -1 est le code d'erreur choisi
    FinSi

    // Calcul du prix de base
    Si distance <= 2 Alors
        prix <- 300
    Sinon
        prix <- 300 + 100 * (distance - 2)
    FinSi

    // Majoration de nuit : de 22h inclus a 5h exclu
    Si (heure >= 22) Ou (heure < 5) Alors
        prix <- prix * 1.5
    FinSi

    Retourner prix
Fin
\end{lstlisting}

\textbf{Étape 3 : Modularisation avec la procédure d'affichage.}\par
\begin{lstlisting}
Procedure AfficherFacture(prix : REEL)
Debut
    Si prix = -1 Alors
        Ecrire("Erreur : distance invalide, course refusee.")
    Sinon
        Ecrire("Prix de votre course : ", prix, " FCFA")
    FinSi
Fin
\end{lstlisting}
Séparer le calcul (\texttt{CalculerTarif}) de l'affichage (\texttt{AfficherFacture}) permet de modifier la présentation de la facture sans toucher aux règles de tarification, et inversement : c'est le principe \og diviser pour régner \fg, qui rend le logiciel plus facile à tester et à faire évoluer.

\medskip
\textbf{Étape 4 : Validation par la table de tests.}\par
\begin{center}
\begin{tabularx}{\linewidth}{|c|l|c|c|X|}
\hline
\rowcolor{popblueL}
\textbf{N°} & \textbf{Cas testé} & \textbf{distance} & \textbf{heure} & \textbf{Résultat attendu} \\
\hline
1 & Course courte de jour & \num{1} & \num{14} & \num{300} FCFA \\
\hline
2 & Course moyenne de jour & \num{5} & \num{10} & \num{600} FCFA \\
\hline
3 & Course courte de nuit & \num{2} & \num{23} & \num{450} FCFA \\
\hline
4 & Longue course de nuit & \num{10} & \num{2} & \num{1650} FCFA \\
\hline
5 & Distance invalide & \num{-3} & \num{12} & erreur ($-1$) \\
\hline
\end{tabularx}
\end{center}

Vérification par la trace d'exécution :
\begin{align*}
\text{Cas 2 : } & 300 + 100 \times (5 - 2) = 300 + 300 = 600 \text{ FCFA (jour, pas de majoration)} \\
\text{Cas 3 : } & 300 \times 1{,}5 = 450 \text{ FCFA (2 km pile, donc tarif de base, mais 23 h = nuit)} \\
\text{Cas 4 : } & (300 + 100 \times (10 - 2)) \times 1{,}5 = 1100 \times 1{,}5 = 1650 \text{ FCFA}
\end{align*}
Les cinq résultats calculés à la main correspondent exactement à ce que produit l'algorithme : le jeu de tests est validé.

\medskip
\textbf{Étape 5 : Justification des choix.}\par
(a) La borne des \num{2} km est traitée par \texttt{distance <= 2} car la grille dit que la prise en charge \emph{inclut} les 2 premiers kilomètres : à exactement \num{2} km, on ne paie encore que \num{300} FCFA, et le supplément ne commence qu'au-delà. (b) La nuit ne peut pas s'écrire \texttt{heure >= 22 ET heure < 5}, car aucun entier n'est à la fois supérieur ou égal à \num{22} et strictement inférieur à \num{5} : cette condition serait toujours fausse et la majoration ne s'appliquerait jamais. La plage horaire \og traverse minuit \fg, il faut donc le \texttt{Ou}. (c) Le test de la distance négative est indispensable : sans lui, une erreur de saisie (par exemple $-3$ km) produirait un prix inférieur à \num{300} FCFA, voire négatif, ce qui serait absurde et ferait perdre de l'argent au conducteur. Tester les cas d'erreur fait partie du travail professionnel de validation.
\end{exoresolu}

\begin{attention}[Pièges fréquents rencontrés dans cette activité]
\begin{itemize}
\item Écrire \texttt{distance < 2} au lieu de \texttt{distance <= 2} : le client de \num{2} km pile se verrait alors facturer \num{400} FCFA au lieu de \num{300} FCFA. Les \cle{cas limites} sont les premiers à tester !
\item Utiliser \texttt{Et} au lieu de \texttt{Ou} pour la plage de nuit (voir justification ci-dessus).
\item Appliquer la majoration de \num{50} \% en ajoutant seulement \num{50} FCFA : une majoration de \num{50} \% est une \emph{multiplication} par \num{1,5}, pas une addition fixe.
\item Oublier de traiter la distance négative, ou la traiter \emph{après} le calcul : le test d'erreur doit venir \emph{en premier}.
\end{itemize}
\end{attention}

\courssub{Bilan de l'activité}

Cette activité t'a fait parcourir, à petite échelle mais de manière complète, le \cle{cycle de vie d'un logiciel} : analyser un besoin réel (les conflits de prix entre clients et benskineurs), concevoir une solution algorithmique, la structurer en modules (une fonction et une procédure), puis la valider par un jeu de tests construit méthodiquement.

Elle a mis en évidence plusieurs idées essentielles du développement logiciel :

\begin{itemize}
\item Un bon programme commence par une \cle{analyse sur papier} : identifier entrées, sorties et constantes avant d'écrire la moindre ligne évite la plupart des erreurs ;
\item Les \cle{structures conditionnelles} traduisent directement les règles de la vie réelle (bornes tarifaires, plages horaires), mais exigent une attention particulière aux \cle{cas limites} (\num{2} km exactement, \num{22} h exactement) ;
\item La \cle{modularité} (fonction de calcul $+$ procédure d'affichage) rend le logiciel plus clair, plus facile à tester et plus facile à modifier si la grille tarifaire change l'année prochaine ;
\item Un logiciel n'est jamais considéré comme terminé tant qu'il n'a pas été \cle{validé par des tests}, y compris des tests de cas absurdes (distance négative) : c'est la démarche qualité des développeurs professionnels.
\end{itemize}

Tu peux désormais transposer cette démarche à d'autres situations camerounaises : calcul du prix d'un trajet en car interurbain, facturation d'une borne de recharge téléphonique, ou calcul d'une moyenne de notes avec mention. La méthode reste la même : analyser, concevoir, modulariser, tester, justifier.

\newpage

\coursec{Méthodes et exercices résolus}

\courssub{Rappels essentiels : du besoin au programme}

\begin{aretenir}[Le cycle de vie d'un logiciel]
Un logiciel ne naît pas d'un coup : il suit un \cle{cycle de vie} en étapes :
\begin{enumerate}
\item \textbf{Analyse du besoin} : comprendre ce que le logiciel doit faire ;
\item \textbf{Conception} : écrire l'\cle{algorithme}, c'est-à-dire la suite ordonnée d'étapes qui résout le problème ;
\item \textbf{Codage} : traduire l'algorithme dans un \cle{langage de programmation} (JavaScript, C, Python...) pour obtenir le \cle{programme} ;
\item \textbf{Tests et débogage} : vérifier que le programme donne les bons résultats et corriger les erreurs (\cle{bugs}) ;
\item \textbf{Mise à jour et maintenance} : faire évoluer le logiciel selon les nouveaux besoins.
\end{enumerate}
L'ordinateur ne comprend que le \cle{langage machine} (suites de 0 et de 1) : le programme écrit par le programmeur (le \cle{code source}) doit donc être \emph{traduit} par un logiciel spécial appelé \cle{compilateur} ou \cle{interpréteur} avant de pouvoir s'exécuter.
\end{aretenir}

\courssub{Méthode 1 : Appliquer les notions de développement logiciel à une situation concrète}

\begin{methode}[Résoudre un problème de la vie courante par le développement logiciel]
Pour passer d'une situation concrète (boutique, cantine scolaire, club sportif) à un logiciel, suis toujours la même démarche :
\begin{enumerate}
\item \textbf{Comprendre le besoin} : formuler en une phrase ce que le logiciel doit faire (exemple : « calculer la moyenne d'un élève »).
\item \textbf{Identifier les données} : distinguer les \cle{données d'entrée} (ce que l'utilisateur fournit), les \cle{traitements} (calculs, comparaisons) et les \cle{résultats} (ce que le logiciel affiche).
\item \textbf{Choisir les variables} : donner à chaque donnée un nom clair et un \cle{type} adapté (entier, réel, chaîne de caractères).
\item \textbf{Écrire l'algorithme} : décrire les étapes en pseudo-code, dans l'ordre : Déclaration, Saisie, Traitement, Affichage.
\item \textbf{Traduire en programme} : choisir un langage (par exemple JavaScript) et transcrire l'algorithme en respectant la syntaxe.
\item \textbf{Tester} : exécuter le programme avec des valeurs connues dont on connaît le résultat à l'avance, puis corriger les erreurs (\cle{débogage}).
\end{enumerate}
\textbf{Réflexe} : un bon programme se lit dans l'ordre \texttt{Saisie} $\rightarrow$ \texttt{Traitement} $\rightarrow$ \texttt{Affichage}. Si l'une de ces trois étapes manque, la solution est incomplète.
\end{methode}

\begin{exoresolu}[La boutique de Mme Etoundi]
Mme Etoundi tient une boutique à Yaoundé. Elle vend des cahiers à \num{500} FCFA l'unité. Elle souhaite un petit programme qui demande le nombre de cahiers achetés par un client et affiche le montant total à payer.
\begin{enumerate}
\item Identifier les données d'entrée, le traitement et le résultat.
\item Écrire l'algorithme correspondant.
\item Traduire cet algorithme en JavaScript.
\end{enumerate}
\tcblower
\textbf{1. Analyse du problème.}
\begin{itemize}
\item \textbf{Donnée d'entrée} : le nombre de cahiers, que l'on nomme \texttt{n} (type : entier).
\item \textbf{Traitement} : calcul du montant total : $\text{total} = n \times 500$.
\item \textbf{Résultat} : affichage du montant total en FCFA.
\end{itemize}
On a besoin de deux variables : \texttt{n} (entier) et \texttt{total} (entier), ainsi que de la constante \num{500} (prix unitaire).

\textbf{2. Algorithme.}
\begin{lstlisting}
Algorithme Boutique
Variables n, total : Entier
Début
    Écrire("Entrez le nombre de cahiers : ")
    Lire(n)
    total <- n * 500
    Écrire("Le montant à payer est : ", total, " FCFA")
Fin
\end{lstlisting}
L'ordre est respecté : saisie de \texttt{n}, puis traitement (calcul de \texttt{total}), puis affichage du résultat.

\textbf{3. Traduction en JavaScript.}
\begin{lstlisting}[language=JavaScript]
// Programme de calcul du montant à payer
var n = parseInt(prompt("Entrez le nombre de cahiers : "));
var total = n * 500; // prix unitaire : 500 FCFA
alert("Le montant à payer est : " + total + " FCFA");
\end{lstlisting}
\textbf{Vérification} : pour $n = 4$ cahiers, le programme calcule $4 \times 500 = \num{2000}$ et affiche « Le montant à payer est : 2000 FCFA », ce qui est correct.

\textbf{Conclusion} : en suivant la démarche Saisie $\rightarrow$ Traitement $\rightarrow$ Affichage, on obtient un programme simple, correct et testé, qui répond exactement au besoin de Mme Etoundi.
\end{exoresolu}

\courssub{Méthode 2 : Choisir correctement les variables et leurs types}

\begin{methode}[Déclarer les bonnes variables avec le bon type]
\begin{enumerate}
\item \textbf{Repérer chaque information} manipulée par le problème (une information = une variable).
\item \textbf{Choisir un nom explicite} : \texttt{moyenne}, \texttt{age}, \texttt{nomEleve} plutôt que \texttt{x}, \texttt{y}, \texttt{z}.
\item \textbf{Choisir le type} selon la nature de la valeur :
\begin{itemize}
\item \cle{Entier} : un nombre sans virgule (âge, quantité, nombre d'élèves) ;
\item \cle{Réel} : un nombre à virgule (moyenne, prix, taille) ;
\item \cle{Chaîne de caractères} : un texte (nom, ville, mention) ;
\item \cle{Booléen} : vrai ou faux (admis ou non).
\end{itemize}
\item \textbf{Initialiser} la variable (lui donner une première valeur) avant de l'utiliser dans un calcul.
\end{enumerate}
\textbf{À retenir} : une moyenne sur 20 est un \emph{réel} (elle peut valoir 12,75), jamais un entier ; un nom d'élève est une \emph{chaîne}, jamais un nombre.
\end{methode}

\begin{exoresolu}[Les variables du bulletin de notes]
Le proviseur d'un collège de Douala veut un programme qui gère les informations suivantes pour chaque élève : son nom complet, son âge, sa moyenne générale sur 20 et le fait qu'il soit admis ou non en classe supérieure.
\begin{enumerate}
\item Pour chaque information, proposer un nom de variable et son type, en justifiant.
\item Écrire la partie « Déclaration » de l'algorithme.
\end{enumerate}
\tcblower
\textbf{1. Choix des variables.}
\begin{center}
\begin{tabularx}{\linewidth}{|l|l|l|X|}
\hline
\rowcolor{popblueL} \textbf{Information} & \textbf{Nom de variable} & \textbf{Type} & \textbf{Justification} \\
\hline
Nom complet & \texttt{nomEleve} & Chaîne & c'est un texte (lettres) \\
\hline
Âge & \texttt{age} & Entier & un âge est un nombre sans virgule \\
\hline
Moyenne sur 20 & \texttt{moyenne} & Réel & elle peut avoir des décimales (ex. 13,50) \\
\hline
Admis ou non & \texttt{admis} & Booléen & seules deux valeurs possibles : vrai ou faux \\
\hline
\end{tabularx}
\end{center}

\textbf{2. Déclaration dans l'algorithme.}
\begin{lstlisting}
Algorithme Bulletin
Variables
    nomEleve : Chaîne
    age : Entier
    moyenne : Réel
    admis : Booléen
Début
    ...
Fin
\end{lstlisting}
\textbf{Conclusion} : chaque information possède un nom clair et un type adapté à sa nature ; c'est la première étape indispensable avant tout traitement.
\end{exoresolu}

\courssub{Méthode 3 : Utiliser les structures de contrôle (alternative et répétitive)}

\begin{methode}[Choisir entre Si...Alors, Pour et Tant que]
\begin{enumerate}
\item \textbf{Un choix à faire ?} Utilise la structure \cle{alternative} : \texttt{Si condition Alors ... Sinon ... FinSi}. La condition doit être une comparaison claire ($=$, $<$, $>$, $\leq$, $\geq$, $\neq$).
\item \textbf{Une action à répéter un nombre de fois connu à l'avance ?} Utilise la boucle \cle{Pour} : \texttt{Pour i allant de 1 à n Faire ... FinPour}.
\item \textbf{Une action à répéter tant qu'une condition est vraie (nombre de tours inconnu) ?} Utilise la boucle \cle{Tant que} : \texttt{Tant que condition Faire ... FinTantQue}.
\item \textbf{Vérifier la sortie de boucle} : dans un \texttt{Tant que}, quelque chose dans le corps de la boucle doit modifier la condition, sinon le programme ne s'arrête jamais (boucle infinie).
\end{enumerate}
\textbf{Réflexe} : pour calculer une somme ou un total, initialise toujours l'accumulateur à 0 \emph{avant} la boucle, puis ajoute chaque valeur \emph{dans} la boucle.
\end{methode}

\begin{exoresolu}[La cantine du collège]
La cantinière d'un collège de Bafoussam veut calculer la recette du jour. Elle saisit le montant payé par chacun des 30 élèves inscrits, et le programme affiche la somme totale encaissée.
\begin{enumerate}
\item Quelle structure de contrôle faut-il utiliser ? Justifier.
\item Écrire l'algorithme complet.
\item Qu'affiche le programme si les trois premiers montants sont 200, 150 et 250 FCFA et que tous les autres élèves paient 100 FCFA ?
\end{enumerate}
\tcblower
\textbf{1. Choix de la structure.} Le nombre de répétitions est connu à l'avance : exactement 30 élèves. On utilise donc une boucle \cle{Pour} (et non un \texttt{Tant que}, réservé aux cas où le nombre de tours est inconnu).

\textbf{2. Algorithme.}
\begin{lstlisting}
Algorithme RecetteCantine
Variables i, montant, total : Entier
Début
    total <- 0
    Pour i allant de 1 à 30 Faire
        Écrire("Montant payé par l'élève ", i, " : ")
        Lire(montant)
        total <- total + montant
    FinPour
    Écrire("Recette totale : ", total, " FCFA")
Fin
\end{lstlisting}
Remarque l'initialisation \texttt{total <- 0} avant la boucle : sans elle, le résultat serait faux.

\textbf{3. Calcul de la recette.}
\begin{itemize}
\item Les trois premiers élèves paient : $200 + 150 + 250 = 600$ FCFA.
\item Les $30 - 3 = 27$ autres élèves paient chacun 100 FCFA : $27 \times 100 = \num{2700}$ FCFA.
\item Total : $600 + \num{2700} = \num{3300}$ FCFA.
\end{itemize}
\textbf{Conclusion} : le programme affiche « Recette totale : 3300 FCFA ». La boucle \texttt{Pour} est la structure adaptée car le nombre d'élèves (30) était connu avant le début du traitement.
\end{exoresolu}

\courssub{Méthode 4 : Modulariser avec des fonctions et des procédures}

\begin{methode}[Diviser un programme en sous-programmes]
Principe « \cle{diviser pour régner} » : un gros problème se découpe en petites tâches indépendantes.
\begin{enumerate}
\item \textbf{Repérer les tâches} répétées ou autonomes (un calcul, un affichage) : chacune devient un sous-programme.
\item \textbf{Choisir entre fonction et procédure} :
\begin{itemize}
\item \cle{Fonction} : elle \emph{calcule} et \emph{retourne} un résultat (mot-clé \texttt{Retourner} / \texttt{return}) ;
\item \cle{Procédure} : elle \emph{effectue une action} (affichage, saisie) sans retourner de valeur.
\end{itemize}
\item \textbf{Définir les paramètres} : ce sont les données dont le sous-programme a besoin pour travailler.
\item \textbf{Appeler} le sous-programme depuis le programme principal autant de fois que nécessaire.
\end{enumerate}
\textbf{Avantages à citer dans une copie} : programme plus lisible, erreurs plus faciles à trouver, réutilisation du code, travail en équipe possible.
\end{methode}

\begin{exoresolu}[La moyenne de la classe]
Un enseignant veut un programme qui calcule la moyenne de deux notes pour chaque élève de sa classe.
\begin{enumerate}
\item Écrire une fonction \texttt{moyenne} qui reçoit deux notes et retourne leur moyenne.
\item Écrire le programme principal qui demande les deux notes d'un élève, appelle la fonction et affiche le résultat.
\item Citer deux avantages de cette organisation.
\end{enumerate}
\tcblower
\textbf{1. La fonction.} Elle calcule un résultat : c'est donc bien une \emph{fonction} (et non une procédure). Elle a deux paramètres réels et retourne un réel.
\begin{lstlisting}
Fonction moyenne(a : Réel, b : Réel) : Réel
Début
    Retourner (a + b) / 2
Fin
\end{lstlisting}

\textbf{2. Le programme principal.}
\begin{lstlisting}
Algorithme MoyenneClasse
Variables n1, n2, m : Réel
Début
    Écrire("Note 1 : ")
    Lire(n1)
    Écrire("Note 2 : ")
    Lire(n2)
    m <- moyenne(n1, n2)
    Écrire("La moyenne est : ", m)
Fin
\end{lstlisting}
\textbf{Vérification} : pour $n1 = 12$ et $n2 = 16$, la fonction retourne $(12 + 16) / 2 = 14$, et le programme affiche « La moyenne est : 14 ».

\textbf{3. Avantages.} (a) Le programme principal est court et lisible ; (b) la fonction \texttt{moyenne} peut être réutilisée pour tous les élèves sans réécrire le calcul, et une éventuelle erreur ne se corrige qu'à un seul endroit.

\textbf{Conclusion} : la modularité consiste à isoler chaque tâche dans un sous-programme : le programme devient plus clair, plus fiable et plus facile à faire évoluer.
\end{exoresolu}

\courssub{Méthode 5 : Tester, déboguer et documenter un programme}

\begin{methode}[Valider un programme avant de le livrer]
\begin{enumerate}
\item \textbf{Préparer un jeu de tests} : choisir des valeurs d'entrée \emph{dont on connaît le résultat exact à l'avance}, en prévoyant des cas normaux et des cas limites (0, valeurs maximales, valeurs particulières).
\item \textbf{Exécuter et comparer} : si le résultat affiché diffère du résultat attendu, il y a un \cle{bug}.
\item \textbf{Déboguer} : relire le code étape par étape (ou suivre les variables dans un tableau d'exécution) pour localiser l'instruction fautive, puis la corriger.
\item \textbf{Re-tester} après chaque correction : une correction peut créer une nouvelle erreur.
\item \textbf{Documenter} : écrire des \cle{commentaires} dans le code (rôle de chaque variable, de chaque bloc) et une petite notice pour l'utilisateur.
\end{enumerate}
\textbf{Réflexe} : un tableau d'exécution (une colonne par variable, une ligne par étape) est l'outil le plus sûr pour traquer une erreur de logique.
\end{methode}

\begin{exoresolu}[Le programme qui se trompe de mention]
Voici un programme écrit par un élève pour attribuer une mention selon la moyenne :
\begin{lstlisting}
Algorithme Mention
Variables m : Réel
Début
    Lire(m)
    Si m >= 10 Alors
        Écrire("Mention Passable")
    FinSi
    Si m >= 12 Alors
        Écrire("Mention Assez Bien")
    FinSi
Fin
\end{lstlisting}
\begin{enumerate}
\item Tester ce programme avec $m = 14$. Que constate-t-on ?
\item Expliquer l'erreur et proposer un programme corrigé.
\item Refaire le test avec $m = 14$ et $m = 8$ pour valider la correction.
\end{enumerate}
\tcblower
\textbf{1. Test avec $m = 14$.} La première condition $14 \geq 10$ est vraie : le programme affiche « Mention Passable ». Puis la deuxième condition $14 \geq 12$ est aussi vraie : il affiche « Mention Assez Bien ». Le programme affiche donc \emph{deux mentions contradictoires} : c'est un bug.

\textbf{2. Explication et correction.} L'élève a utilisé deux tests indépendants au lieu d'une alternative avec \texttt{Sinon}. Il faut tester d'abord la condition la plus forte, puis utiliser \texttt{Sinon} pour les autres cas :
\begin{lstlisting}
Algorithme MentionCorrige
Variables m : Réel
Début
    Lire(m)
    Si m >= 12 Alors
        Écrire("Mention Assez Bien")
    Sinon
        Si m >= 10 Alors
            Écrire("Mention Passable")
        Sinon
            Écrire("Moyenne insuffisante")
        FinSi
    FinSi
Fin
\end{lstlisting}

\textbf{3. Nouveaux tests.}
\begin{itemize}
\item Avec $m = 14$ : $14 \geq 12$ est vrai $\rightarrow$ affichage unique « Mention Assez Bien ». Correct.
\item Avec $m = 8$ : $8 \geq 12$ est faux, $8 \geq 10$ est faux $\rightarrow$ affichage « Moyenne insuffisante ». Correct.
\end{itemize}
\textbf{Conclusion} : le test avec une valeur connue a révélé le bug ; le débogage a montré que des conditions qui se recouvrent exigent une structure \texttt{Si ... Sinon ... FinSi} ; la re-validation confirme que le programme corrigé donne un seul résultat juste dans tous les cas.
\end{exoresolu}

\courssub{Méthode 6 : Rédiger et justifier une démarche ou une conclusion}

\begin{methode}[Rédiger une réponse complète et justifiée au BEPC]
\begin{enumerate}
\item \textbf{Annoncer} ce que l'on va faire : « Pour résoudre ce problème, j'utilise une boucle Pour car... ».
\item \textbf{Dérouler} la démarche dans l'ordre : analyse des données, choix des variables et des types, algorithme, traduction, test.
\item \textbf{Justifier chaque choix} par une règle du cours : type choisi selon la nature de la donnée, structure choisie selon que le nombre de répétitions est connu ou non, fonction choisie parce qu'un résultat est retourné.
\item \textbf{Valider} par un test chiffré : donner un exemple d'exécution avec le calcul complet.
\item \textbf{Conclure} par une phrase qui répond exactement à la question posée : « Le programme affiche donc... ».
\end{enumerate}
\textbf{Réflexe} : une réponse sans justification ni exemple de test perd des points, même si l'algorithme est correct.
\end{methode}

\begin{exoresolu}[Accès à la salle informatique]
Le club informatique d'un collège de Bafoussam réserve sa salle aux élèves de 12 ans et plus. Rédige une réponse complète et justifiée à la question : « Écrire un programme qui demande l'âge d'un élève et affiche s'il peut entrer dans la salle. »
\tcblower
\textbf{Démarche.} Le programme doit prendre une \emph{décision} selon une condition : il faut donc une structure alternative \texttt{Si ... Sinon ... FinSi}.

\textbf{Choix des variables.} Une seule donnée d'entrée : l'âge, noté \texttt{age}, de type \cle{Entier} car un âge est un nombre sans virgule.

\textbf{Algorithme.}
\begin{lstlisting}
Algorithme AccesSalle
Variables age : Entier
Début
    Écrire("Quel est votre âge ? ")
    Lire(age)
    Si age >= 12 Alors
        Écrire("Accès autorisé")
    Sinon
        Écrire("Accès refusé")
    FinSi
Fin
\end{lstlisting}
\textbf{Justification.} La condition « 12 ans et plus » se traduit par $\texttt{age} \geq 12$ (et non $\texttt{age} > 12$, qui exclurait les élèves de 12 ans exactement). Le \texttt{Sinon} garantit qu'un message s'affiche dans tous les cas.

\textbf{Validation par des tests.}
\begin{itemize}
\item Avec $\texttt{age} = 14$ : $14 \geq 12$ est vrai $\rightarrow$ « Accès autorisé ». Correct.
\item Avec $\texttt{age} = 12$ : $12 \geq 12$ est vrai $\rightarrow$ « Accès autorisé ». Correct (cas limite).
\item Avec $\texttt{age} = 10$ : $10 \geq 12$ est faux $\rightarrow$ « Accès refusé ». Correct.
\end{itemize}
\textbf{Conclusion} : le programme ci-dessus répond au besoin du club : il autorise l'accès à tout élève de 12 ans et plus, et le refuse aux autres, comme le prouvent les trois tests effectués, y compris le cas limite de 12 ans.
\end{exoresolu}

\begin{savaistu}[D'où vient le mot « bug » ?]
En 1947, les ingénieurs de l'ordinateur Harvard Mark II découvrirent qu'une panne était causée par un véritable insecte (une mite) coincé dans un relais de la machine. Depuis, le mot anglais \emph{bug} (« insecte ») désigne toute erreur dans un programme, et \emph{déboguer} signifie « enlever les insectes », c'est-à-dire corriger les erreurs. Aujourd'hui, au Cameroun comme partout dans le monde, les développeurs passent une grande partie de leur temps à tester et déboguer leurs applications avant de les mettre en service.
\end{savaistu}

\begin{attention}[Les 5 erreurs qui coûtent des points au BEPC]
\begin{enumerate}
\item \textbf{Oublier d'initialiser une variable} (par exemple \texttt{total <- 0} avant une boucle de somme) : le calcul devient faux ou imprévisible. Toute variable utilisée dans un calcul doit avoir reçu une valeur avant.
\item \textbf{Confondre $>$ et $\geq$} dans une condition : « 12 ans et plus » s'écrit $\texttt{age} \geq 12$. Tester systématiquement le cas limite (ici 12) pour vérifier le bon sens de l'inégalité.
\item \textbf{Choisir un mauvais type de variable} : déclarer une moyenne en Entier fait perdre les décimales (12,75 devient 12) ; déclarer un nom en nombre est absurde. Justifie toujours le type par la nature de la donnée.
\item \textbf{Écrire des tests indépendants qui se recouvrent} au lieu d'un \texttt{Si ... Sinon ... FinSi} : le programme affiche alors plusieurs messages contradictoires (voir l'exercice sur les mentions).
\item \textbf{Rendre un algorithme sans test ni conclusion} : une copie doit contenir la démarche (analyse, variables, algorithme), un exemple d'exécution chiffré qui valide le résultat, et une phrase de conclusion qui répond à la question posée.
\end{enumerate}
\end{attention}

\newpage

\coursec{Exercices}

\courssub{Je vérifie mes connaissances}

\begin{exercice}[QCM : Cycle de vie et vocabulaire]
Parmi les propositions suivantes, une seule est exacte pour chaque question. Recopie le numéro de la question et la lettre de la bonne réponse.

\begin{enumerate}
    \item Quelle phase du cycle de vie en V précède immédiatement la phase de \emph{codage} (programmation) ?
    \begin{itemize}
        \item[A.] Recueil des besoins
        \item[B.] Conception détaillée (algorithmes, structures de données)
        \item[C.] Tests d'intégration
        \item[D.] Maintenance corrective
    \end{itemize}

    \item Dans un algorithme, l'instruction \texttt{Afficher} correspond à quelle action concrète lors de l'exécution du programme ?
    \begin{itemize}
        \item[A.] Stocker une valeur en mémoire vive (RAM)
        \item[B.] Envoyer un résultat vers la sortie standard (l'écran)
        \item[C.] Lire une valeur saisie au clavier
        \item[D.] Effectuer un calcul arithmétique
    \end{itemize}

    \item La structure de contrôle qui permet de répéter un bloc d'instructions \emph{tant qu'une condition reste vraie} s'appelle :
    \begin{itemize}
        \item[A.] \texttt{SI ... ALORS ... SINON ... FIN SI}
        \item[B.] \texttt{POUR i DE 1 A N FAIRE ... FIN POUR}
        \item[C.] \texttt{TANT QUE condition FAIRE ... FIN TANT QUE}
        \item[D.] \texttt{SELON ... FAIRE ... FIN SELON}
    \end{itemize}

    \item Une \emph{fonction} se distingue d'une \emph{procédure} principalement parce qu'elle :
    \begin{itemize}
        \item[A.] Ne reçoit aucun paramètre en entrée
        \item[B.] Ne peut pas modifier de variables globales
        \item[C.] Renvoie une valeur (un résultat) utilisable dans une expression
        \item[D.] S'exécute plus rapidement
    \end{itemize}
\end{enumerate}
\end{exercice}

\begin{exercice}[Vrai ou Faux ? Justifie]
Pour chaque affirmation, réponds \textbf{Vrai} ou \textbf{Faux} et justifie ta réponse en une phrase précise, en utilisant le vocabulaire du cours.

\begin{enumerate}
    \item « L'algorithme est le code source écrit en langage C ou Python. »
    \item « Une variable de type \texttt{Entier} peut stocker la valeur \num{12,5}. »
    \item « La phase de \emph{tests unitaires} consiste à vérifier que le logiciel complet répond au cahier des charges du client. »
    \item « Dans une boucle \texttt{POUR i DE 1 A 10 FAIRE}, la variable \texttt{i} prend successivement les valeurs 1, 2, ..., 10. »
    \item « Le débogage consiste à écrire la documentation utilisateur du logiciel. »
\end{enumerate}
\end{exercice}

\begin{exercice}[Définitions et correspondance]
Associe chaque concept de la colonne de gauche à sa définition exacte dans la colonne de droite (écris les paires sous la forme : 1--c, 2--a, etc.).

\begin{center}
\begin{tabularx}{\linewidth}{|l|X|}
\hline
\rowcolor{popblueL}\textbf{Concept} & \textbf{Définition} \\
\hline
1. Variable & a. Phase qui transforme le code source en programme exécutable \\
\hline
2. Compilation & b. Suite ordonnée et finie d'opérations permettant de résoudre un problème \\
\hline
3. Algorithme & c. Zone de mémoire nommée, d'un type défini, dont la valeur peut changer \\
\hline
4. Test unitaire & d. Vérification isolée du bon fonctionnement d'une fonction ou d'une procédure \\
\hline
5. Cahier des charges & e. Document qui décrit précisément ce que le logiciel \emph{doit} faire \\
\hline
\end{tabularx}
\end{center}
\end{exercice}

\begin{exercice}[Lecture d'algorithme : calcul d'une somme]
Soit l'algorithme suivant écrit en pseudo-code :
\begin{lstlisting}
Algorithme SommePremiersEntiers
Variable
    n, i, somme : Entier
Debut
    Afficher("Entrez un entier positif N : ")
    Saisir(n)
    somme <- 0
    Pour i de 1 a n faire
        somme <- somme + i
    Fin Pour
    Afficher("La somme des ", n, " premiers entiers est : ", somme)
Fin
\end{lstlisting}
\begin{enumerate}
    \item Quelle est la valeur finale de \texttt{somme} si l'utilisateur saisit \texttt{5} ? Détaille le calcul.
    \item Quelle structure de contrôle est utilisée pour la répétition ?
    \item Quel serait l'affichage si l'utilisateur saisit \texttt{0} ? (On admet que dans ce cas la boucle \texttt{POUR} n'est pas exécutée.)
\end{enumerate}
\end{exercice}

\courssub{Je m'entraîne}

\begin{exercice}[Compléter un organigramme : maximum de trois nombres]
L'organigramme ci-dessous détermine le maximum de trois nombres \texttt{A}, \texttt{B}, \texttt{C}. Les losanges (tests) et les rectangles (actions) sont vides : à toi de les compléter.

\begin{popfigure}
\begin{tikzpicture}[
    node distance=1.2cm and 2.5cm,
    startstop/.style={rectangle, rounded corners, minimum width=3cm, minimum height=1cm, text centered, draw=popblue, fill=popblueL, thick},
    process/.style={rectangle, minimum width=3cm, minimum height=1cm, text centered, draw=popdark, fill=popcream, thick},
    decision/.style={diamond, aspect=2, minimum width=3.5cm, minimum height=1.5cm, text centered, draw=poporange, fill=poporangeL, thick, inner sep=0pt},
    arrow/.style={thick,->,>=Stealth, draw=popdark}
]
\node (start) [startstop] {Début};
\node (inA) [process, below of=start] {Saisir A, B, C};
\node (maxA) [process, right of=inA, xshift=3cm] {? ? ?};
\node (testB) [decision, below of=inA] {? ? ?};
\node (maxB) [process, right of=testB] {? ? ?};
\node (testC) [decision, below of=testB] {? ? ?};
\node (maxC) [process, right of=testC] {? ? ?};
\node (aff) [process, below of=testC] {Afficher max};
\node (end) [startstop, below of=aff] {Fin};

\draw [arrow] (start) -- (inA);
\draw [arrow] (inA) -- (maxA);
\draw [arrow] (maxA) |- (testB);
\draw [arrow] (testB) -- node[anchor=south] {Oui} (maxB);
\draw [arrow] (maxB) |- (testC);
\draw [arrow] (testB) -- node[anchor=east] {Non} (testC);
\draw [arrow] (testC) -- node[anchor=south] {Oui} (maxC);
\draw [arrow] (maxC) |- (aff);
\draw [arrow] (testC) -- node[anchor=east] {Non} (aff);
\draw [arrow] (aff) -- (end);
\end{tikzpicture}
\legende{Organigramme à compléter : recherche du maximum de trois nombres}
\end{popfigure}

Recopie l'organigramme sur ta copie et complète les zones vides (tests et actions) pour que l'algorithme soit correct. \emph{Indication : on utilise une variable \texttt{max} initialisée avec la valeur de \texttt{A}, puis on la compare successivement à \texttt{B} puis à \texttt{C}.}
\end{exercice}

\begin{exercice}[Écriture de fonction : année bissextile]
Écris en pseudo-code (syntaxe du cours : \texttt{Fonction}, \texttt{Debut}, \texttt{Fin}, \texttt{Retourner}) une fonction \texttt{EstBissextile(annee : Entier) : Booleen} qui retourne \texttt{VRAI} si l'année passée en paramètre est bissextile, \texttt{FAUX} sinon.

Rappel de la règle du calendrier grégorien :
\begin{itemize}
    \item une année est bissextile si elle est divisible par 4 \textbf{et} non divisible par 100 ;
    \item \textbf{sauf} si elle est divisible par 400 (alors elle est bissextile).
\end{itemize}
\emph{Indication :} le reste de la division euclidienne s'écrit avec l'opérateur \texttt{MOD} ; par exemple, « \texttt{annee} est divisible par 4 » s'écrit \texttt{annee MOD 4 = 0}.
\end{exercice}

\begin{exercice}[Conception de jeux de tests : décision d'admission]
On considère l'algorithme suivant qui détermine la décision du conseil de classe selon la moyenne d'un élève (sur 20) :
\begin{lstlisting}
Algorithme Mention
Variable
    note : Reel
    mention : Chaine
Debut
    Saisir(note)
    Si note >= 10 Alors
        mention <- "Admis"
    Sinon
        mention <- "Recale"
    Fin Si
    Afficher(mention)
Fin
\end{lstlisting}
Propose un jeu de \textbf{4 tests pertinents} sous forme de tableau. Pour chaque test, indique : le numéro du test, la valeur d'entrée (\texttt{note}), le résultat attendu (\texttt{mention}) et le \textbf{type de test} (cas normal, valeur limite basse, valeur limite haute, valeur invalide). Justifie brièvement le choix de chaque valeur.
\end{exercice}

\begin{exercice}[Algorithme complet : moyenne de la classe]
Écris un algorithme complet (avec \texttt{Algorithme}, \texttt{Variable}, \texttt{Debut}, \texttt{Fin}) qui :
\begin{enumerate}
    \item demande à l'utilisateur le nombre d'élèves \texttt{N} (entier strictement positif) ;
    \item pour chaque élève de 1 à \texttt{N} (boucle \texttt{POUR}), demande sa note (réel compris entre 0 et 20) et l'ajoute à une variable \texttt{somme} ;
    \item calcule la moyenne : \texttt{moyenne <- somme / N} ;
    \item affiche la moyenne de la classe avec un message clair.
\end{enumerate}
Veille à initialiser correctement la variable \texttt{somme} avant la boucle.
\end{exercice}

\courssub{Je me prépare au BEPC}

\begin{exercice}[Sujet type BEPC : gestion d'un stock de boutique — 12 pts]
M. Kamga gère une petite boutique à Yaoundé. Il souhaite un programme simple pour suivre le stock de ses 5 articles principaux. On modélise le stock par deux tableaux de taille 5 :
\begin{itemize}
    \item \texttt{noms[5]} : chaînes de caractères (exemples : "Riz", "Huile", "Savon", "Pates", "Lait") ;
    \item \texttt{qt[5]} : entiers (quantités en stock).
\end{itemize}
Les indices des tableaux vont de 0 à 4.

\begin{enumerate}
    \item \textbf{Initialisation (2 pts).} Écris une procédure \texttt{InitialiserStock()} qui remplit \texttt{noms} avec les 5 articles ci-dessus et met toutes les quantités \texttt{qt} à 0.
    \item \textbf{Affichage (2 pts).} Écris une procédure \texttt{AfficherStock()} qui affiche proprement le contenu du stock, ligne par ligne (exemple : \texttt{0 : Riz -- 50 kg}).
    \item \textbf{Vente (3 pts).} Écris une fonction \texttt{Vendre(indice : Entier, qteVendue : Entier) : Booleen} qui :
    \begin{itemize}
        \item vérifie que \texttt{indice} est valide (compris entre 0 et 4) et que \texttt{qteVendue} $> 0$ ;
        \item vérifie que le stock est suffisant (\texttt{qt[indice] >= qteVendue}) ;
        \item si tout est correct, diminue le stock et retourne \texttt{VRAI} ; sinon affiche un message d'erreur adapté et retourne \texttt{FAUX}.
    \end{itemize}
    \item \textbf{Algorithme principal (3 pts).} Écris l'algorithme principal qui :
    \begin{itemize}
        \item appelle \texttt{InitialiserStock()} ;
        \item affiche le stock initial ;
        \item demande à l'utilisateur un indice d'article et une quantité à vendre ;
        \item appelle \texttt{Vendre} et affiche « Vente enregistrée » ou « Échec de la vente » selon la valeur retournée ;
        \item affiche le stock final.
    \end{itemize}
    \item \textbf{Analyse (2 pts).} Pourquoi est-il préférable d'utiliser une \emph{fonction} retournant un booléen pour \texttt{Vendre} plutôt qu'une simple procédure ? Quel avantage de la modularité cela illustre-t-il ?
\end{enumerate}
\end{exercice}

\begin{exercice}[Sujet type BEPC : mini-jeu « Pierre--Feuille--Ciseaux » — 10 pts]
On veut programmer le jeu classique contre l'ordinateur. Règles : la Pierre bat les Ciseaux, les Ciseaux battent la Feuille, la Feuille bat la Pierre. Si les deux choix sont identiques, il y a match nul.

\begin{enumerate}
    \item \textbf{Choix de l'ordinateur (2 pts).} Écris une fonction \texttt{ChoixOrdi() : Entier} qui retourne un entier aléatoire entre 1 et 3 (1 = Pierre, 2 = Feuille, 3 = Ciseaux). \emph{Indication :} on dispose de la fonction \texttt{ALEA(1, 3)} qui fournit un entier au hasard entre 1 et 3.
    \item \textbf{Comparaison (4 pts).} Écris une fonction \texttt{Resultat(joueur : Entier, ordi : Entier) : Chaine} qui retourne \texttt{"Gagne"}, \texttt{"Perdu"} ou \texttt{"Nul"} selon les règles du jeu. Utilise une structure \texttt{SI ... ALORS ... SINON} claire (éventuellement avec des \texttt{SI} imbriqués).
    \item \textbf{Boucle de jeu (2 pts).} Écris l'algorithme principal qui :
    \begin{itemize}
        \item initialise \texttt{scoreJoueur} et \texttt{scoreOrdi} à 0 ;
        \item répète (boucle \texttt{TANT QUE}) : demande le choix du joueur (1, 2 ou 3), appelle \texttt{ChoixOrdi}, appelle \texttt{Resultat}, met à jour les scores et affiche le résultat du tour ;
        \item la boucle s'arrête dès que l'un des deux scores atteint 3 ;
        \item affiche le grand gagnant final.
    \end{itemize}
    \item \textbf{Amélioration (2 pts).} Propose une modification pour gérer le cas où l'utilisateur saisit une valeur invalide (par exemple 5 ou $-1$) sans fausser le jeu. Décris la logique en français ou en pseudo-code (inutile de réécrire tout le programme).
\end{enumerate}
\end{exercice}

\begin{exercice}[Sujet type BEPC : analyse d'un algorithme de tri simple — 8 pts]
On donne l'algorithme suivant qui trie un tableau \texttt{T} de \texttt{N} entiers par ordre croissant (tri par sélection) :

\begin{lstlisting}
Algorithme TriSelection
Variable
    T : Tableau[0..N-1] de Entier
    N, i, j, minIdx, temp : Entier
Debut
    // On suppose N et le tableau T deja saisis
    Pour i de 0 a N-2 faire
        minIdx <- i
        Pour j de i+1 a N-1 faire
            Si T[j] < T[minIdx] Alors
                minIdx <- j
            Fin Si
        Fin Pour
        Si minIdx <> i Alors
            temp <- T[i]
            T[i] <- T[minIdx]
            T[minIdx] <- temp
        Fin Si
    Fin Pour
    Afficher("Tableau trie")
Fin
\end{lstlisting}

\begin{enumerate}
    \item \textbf{Exécution pas à pas (3 pts).} Soit \texttt{N = 5} et \texttt{T = [4, 2, 5, 1, 3]}. Complète le tableau de trace ci-dessous pour les deux premières itérations de la boucle externe (\texttt{i = 0} puis \texttt{i = 1}). Indique la valeur finale de \texttt{minIdx}, l'échange éventuellement effectué (indices concernés) et l'état du tableau \texttt{T} après l'itération.

    \begin{center}
    \begin{tabularx}{\linewidth}{|c|c|X|X|}
    \hline
    \rowcolor{popblueL}\textbf{i} & \textbf{minIdx final} & \textbf{Échange ? (indices)} & \textbf{T après l'itération} \\
    \hline
    0 & & & \\
    \hline
    1 & & & \\
    \hline
    \end{tabularx}
    \end{center}

    \item \textbf{Nombre de comparaisons (2 pts).} Combien de comparaisons \texttt{T[j] < T[minIdx]} sont effectuées au total pour un tableau de taille \texttt{N} ? Donne la formule sous forme d'une somme en fonction de \texttt{N}, puis sa forme simplifiée.
    \item \textbf{Propriété de la boucle (2 pts).} Énonce en français clair la propriété vérifiée au début de chaque itération \texttt{i} de la boucle externe : que peut-on dire de la partie gauche du tableau (indices $0$ à $i-1$) ?
    \item \textbf{Correction (1 pt).} Si l'on remplace la condition \texttt{Si minIdx <> i Alors} par \texttt{Si T[minIdx] < T[i] Alors}, l'algorithme reste-t-il correct ? Justifie ta réponse.
\end{enumerate}
\end{exercice}

\newpage

\coursec{Corrigés des exercices}

\begin{corrige}[Exercice 1 — QCM : Cycle de vie et vocabulaire]
\begin{enumerate}
    \item \textbf{B.} Dans le cycle de vie en V, la phase de \emph{conception détaillée} (choix des algorithmes et des structures de données) précède immédiatement le codage : on ne programme que ce qui a été conçu. Le recueil des besoins (A) est la toute première phase ; les tests (C) et la maintenance (D) viennent \emph{après} le codage.
    \item \textbf{B.} \texttt{Afficher} envoie un résultat vers la sortie standard, c'est-à-dire l'écran. Stocker une valeur (A) correspond à l'affectation \texttt{<-} ; lire au clavier (C) correspond à \texttt{Saisir}.
    \item \textbf{C.} La boucle \texttt{TANT QUE condition FAIRE ... FIN TANT QUE} répète le bloc tant que la condition reste vraie (le nombre de répétitions n'est pas connu à l'avance). La boucle \texttt{POUR} (B) s'utilise quand le nombre d'itérations est connu ; \texttt{SI} (A) et \texttt{SELON} (D) sont des structures conditionnelles, pas des répétitions.
    \item \textbf{C.} Une fonction \emph{retourne une valeur} (avec \texttt{Retourner}) utilisable dans une expression, par exemple \texttt{m <- Max(a, b)}. Une procédure effectue des actions mais ne renvoie rien.
\end{enumerate}
\textbf{Réponses : 1--B, 2--B, 3--C, 4--C.}
\end{corrige}

\begin{corrige}[Exercice 2 — Vrai ou Faux ? Justifie]
\begin{enumerate}
    \item \textbf{Faux.} L'algorithme est une suite ordonnée et finie d'opérations décrite en pseudo-code, indépendante de tout langage ; le code source en C ou en Python est la \emph{traduction} (l'implémentation) de cet algorithme dans un langage de programmation.
    \item \textbf{Faux.} Une variable de type \texttt{Entier} ne peut stocker que des nombres sans partie décimale ; la valeur \num{12,5} exige le type \texttt{Reel}.
    \item \textbf{Faux.} Les tests unitaires vérifient \emph{isolément} chaque fonction ou procédure ; c'est la phase de \emph{validation} (ou recette) qui vérifie que le logiciel complet répond au cahier des charges du client.
    \item \textbf{Vrai.} La boucle \texttt{POUR} fait prendre au compteur \texttt{i} toutes les valeurs entières de la borne initiale (1) à la borne finale (10), en augmentant de 1 à chaque tour.
    \item \textbf{Faux.} Le débogage consiste à rechercher et corriger les erreurs (bogues) d'un programme, souvent à l'aide d'un débogueur ; la rédaction de la documentation utilisateur relève de la phase de documentation, qui est distincte.
\end{enumerate}
\end{corrige}

\begin{corrige}[Exercice 3 — Définitions et correspondance]
\begin{itemize}
    \item \textbf{1--c} : une variable est une zone de mémoire nommée, d'un type défini, dont la valeur peut changer au cours de l'exécution.
    \item \textbf{2--a} : la compilation est la phase qui transforme le code source en programme exécutable (langage machine).
    \item \textbf{3--b} : un algorithme est une suite ordonnée et finie d'opérations permettant de résoudre un problème.
    \item \textbf{4--d} : un test unitaire est la vérification isolée du bon fonctionnement d'une fonction ou d'une procédure.
    \item \textbf{5--e} : le cahier des charges est le document qui décrit précisément ce que le logiciel doit faire.
\end{itemize}
\textbf{Réponses : 1--c, 2--a, 3--b, 4--d, 5--e.}
\end{corrige}

\begin{corrige}[Exercice 4 — Lecture d'algorithme : calcul d'une somme]
\begin{enumerate}
    \item Pour \texttt{n = 5}, la variable \texttt{somme} est initialisée à 0, puis la boucle ajoute successivement les valeurs de \texttt{i} :
    \begin{align*}
        i = 1 &: \quad somme = 0 + 1 = 1\\
        i = 2 &: \quad somme = 1 + 2 = 3\\
        i = 3 &: \quad somme = 3 + 3 = 6\\
        i = 4 &: \quad somme = 6 + 4 = 10\\
        i = 5 &: \quad somme = 10 + 5 = 15
    \end{align*}
    La valeur finale de \texttt{somme} est donc \textbf{15} (on vérifie : $1+2+3+4+5 = 15$).
    \item La structure de contrôle utilisée est la boucle \textbf{\texttt{POUR}} (boucle itérative à compteur), car le nombre de répétitions (\texttt{n}) est connu avant d'entrer dans la boucle.
    \item Si l'utilisateur saisit \texttt{0}, la boucle n'est pas exécutée et \texttt{somme} conserve sa valeur initiale 0. L'affichage est : \texttt{La somme des 0 premiers entiers est : 0}. Le programme reste correct grâce à l'initialisation \texttt{somme <- 0} avant la boucle.
\end{enumerate}
\end{corrige}

\begin{corrige}[Exercice 5 — Compléter un organigramme : maximum de trois nombres]
On suit l'indication : on initialise \texttt{max} avec \texttt{A}, puis on le compare à \texttt{B} puis à \texttt{C}.
\begin{itemize}
    \item Premier rectangle vide (après la saisie) : \textbf{\texttt{max <- A}} ;
    \item Premier losange : \textbf{\texttt{B > max ?}} — si Oui, rectangle : \textbf{\texttt{max <- B}} ; si Non, on passe directement au test suivant ;
    \item Deuxième losange : \textbf{\texttt{C > max ?}} — si Oui, rectangle : \textbf{\texttt{max <- C}} ; si Non, on passe directement à l'affichage ;
    \item Ensuite : \texttt{Afficher max}, puis \texttt{Fin}.
\end{itemize}
L'algorithme équivalent en pseudo-code est :
\begin{lstlisting}
Algorithme Maximum3
Variable
    A, B, C, max : Reel
Debut
    Saisir(A, B, C)
    max <- A
    Si B > max Alors
        max <- B
    Fin Si
    Si C > max Alors
        max <- C
    Fin Si
    Afficher("Le maximum est : ", max)
Fin
\end{lstlisting}
\textbf{Vérification} avec $A = 7$, $B = 12$, $C = 5$ : \texttt{max} vaut 7, puis 12 (car $12 > 7$), et reste 12 (car $5 < 12$). Affichage : 12, correct.
\end{corrige}

\begin{corrige}[Exercice 6 — Écriture de fonction : année bissextile]
\begin{lstlisting}
Fonction EstBissextile(annee : Entier) : Booleen
Debut
    Si (annee MOD 400 = 0) Alors
        Retourner VRAI
    Sinon
        Si (annee MOD 4 = 0) ET (annee MOD 100 <> 0) Alors
            Retourner VRAI
        Sinon
            Retourner FAUX
        Fin Si
    Fin Si
Fin
\end{lstlisting}
\textbf{Justification.} On teste d'abord le cas exceptionnel : une année divisible par 400 est toujours bissextile (ex. 2000). Sinon, il faut être divisible par 4 \emph{sans} être divisible par 100 (ex. 2024 est bissextile ; 1900 ne l'est pas, car $1900 \bmod 100 = 0$). Tous les autres cas retournent \texttt{FAUX}.

\textbf{Jeux de vérification} : \texttt{EstBissextile(2024)} = VRAI ; \texttt{EstBissextile(1900)} = FAUX ; \texttt{EstBissextile(2000)} = VRAI ; \texttt{EstBissextile(2023)} = FAUX.
\end{corrige}

\begin{corrige}[Exercice 7 — Conception de jeux de tests : décision d'admission]
\begin{center}
\begin{tabularx}{\linewidth}{|c|c|c|c|X|}
\hline
\rowcolor{popblueL}\textbf{N°} & \textbf{note} & \textbf{Résultat attendu} & \textbf{Type de test} & \textbf{Justification} \\
\hline
1 & 12 & Admis & Cas normal & Valeur courante au-dessus du seuil : vérifie le comportement général. \\
\hline
2 & 10 & Admis & Valeur limite basse & Le seuil exact : la condition est \texttt{>= 10}, donc 10 doit donner « Admis ». C'est le test le plus important (erreur classique : écrire \texttt{>} au lieu de \texttt{>=}). \\
\hline
3 & 9,99 & Recale & Valeur limite haute & Juste en dessous du seuil : vérifie que la frontière est bien placée. \\
\hline
4 & $-3$ & Recale (ou message d'erreur) & Valeur invalide & Une note négative est impossible : le programme ne doit pas « admettre » une telle valeur ; idéalement il devrait la rejeter. \\
\hline
\end{tabularx}
\end{center}
\textbf{Points de vigilance.} Un bon jeu de tests explore toujours les \emph{valeurs limites} (autour du seuil 10) et pas seulement des cas « évidents ». On peut ajouter 20 (note maximale) comme valeur limite haute du domaine valide, et 25 comme valeur invalide supérieure.
\end{corrige}

\begin{corrige}[Exercice 8 — Algorithme complet : moyenne de la classe]
\begin{lstlisting}
Algorithme MoyenneClasse
Variable
    N, i : Entier
    note, somme, moyenne : Reel
Debut
    Afficher("Nombre d'eleves : ")
    Saisir(N)
    somme <- 0
    Pour i de 1 a N faire
        Afficher("Note de l'eleve ", i, " : ")
        Saisir(note)
        somme <- somme + note
    Fin Pour
    moyenne <- somme / N
    Afficher("La moyenne de la classe est : ", moyenne, " / 20")
Fin
\end{lstlisting}
\textbf{Points de vigilance.}
\begin{itemize}
    \item \texttt{somme} est déclarée en \texttt{Reel} (les notes peuvent être décimales) et \textbf{initialisée à 0 avant la boucle} : sans cela, le résultat serait imprévisible.
    \item La division \texttt{somme / N} est placée \emph{après} la boucle, jamais à l'intérieur.
    \item \texttt{N} étant supposé strictement positif (énoncé), il n'y a pas de division par zéro ; dans un programme réel, on vérifierait \texttt{N > 0} avant de calculer.
\end{itemize}
\textbf{Vérification} : pour $N = 3$ et les notes 10, 12, 14 : $somme = 36$, $moyenne = 36 / 3 = 12$. Correct.
\end{corrige}

\begin{corrige}[Exercice 9 — Sujet type BEPC : gestion d'un stock de boutique]
\textbf{1. Initialisation (2 pts).}
\begin{lstlisting}
Procedure InitialiserStock()
Variable
    i : Entier
Debut
    noms[0] <- "Riz"
    noms[1] <- "Huile"
    noms[2] <- "Savon"
    noms[3] <- "Pates"
    noms[4] <- "Lait"
    Pour i de 0 a 4 faire
        qt[i] <- 0
    Fin Pour
Fin
\end{lstlisting}

\textbf{2. Affichage (2 pts).}
\begin{lstlisting}
Procedure AfficherStock()
Variable
    i : Entier
Debut
    Pour i de 0 a 4 faire
        Afficher(i, " : ", noms[i], " -- ", qt[i], " kg")
    Fin Pour
Fin
\end{lstlisting}

\textbf{3. Vente (3 pts).}
\begin{lstlisting}
Fonction Vendre(indice : Entier, qteVendue : Entier) : Booleen
Debut
    Si (indice < 0) OU (indice > 4) Alors
        Afficher("Erreur : indice invalide")
        Retourner FAUX
    Fin Si
    Si qteVendue <= 0 Alors
        Afficher("Erreur : quantite invalide")
        Retourner FAUX
    Fin Si
    Si qt[indice] < qteVendue Alors
        Afficher("Erreur : stock insuffisant")
        Retourner FAUX
    Fin Si
    qt[indice] <- qt[indice] - qteVendue
    Retourner VRAI
Fin
\end{lstlisting}

\textbf{4. Algorithme principal (3 pts).}
\begin{lstlisting}
Algorithme GestionStock
Variable
    noms : Tableau[0..4] de Chaine
    qt : Tableau[0..4] de Entier
    indice, qte : Entier
    ok : Booleen
Debut
    InitialiserStock()
    Afficher("--- Stock initial ---")
    AfficherStock()
    Afficher("Indice de l'article (0 a 4) : ")
    Saisir(indice)
    Afficher("Quantite a vendre : ")
    Saisir(qte)
    ok <- Vendre(indice, qte)
    Si ok = VRAI Alors
        Afficher("Vente enregistree")
    Sinon
        Afficher("Echec de la vente")
    Fin Si
    Afficher("--- Stock final ---")
    AfficherStock()
Fin
\end{lstlisting}

\textbf{5. Analyse (2 pts).} Une fonction retournant un booléen permet au programme appelant de \emph{connaître le résultat de l'opération} et de décider quoi faire ensuite (afficher « Vente enregistrée » ou « Échec ») : la décision est réutilisable et testable indépendamment. Une simple procédure effectuerait l'action « en silence » sans informer l'appelant. Cela illustre l'avantage de la \textbf{modularité} : chaque fonction a un rôle précis, peut être testée isolément (test unitaire) et réutilisée dans d'autres parties du programme.

\textbf{Barème indicatif} : 1. 2 pts (affectations 1 pt, boucle de mise à zéro 1 pt) ; 2. 2 pts ; 3. 3 pts (1 pt par vérification, retour booléen compris) ; 4. 3 pts (appels et enchaînement corrects) ; 5. 2 pts (rôle du booléen 1 pt, modularité 1 pt).

\textbf{Points de vigilance.} Les indices vont de 0 à 4 (pas de 1 à 5) ; les tableaux sont des variables globales accessibles dans les procédures ; la modification du stock ne doit avoir lieu \emph{qu'après} toutes les vérifications.
\end{corrige}

\begin{corrige}[Exercice 10 — Sujet type BEPC : mini-jeu « Pierre--Feuille--Ciseaux »]
\textbf{1. Choix de l'ordinateur (2 pts).}
\begin{lstlisting}
Fonction ChoixOrdi() : Entier
Debut
    Retourner ALEA(1, 3)
Fin
\end{lstlisting}

\textbf{2. Comparaison (4 pts).}
\begin{lstlisting}
Fonction Resultat(joueur : Entier, ordi : Entier) : Chaine
Debut
    Si joueur = ordi Alors
        Retourner "Nul"
    Sinon
        Si (joueur = 1 ET ordi = 3) OU (joueur = 2 ET ordi = 1) OU (joueur = 3 ET ordi = 2) Alors
            Retourner "Gagne"
        Sinon
            Retourner "Perdu"
        Fin Si
    Fin Si
Fin
\end{lstlisting}
\textbf{Justification.} Le joueur gagne dans exactement trois cas : Pierre (1) contre Ciseaux (3), Feuille (2) contre Pierre (1), Ciseaux (3) contre Feuille (2). L'égalité étant traitée avant, tous les cas restants sont perdus.

\textbf{3. Boucle de jeu (2 pts).}
\begin{lstlisting}
Algorithme PierreFeuilleCiseaux
Variable
    joueur, ordi, scoreJoueur, scoreOrdi : Entier
    res : Chaine
Debut
    scoreJoueur <- 0
    scoreOrdi <- 0
    Tant Que (scoreJoueur < 3) ET (scoreOrdi < 3) Faire
        Afficher("Votre choix (1=Pierre, 2=Feuille, 3=Ciseaux) : ")
        Saisir(joueur)
        ordi <- ChoixOrdi()
        res <- Resultat(joueur, ordi)
        Si res = "Gagne" Alors
            scoreJoueur <- scoreJoueur + 1
        Fin Si
        Si res = "Perdu" Alors
            scoreOrdi <- scoreOrdi + 1
        Fin Si
        Afficher("Tour : ", res, " -- Score : ", scoreJoueur, " / ", scoreOrdi)
    Fin Tant Que
    Si scoreJoueur = 3 Alors
        Afficher("Vous etes le grand gagnant !")
    Sinon
        Afficher("L'ordinateur gagne la partie.")
    Fin Si
Fin
\end{lstlisting}
\textbf{Vérification de la condition d'arrêt} : la boucle continue tant qu'\emph{aucun} des deux scores n'a atteint 3 ; dès que l'un atteint 3, la condition \texttt{(scoreJoueur < 3) ET (scoreOrdi < 3)} devient fausse et la boucle s'arrête.

\textbf{4. Amélioration (2 pts).} On contrôle la saisie avec une boucle \texttt{TANT QUE} : tant que la valeur saisie n'est pas comprise entre 1 et 3, on affiche un message d'erreur et on redemande le choix. Ainsi, une valeur invalide (5, $-1$) ne déclenche ni tour de jeu ni modification des scores.
\begin{lstlisting}
Saisir(joueur)
Tant Que (joueur < 1) OU (joueur > 3) Faire
    Afficher("Choix invalide, recommencez : ")
    Saisir(joueur)
Fin Tant Que
\end{lstlisting}

\textbf{Barème indicatif} : 1. 2 pts ; 2. 4 pts (cas nul 1 pt, trois cas gagnants 2 pts, cas perdus 1 pt) ; 3. 2 pts (condition d'arrêt 1 pt, mise à jour des scores 1 pt) ; 4. 2 pts.

\textbf{Points de vigilance.} Bien initialiser les scores \emph{avant} la boucle ; ne pas oublier que le match nul ne modifie aucun score ; la condition du \texttt{TANT QUE} utilise \texttt{ET} (les deux scores doivent être $< 3$ pour continuer).
\end{corrige}

\begin{corrige}[Exercice 11 — Sujet type BEPC : analyse d'un algorithme de tri simple]
\textbf{1. Exécution pas à pas (3 pts).} Tableau initial : \texttt{T = [4, 2, 5, 1, 3]}.

\emph{Itération i = 0} : on cherche le minimum dans \texttt{T[0..4]}. Les valeurs sont 4, 2, 5, 1, 3 ; le minimum est 1, à l'indice 3, donc \texttt{minIdx = 3}. Comme $minIdx \neq i$, on échange \texttt{T[0]} et \texttt{T[3]} : \texttt{T = [1, 2, 5, 4, 3]}.

\emph{Itération i = 1} : on cherche le minimum dans \texttt{T[1..4]}, c'est-à-dire parmi 2, 5, 4, 3. Le minimum est 2, déjà à l'indice 1, donc \texttt{minIdx = 1}. Comme $minIdx = i$, aucun échange : \texttt{T = [1, 2, 5, 4, 3]}.

\begin{center}
\begin{tabularx}{\linewidth}{|c|c|X|X|}
\hline
\rowcolor{popblueL}\textbf{i} & \textbf{minIdx final} & \textbf{Échange ? (indices)} & \textbf{T après l'itération} \\
\hline
0 & 3 & Oui : indices 0 et 3 & [1, 2, 5, 4, 3] \\
\hline
1 & 1 & Non (minIdx = i) & [1, 2, 5, 4, 3] \\
\hline
\end{tabularx}
\end{center}

\textbf{2. Nombre de comparaisons (2 pts).} Pour chaque valeur de \texttt{i}, la boucle interne effectue une comparaison pour chaque \texttt{j} de $i+1$ à $N-1$, soit $(N-1) - (i+1) + 1 = N - 1 - i$ comparaisons. Au total :
\[ \sum_{i=0}^{N-2} (N - 1 - i) = (N-1) + (N-2) + \dots + 1 = \frac{N(N-1)}{2}. \]
\textbf{Vérification} avec $N = 5$ : $\dfrac{5 \times 4}{2} = 10$ comparaisons ($4 + 3 + 2 + 1 = 10$). Correct.

\textbf{3. Propriété de la boucle (2 pts).} Au début de chaque itération \texttt{i}, la partie gauche du tableau (indices $0$ à $i-1$) contient les $i$ plus petits éléments du tableau, \textbf{déjà triés par ordre croissant et placés à leur position définitive}. C'est l'\emph{invariant de boucle} : chaque itération agrandit cette zone triée d'un élément.

\textbf{4. Correction (1 pt).} \textbf{Oui, l'algorithme reste correct.} Par construction, \texttt{T[minIdx]} est le minimum de la zone \texttt{T[i..N-1]}, donc on a toujours \texttt{T[minIdx]} $\leq$ \texttt{T[i]}. Deux cas : si \texttt{minIdx = i}, alors \texttt{T[minIdx] = T[i]}, la condition \texttt{T[minIdx] < T[i]} est fausse et l'on n'échange pas (comme avant) ; si \texttt{minIdx} $\neq$ \texttt{i}, alors \texttt{T[minIdx] < T[i]} et l'échange a lieu (comme avant). Le comportement est donc identique.

\textbf{Barème indicatif} : 1. 3 pts (1,5 pt par ligne : minIdx, échange, état du tableau) ; 2. 2 pts (somme 1 pt, forme simplifiée 1 pt) ; 3. 2 pts ; 4. 1 pt.

\textbf{Points de vigilance.} La boucle externe s'arrête à $N-2$ (le dernier élément est automatiquement à sa place) ; la recherche du minimum porte sur la zone \emph{non triée} (de \texttt{i} à $N-1$) ; dans la formule, ne pas oublier que le nombre de termes de la somme est $N - 1$.
\end{corrige}

\newpage

\coursec{Fiche bilan}

\begin{popfigure}
\centering
\begin{tikzpicture}[
    mindmap,
    grow cyclic,
    every node/.style={concept, execute at begin node=\hskip0pt},
    concept color=popblue,
    level 1/.append style={level distance=4.5cm, sibling angle=72},
    level 2/.append style={level distance=3cm, sibling angle=45},
    branch1/.style={concept color=poporange, font=\small\bfseries},
    branch2/.style={concept color=popgreen, font=\small\bfseries},
    branch3/.style={concept color=poppurple, font=\small\bfseries},
    branch4/.style={concept color=poppink, font=\small\bfseries},
    branch5/.style={concept color=popgold, font=\small\bfseries},
    branch6/.style={concept color=popdark, font=\small\bfseries},
]
\node[concept, text width=2.8cm, font=\bfseries\large] {Développement\\ Logiciel}
    child[branch1] { node[concept, text width=2.5cm] {1. Cycle de vie\\ (Analyse $\to$ Maintenance)} }
    child[branch2] { node[concept, text width=2.5cm] {2. Algo $\to$ Programme\\ (Langages, Traduction)} }
    child[branch3] { node[concept, text width=2.5cm] {3. Briques de base\\ (Variables, Types, Contrôle)} }
    child[branch4] { node[concept, text width=2.5cm] {4. Modularité\\ (Fonctions, Procédures)} }
    child[branch5] { node[concept, text width=2.5cm] {5. Qualité\\ (Tests, Débogage, Doc)} }
    child[branch6] { node[concept, text width=2.5cm] {6. Outils\\ (Éditeur, IDE, Compilateur)} };
\end{tikzpicture}
\legende{Carte mentale : Concepts de base du développement logiciel (Leçon 15)}
\end{popfigure}

\begin{bilan}

\courssub{Définitions clés}

\begin{itemize}
    \item \cle{Logiciel} : Ensemble de programmes, procédures, règles et documentation relatifs au fonctionnement d'un système de traitement de l'information.
    \item \cle{Algorithme} : Suite finie et non ambiguë d'instructions (opérations élémentaires) permettant de résoudre un problème ou d'obtenir un résultat. Écrit en langage naturel ou pseudo-code.
    \item \cle{Programme} : Traduction d'un algorithme dans un \cle{langage de programmation} (C, Python, JavaScript, etc.) compréhensible par la machine après traduction.
    \item \cle{Code source} : Texte du programme écrit par le développeur (lisible par l'humain).
    \item \cle{Code objet / Exécutable} : Résultat de la traduction (binaire), directement exécutable par le processeur.
    \item \cle{Variable} : Zone mémoire nommée, identifiée par un \cle{type} (entier, réel, chaîne, booléen), dont la valeur peut changer pendant l'exécution.
    \item \cle{Structure de contrôle} : Instruction qui détermine l'ordre d'exécution des autres instructions (Séquence, Sélection \texttt{Si}, Itération \texttt{Pour}/\texttt{TantQue}).
    \item \cle{Fonction} : Sous-programme qui effectue un traitement et \emph{renvoie une valeur} (un seul résultat).
    \item \cle{Procédure} : Sous-programme qui effectue un traitement \emph{sans renvoyer de valeur} (effet de bord : affichage, modification de variables globales/paramètres).
    \item \cle{Bogue (Bug)} : Erreur dans le code source provoquant un comportement inattendu. \cle{Débogage} : Processus de localisation et de correction des bogues.
\end{itemize}

\courssub{Repères essentiels (à connaître par cœur)}

\begin{center}
\begin{tabularx}{\linewidth}{|>{\bfseries}l|X|X|}
\hline
\rowcolor{popblueL}
\textbf{Concept} & \textbf{Description / Règle} & \textbf{Exemple / Note} \\
\hline
\textbf{Cycle de vie (en V ou cascade)} & 1. Analyse des besoins \newline 2. Conception (Algo/MCD) \newline 3. Codage (Programmation) \newline 4. Tests (Unitaires, Intégration, Validation) \newline 5. Déploiement \& Maintenance & Itératif dans la pratique moderne (Agile). \\
\hline
\textbf{Traduction} & \textbf{Compilateur} : Traduit tout le code source en exécutable \emph{avant} exécution (C, C++). Rapide à l'exécution. \newline \textbf{Interpréteur} : Traduit et exécute instruction par instruction (Python, JS, BASIC). Plus lent, mais interactif. & Erreur de compilation = syntaxe. Erreur à l'exécution = logique/exception. \\
\hline
\textbf{Types de base} & \texttt{Entier} (\texttt{int}), \texttt{Réel} (\texttt{float/double}), \texttt{Caractère} (\texttt{char}), \texttt{Chaîne} (\texttt{string}), \texttt{Booléen} (\texttt{bool} : Vrai/Faux) & Typage fort (C) vs faible/dynamique (Python, JS). En 3ème : on précise le type en pseudo-code. \\
\hline
\textbf{Structures de contrôle} & \textbf{Séquence} : Instructions les unes après les autres. \newline \textbf{Sélection} : \texttt{Si ... Alors ... Sinon ... FinSi} (choix). \newline \textbf{Itération} : \texttt{Pour} (nb tours connu) / \texttt{TantQue} (condition). & Éviter les boucles infinies (condition de sortie modifiée dans le corps). \\
\hline
\textbf{Portée des variables} & \textbf{Locale} : Déclarée dans un bloc (fonction/procédure), invisible dehors. \newline \textbf{Globale} : Déclarée hors blocs, visible partout (à limiter). & Privilégier le passage par \emph{paramètres} (entrées/sorties). \\
\hline
\textbf{Test / Qualité} & \textbf{Jeu de tests} : Cas normaux, limites (bords), invalides/erroneux. \newline \textbf{Traces d'exécution} : Tableau (Variable | Valeur) pour suivre l'algo. \newline \textbf{Documentation} : Commentaires (//, /* */), guide utilisateur, doc technique. & Un logiciel non testé est un logiciel qui ne marche pas. \\
\hline
\end{tabularx}
\end{center}

\courssub{Méthodes express}

\begin{methode}[Écrire un algorithme structuré (Pseudo-code)]
\begin{enumerate}
    \item \textbf{Analyser} : Identifier les données d'entrée (Données), le résultat attendu (Résultats) et les traitements (Traitements).
    \item \textbf{Déclarer} les variables (Nom : Type) dans la section \texttt{Variables}.
    \item \textbf{Écrire} le corps (\texttt{Début} ... \texttt{Fin}) en utilisant \textbf{uniquement} les 3 structures de contrôle (Séquence, Si, Pour/TantQue).
    \item \textbf{Modulariser} : Si une tâche est répétée ou complexe $\to$ Créer une \texttt{Fonction} (retourne valeur) ou \texttt{Procédure} (action).
    \item \textbf{Tester} sur papier avec un \textbf{jeu de tests} (tableau de traces) avant de coder.
\end{enumerate}
\end{methode}

\begin{methode}[Déboguer efficacement (Premier niveau)]
\begin{enumerate}
    \item Lire le message d'erreur (ligne, type : syntaxe, exécution, logique).
    \item Utiliser des \texttt{Afficher} (prints) stratégiques pour suivre les valeurs des variables (trace pas à pas).
    \item Isoler le bloc fautif (commenter des parties, tester des fonctions seules).
    \item Corriger \textbf{une seule chose} à la fois, recompiler, retester.
\end{enumerate}
\end{methode}

\end{bilan}

\begin{aretenir}[Checklist avant le Bac]
\begin{itemize}
    \item $\square$ Je sais définir : Logiciel, Algorithme, Programme, Variable, Type, Fonction, Procédure, Bogue.
    \item $\square$ Je cite les 5 grandes phases du cycle de vie d'un logiciel (Analyse $\to$ Maintenance).
    \item $\square$ Je distingue Compilateur (ex: C) et Interpréteur (ex: Python, JS) et leurs impacts.
    \item $\square$ Je connais les 3 structures de contrôle fondamentales : Séquence, Sélection (\texttt{Si}), Itération (\texttt{Pour}, \texttt{TantQue}).
    \item $\square$ Je sais déclarer une variable avec un type adapté (\texttt{Entier}, \texttt{Réel}, \texttt{Chaîne}, \texttt{Booléen}).
    \item $\square$ Je sais écrire un algorithme complet en pseudo-code (Variables, Début, Fin, indentation correcte).
    \item $\square$ Je sais construire un \textbf{jeu de tests} pertinent (cas normaux, limites, erreurs).
    \item $\square$ Je sais réaliser une \textbf{table de traces} (exécution pas à pas) pour valider un algorithme.
    \item $\square$ Je comprends l'intérêt de la modularité (Fonctions/Procédures) : lisibilité, réutilisation, maintenance.
    \item $\square$ Je sais expliquer la différence entre paramètre \emph{entrée} (valeur) et \emph{sortie} (référence/adresse) -- notion de passage par valeur/adresse.
    \item $\square$ Je connais les bases de la documentation : commentaires code, nommage clair (camelCase/snake\_case), guide utilisateur.
    \item $\square$ Je suis capable de proposer une correction simple face à un bogue classique (boucle infinie, décalage d'indice, condition inversée).
\end{itemize}
\end{aretenir}

\findecours

\end{document}
